% 高一95_金汇中学_10月月考.tex
%   —— 高一上数学“2026-2027 年金汇中学高一上 10 月月考”（讲评版）
% 试卷版：xelatex 高一95_金汇中学_10月月考.tex
% 答案版：xelatex "\def\isanswer{1}% 高一95_金汇中学_10月月考.tex
%   —— 高一上数学“2026-2027 年金汇中学高一上 10 月月考”（讲评版）
% 试卷版：xelatex 高一95_金汇中学_10月月考.tex
% 答案版：xelatex "\def\isanswer{1}% 高一95_金汇中学_10月月考.tex
%   —— 高一上数学“2026-2027 年金汇中学高一上 10 月月考”（讲评版）
% 试卷版：xelatex 高一95_金汇中学_10月月考.tex
% 答案版：xelatex "\def\isanswer{1}% 高一95_金汇中学_10月月考.tex
%   —— 高一上数学“2026-2027 年金汇中学高一上 10 月月考”（讲评版）
% 试卷版：xelatex 高一95_金汇中学_10月月考.tex
% 答案版：xelatex "\def\isanswer{1}\input{高一95_金汇中学_10月月考.tex}"
% 紧凑版：xelatex "\def\iscompact{1}\input{高一95_金汇中学_10月月考.tex}"
%
% 原始材料：微信公众号“魔都数学加油站”文章，连载“27学年高一95”；
% 原文标题“27学年高一95——2026-2027年上海市金汇中学高一上10月月考”；
% 原文链接：https://mp.weixin.qq.com/s/H3yH73FTin3V7rxWQ-9fjw。
% 正文为 8 张 PNG 页（`原图/高一95_金汇中学_10月月考/01..08.png`），均 4762×6735，按页序读取。
%
% 篇幅与题号：原卷自第 3 题起编号（第 1、2 题不在原文内）。填空题 3—12（10 题）、
% 选择题 13—16（4 题）、解答题 17—21（5 题），共 19 题。
%
% 讲评版：第 4、5 题只印【答案】（“$m>0$”“$x=0$ 或 $x=1$”），本稿只给 \ans{}；
% 第 13、14 题无【解析】、答案“D”“B”直接印在题干括号内，本稿试卷版按惯例留空、
% 答案版提 \ans{}；其余 15 题印【解析】，其中第 15、16 题的答案在解析末句“故选 $A$／$C$”。
% 第 19 题解析带原卷的分步分值（“（3 分）”“（6 分）”等），照录。全卷无插图。
%
% 副标题：原卷未印副标题。本稿按内容标作“集合与常用逻辑用语、不等式”。
%
% 原卷笔误/存疑（均照录未改，正文对应处另有 % 注）：
%   ① 第 17 题 (1) 解析首句作“**第二个**不等式可化为 $(2x+5)(x-3)\leqslant0$”，而
%      $-2x^2+x+15\geqslant0$ 才是**第一个**不等式；原卷误标，照录；
%   ② 第 19 题 (2) 题干问“$x\in B$ 是 $x\in A$ 的充分非必要条件”，解析却从“$A\cup B=A$，
%      得 $B\subseteq A$”起论（少了“真子集”一层）；原卷照录，结论 $a\leqslant-1$ 未改。
%
% 与原卷的排版差异：
%   ① 原卷每题下直接印【答案】/【解析】，本稿按全稿惯例将答案与解析置于答案版；
%   ② 原卷未印副标题，本稿补出；填空横线统一排 \blankline[4.4em]；
%   ③ 原卷数集记号排斜体 $R$、$N$、$Z$，本稿按全稿惯例改用 $\R$、$\N$、$\Z$；
%   ④ 第 17—21 题的子问用行内“（1）（2）（3）”段，解析中的子问标号用 \stepm；
%   ⑤ 解答题按惯例加 \ansspace；选择题括号排 \mbox{（\quad）}。
%
% 行数核对：逐题数 \step/\stepm 条数，合计 87 条（Q3：2；Q4：0；Q5：0；Q6：1；Q7：2；
% Q8：4；Q9：3；Q10：3；Q11：7；Q12：3；Q13：0；Q14：0；Q15：2；Q16：14；Q17：7；
% Q18：5；Q19：12；Q20：10；Q21：12）。第 4、5 题只印答案，第 13、14 题无解析，
% 均不计入 \step 账。原卷解析的换行按语意合并，未逐行强拆。
%
\documentclass[a4paper,11pt]{article}
\usepackage{common}

\begin{document}

\examtitlest[7]{2026--2027 年金汇中学高一上 10 月月考}{集合与常用逻辑用语、不等式}

\bigsec{一、填空题}

\begin{enumerate}
\setcounter{enumi}{2}

\item 命题：如果 $ab>c$，那么 $a>\dfrac cb$，是\blankline[4.4em]（选填“真”、“假”）命题.

\ifanswer
\solbegin
\step{取 $a=1,b=-2,c=-3$，则 $ab=-2>-3=c$，但 $a=1<\dfrac cb=\dfrac32$，}
\step{所以原命题是假命题.}
\solend
\fi

\item 若“$x=0$”是“$x<m$”的充分条件，则实数 $m$ 的取值范围是\blankline[4.4em].

\ifanswer
\ans{$m>0$}
\fi

\item “$x\neq0$ 且 $x\neq1$”的否定形式为\blankline[4.4em].

\ifanswer
\ans{$x=0$ 或 $x=1$}
\fi

\item 若集合 $M=\{(x,y)\mid x+y=3\}$，$N=\left\{(x,y)\left|y=\dfrac12x\right.\right\}$，则 $M\cap N=\blankline[4.4em]$.

\ifanswer
\solbegin
\step{$x+\dfrac12x=3\Rightarrow x=2,y=1$，所以 $M\cap N=\{(2,1)\}$.}
\solend
\fi

\item 集合 $\{x\in\N\mid0<\lvert x-1\rvert<3\}$ 的真子集个数为\blankline[4.4em].

\ifanswer
\solbegin
\step{因为 $\{x\in\N\mid0<\lvert x-1\rvert<3\}=\{x\in\N\mid-3<x-1<3,x\neq1\}$}
\step{$=\{x\in\N\mid-2<x<4,x\neq1\}=\{0,2,3\}$，所以真子集的个数为 $2^3-1=7$.}
\solend
\fi

\item 设 $A=\{x\mid x^2-3x-4=0\}$，$B=\{x\mid ax-1=0\}$，若 $A\cap B=B$，则实数 $a$ 的值可以为\blankline[4.4em].

\ifanswer
\solbegin
\step{$A=\{-1,4\}$，因为 $A\cap B=B$，所以 $B\subseteq A$，}
\step{①$a=0$ 时，$B=\varnothing$，满足 $B\subseteq A$；}
\step{②$a\neq0$ 时，$B=\left\{\dfrac1a\right\}$，则 $\dfrac1a=-1$ 或 $\dfrac1a=4$，所以 $a=-1$ 或 $\dfrac14$；}
\step{所以 $a$ 的值可以为 $0$、$-1$、$\dfrac14$.}
\solend
\fi

\item 关于 $x$ 的不等式 $ax^2-(a+1)x+1\leqslant0(a<0)$ 的解集为\blankline[4.4em].

\ifanswer
\solbegin
\step{原不等式即 $(ax-1)(x-1)\leqslant0$，因为 $a<0$，所以 $\left(x-\dfrac1a\right)(x-1)\geqslant0$.}
\step{又 $\dfrac1a<0<1$，所以 $x\leqslant\dfrac1a$ 或 $x\geqslant1$，}
\step{故解集为 $\left(-\infty,\dfrac1a\right]\cup[1,+\infty)$.}
\solend
\fi

\item 若不等式 $(a-1)x^2+(2a-2)x-3<0$ 的解集是 $\R$，则 $a$ 的取值范围是\blankline[4.4em].

\ifanswer
\solbegin
\step{当 $a=1$ 时，不等式为 $-3<0$，对任意实数 $x$ 均成立.}
\step{当 $a\neq1$ 时，$\begin{cases}a-1<0\\ \Delta=(2a-2)^2+12(a-1)=4(a-1)(a+2)<0\end{cases}$，解得 $-2<a<1$.}
\step{合并得 $a\in(-2,1]$.}
\solend
\fi

\item 若关于 $x$ 的不等式 $x^2-(m+2)x+2m<0$ 的解集中恰有 $3$ 个整数，则实数 $m$ 的取值范围为\blankline[4.4em].

\ifanswer
\solbegin
\step{将不等式 $x^2-(m+2)x+2m<0$ 因式分解为 $(x-2)(x-m)<0$，}
\step{当 $m<2$ 时，解集为 $(m,2)$.}
\step{要使该区间恰有 $3$ 个整数，这 $3$ 个整数为 $-1,0,1$，故 $m$ 需满足 $-2\leqslant m<-1$.}
\step{当 $m>2$ 时，解集为 $(2,m)$. 要使该区间恰有 $3$ 个整数，}
\step{这 $3$ 个整数为 $3$，$4$，$5$，故 $m$ 需满足 $5<m\leqslant6$.}
\step{当 $m=2$ 时，不等式无解，不满足条件.}
\step{综上，实数 $m$ 的取值范围是 $[-2,-1)\cup(5,6]$.}
\solend
\fi

\item 已知集合 $A=\{x\mid x^2-3x+2\leqslant0\},B=\{y\mid y=x^2-2x+a,x\in\R\}$，若对任何一个 $x\in A$，都有 $x\in B$，则实数 $a$ 的取值范围是\blankline[4.4em].

\ifanswer
\solbegin
\step{$(x-1)(x-2)\leqslant0$，所以 $A=[1,2]$.}
\step{由 $y=(x-1)^2+a-1,x\in\R$，得 $B=[a-1,+\infty)$.}
\step{因此 $A\subseteq B\Leftrightarrow a-1\leqslant1\Leftrightarrow a\leqslant2$，所求范围为 $(-\infty,2]$.}
\solend
\fi

\end{enumerate}

\bigsec{二、选择题}

\begin{enumerate}
\setcounter{enumi}{12}

\item 已知 $a=\sqrt2$，集合 $A=\{x\mid x\leqslant2\}$，则下列表示正确的是\mbox{（\quad）}% 原卷括号内已印“D”，试卷版按惯例留空

\keepwithprev
A. $\{a\}\in A$\qquad B. $a\notin A$\qquad C. $a\subset A$\qquad D. $a\in A$

\ifanswer
\ans{$D$}
\fi

\item 下列四个命题中，真命题是\mbox{（\quad）}% 原卷括号内已印“B”，试卷版按惯例留空

\keepwithprev
A. 若 $a>b$，则 $ac^2>bc^2$\qquad B. 若 $a>b,c>d$，则 $a-d>b-c$

C. 若 $\lvert a\rvert>b$，则 $a^2>b^2$\qquad D. 若 $a>b$，则 $\dfrac1a<\dfrac1b$

\ifanswer
\ans{$B$}
\fi

\item “$x>1$”是“$\dfrac1x<1$”的\mbox{（\quad）}

\keepwithprev
A. 充分不必要条件\qquad B. 必要不充分条件

C. 充要条件\qquad D. 既不充分也不必要条件

\ifanswer
\solbegin
\step{$x>1\Rightarrow0<\dfrac1x<1$，充分性成立.}
\step{当 $x=-1$ 时，$\dfrac1x=-1<1$，但 $x>1$ 不成立，必要性不成立. 故选 $A$.}
\solend
\fi

\item 对于集合 $M=\{a\mid a=x^2-y^2,x\in\Z,y\in\Z\}$，给出如下三个结论：①如果 $P=\{b\mid b=2n+1,n\in\Z\}$，那么 $P\subseteq M$；②如果 $c=4n+2,n\in\Z$，那么 $c\notin M$；③如果 $a_1\in M,a_2\in M$，那么 $a_1a_2\in M$. 其中正确结论的个数是\mbox{（\quad）}

\keepwithprev
A. $1$\qquad B. $2$\qquad C. $3$\qquad D. $0$

\ifanswer
\solbegin
\step{对于①：因为 $2n+1=(n+1)^2-n^2,n\in\Z$，所以 $b=2n+1\in M$，故 $P\subseteq M$，}
\step{故①正确；}
\step{对于②：因为 $c=4n+2=2(2n+1),n\in\Z$，所以 $c$ 为偶数，且不能被 $4$ 整除，}
\step{若 $c\in M$，则存在 $x\in\Z,y\in\Z$，得 $c=x^2-y^2=(x+y)(x-y),x\in\Z,y\in\Z$，}
\step{因为 $x+y$ 和 $x-y$ 同奇或同偶，}
\step{若 $x+y$ 和 $x-y$ 同奇，则 $c=(x+y)(x-y)$ 为奇数，矛盾，不符合，}
\step{若 $x+y$ 和 $x-y$ 同偶，则 $c=(x+y)(x-y)$ 能被 $4$ 整除，矛盾，不符合，}
\step{所以 $c\notin M$，故②正确；}
\step{对于③：因为 $a_1\in M,a_2\in M$，}
\step{所以存在 $x_1,x_2\in\Z,y_1,y_2\in\Z$ 使得 $a_1=x_1^2-y_1^2,a_2=x_2^2-y_2^2$，}
\step{所以 $a_1a_2=(x_1^2-y_1^2)(x_2^2-y_2^2)=(x_1x_2)^2+(y_1y_2)^2-(x_1y_2)^2-(x_2y_1)^2$}
\step{$=(x_1x_2+y_1y_2)^2-(x_1y_2+x_2y_1)^2$，}
\step{因为 $x_1x_2+y_1y_2,x_1y_2+x_2y_1\in\Z$，所以 $a_1a_2\in M$，故③正确.}
\step{故选 $C$.}
\solend
\fi

\end{enumerate}

\bigsec{三、解答题}

\begin{enumerate}
\setcounter{enumi}{16}

\item （1）求不等式组 $\begin{cases}-2x^2+x+15\geqslant0\\ (x+1)(3x-1)>0\end{cases}$ 的解集；

（2）讨论解关于 $x$ 的不等式 $mx+4<m^2+2x$.

\ansspace*[12]

\ifanswer
\solbegin
\stepm{（1）}{第二个不等式可化为 $(2x+5)(x-3)\leqslant0$，解得 $-\dfrac52\leqslant x\leqslant3$.}% 原卷把“第一个”误作“第二个”，照录
\step{第二个不等式 $(x+1)(3x-1)>0$ 的解集为 $(-\infty,-1)\cup\left(\dfrac13,+\infty\right)$.}
\step{原不等式组的解集为 $\left[-\dfrac52,-1\right)\cup\left(\dfrac13,3\right]$.}
\stepm{（2）}{移项并因式分解，得 $(m-2)x<m^2-4\Leftrightarrow(m-2)x<(m-2)(m+2)$.}
\step{当 $m>2$ 时，$m-2>0$，解集为 $(-\infty,m+2)$.}
\step{当 $m<2$ 时，$m-2<0$，解集为 $(m+2,+\infty)$.}
\step{当 $m=2$ 时，不等式化为 $0<0$，不成立，解集为 $\varnothing$.}
\solend
\fi

\item （1）设 $x$ 是实数，比较 $x^2+9$ 与 $6x$ 的值的大小；

（2）已知 $a>0$ 且 $a\neq1$，比较 $\dfrac{a^2+1}{a^2-1}$ 与 $\dfrac{a+1}{a-1}$ 的大小.

\ansspace*[11]

\ifanswer
\solbegin
\stepm{（1）}{作差得 $(x^2+9)-6x=(x-3)^2\geqslant0$，因此 $x^2+9\geqslant6x$；}
\stepm{（2）}{$\dfrac{a^2+1}{a^2-1}-\dfrac{a+1}{a-1}=\dfrac{-2a}{a^2-1}$，}
\step{当 $a>1$ 时，$-2a<0$，$a^2-1>0$，则 $\dfrac{-2a}{a^2-1}<0$，即 $\dfrac{a^2+1}{a^2-1}<\dfrac{a+1}{a-1}$；}
\step{当 $0<a<1$ 时，$-2a<0$，$a^2-1<0$，则 $\dfrac{-2a}{a^2-1}>0$，即 $\dfrac{a^2+1}{a^2-1}>\dfrac{a+1}{a-1}$；}
\step{综上，$a>1$ 时，$\dfrac{a^2+1}{a^2-1}<\dfrac{a+1}{a-1}$；$0<a<1$ 时，$\dfrac{a^2+1}{a^2-1}>\dfrac{a+1}{a-1}$.}
\solend
\fi

\item 已知 $a\in\R$，集合 $A=\left\{x\left|\dfrac{2-x}{x+1}\leqslant0\right.\right\},B=\{x\mid x^2+(3-a)x-3a<0\}$.

（1）当 $a=\dfrac13$ 时，求 $A$ 和 $B$；

（2）“$x\in B$”是“$x\in A$”的充分非必要条件，求实数 $a$ 的取值范围.

\ansspace*[13]

\ifanswer
\solbegin
\stepm{（1）}{由 $\dfrac{2-x}{x+1}\leqslant0$，得 $\begin{cases}(2-x)(x+1)\leqslant0\\ x\neq-1\end{cases}$，解得 $x<-1$ 或 $x\geqslant2$，}
\step{则 $A=\{x\mid x<-1$ 或 $x\geqslant2\}$；$\cdots\cdots\cdots\cdots\cdots\cdots\cdots\cdots\cdots\cdots\cdots\cdots\cdots\cdots\cdots\cdots$（3 分）}
\step{$B=\{x\mid x^2+(3-a)x-3a<0\}=\{x\mid(x+3)(x-a)<0\}$，}
\step{当 $a=\dfrac13$ 时，$B=\left\{x\left|(x+2)\left(x-\dfrac13\right)<0\right.\right\}$，}
\step{解得 $B=\left\{x\left|-3<x<\dfrac13\right.\right\}\cdots\cdots\cdots\cdots\cdots\cdots$（6 分）}
\stepm{（2）}{由 $A\cup B=A$，得 $B\subseteq A$，}% 原卷从 A∪B=A 起论，与题干“充分非必要”差一层，照录
\step{①当 $B=\varnothing$ 时，得 $a=-3$，符合题意；$\cdots\cdots\cdots\cdots\cdots$（2 分）}
\step{②当 $B\neq\varnothing$ 时，若 $a>-3$，则 $B=\{x\mid-3<x<a\}$，}
\step{由 $B\subseteq A$，得 $a\leqslant-1$，此时 $-3<a\leqslant-1$；$\cdots\cdots\cdots\cdots\cdots$（4 分）}
\step{若 $a<-3$，则 $B=\{x\mid a<x<-3\}$，此时 $B\subseteq A$ 恒成立，}
\step{故 $a<-3$ 符合题意；$\cdots\cdots\cdots$（6 分）}
\step{综上所述，实数 $a$ 的取值范围为 $a\leqslant-1$. $\cdots\cdots\cdots\cdots\cdots$（8 分）}
\solend
\fi

\item 已知集合 $A=\{x\mid x<-2$ 或 $x\geqslant3\}$，$U=R$.

（1）求 $\overline A$；

（2）若 $B=\{x\mid mx-1>0\}$，当 $B\subseteq A$ 时，求实数 $m$ 的取值范围；

（3）若 $B=\{x\mid2m-6\leqslant x\leqslant m+4\}$，当 $B\cap\overline A=\varnothing$ 时，求实数 $m$ 的取值范围.

\ansspace*[14]

\ifanswer
\solbegin
\stepm{（1）}{因为集合 $A=\{x\mid x<-2$ 或 $x\geqslant3\}$，$U=R$，所以 $\overline A=\{x\mid-2\leqslant x<3\}$；}
\stepm{（2）}{当 $m=0$ 时，$B=\varnothing$，符合题意，}
\step{当 $m>0$ 时，$B=\left\{x\left|x>\dfrac1m\right.\right\}$，则 $\dfrac1m\geqslant3$，解得 $0<m\leqslant\dfrac13$，}
\step{当 $m<0$ 时，$B=\left\{x\left|x<\dfrac1m\right.\right\}$，则 $\dfrac1m\leqslant-2$，解得 $-\dfrac12\leqslant m<0$，}
\step{综上所述，实数 $m$ 的取值范围 $\left[-\dfrac12,\dfrac13\right]$；}
\stepm{（3）}{由题意得 $B=\{x\mid2m-6\leqslant x\leqslant m+4\}$，$\overline A=\{x\mid-2\leqslant x<3\}$，且 $B\cap\overline A=\varnothing$，}
\step{当 $B=\varnothing$ 时，$2m-6>m+4$，解得 $m>10$，}
\step{当 $B\neq\varnothing$ 时，则 $\begin{cases}2m-6\leqslant m+4\\ m+4<-2\end{cases}$ 或 $\begin{cases}2m-6\leqslant m+4\\ 2m-6\geqslant3\end{cases}$，}
\step{解得 $m<-6$ 或 $\dfrac92\leqslant m\leqslant10$，}
\step{综上所述，实数 $m$ 的取值范围为 $(-\infty,-6)\cup\left[\dfrac92,+\infty\right)$.}
\solend
\fi

\item 已知命题 $p$：关于 $x$ 的方程 $x^2+mx+1=0$ 有两个不同的正根；命题 $q$：关于 $x$ 的方程 $4x^2+4(m-2)x=-1$ 有实根.

（1）若命题 $q$ 为假命题，求整数 $m$ 的解集；

（2）若命题 $p$ 为真命题，且关于 $x$ 的方程 $x^2+mx+1=0$ 的两个实数根 $x_1$、$x_2$ 满足 $3x_1+2x_2=5$，求实数 $m$ 的取值集合；

（3）若命题 $p$ 和命题 $q$ 有且仅有一个成立，求实数 $m$ 的取值范围.

\ansspace*[16]

\ifanswer
\solbegin
\stepm{（1）}{命题 $q$：关于 $x$ 的方程 $4x^2+4(m-2)x=-1$ 有实根为假命题，}
\step{即 $4x^2+4(m-2)x+1=0$ 无实根，}
\step{则 $\Delta=16(m-2)^2-16<0$，则 $1<m<3$，则整数 $m$ 的解集为 $\{2\}$；}
\stepm{（2）}{若命题 $p$ 为真命题，且关于 $x$ 的方程 $x^2+mx+1=0$ 的两个实数根 $x_1$、$x_2$}
\step{满足 $3x_1+2x_2=5$，则方程 $x^2+mx+1=0$ 有两个不同的正根，}
\step{则 $\begin{cases}\Delta=m^2-4>0\\ x_1+x_2=-m>0\end{cases}$，则 $m<-2$，又 $\begin{cases}x_1+x_2=-m\\ x_1x_2=1\\ 3x_1+2x_2=5\end{cases}$，则 $\begin{cases}x_1=5+2m\\ x_2=-3m-5\end{cases}$，}
\step{则 $(5+2m)(-3m-5)=1$，得 $m=-2$ 或 $m=-\dfrac{13}6$，}
\step{又因为 $m<-2$，所以 $m=-\dfrac{13}6$，所以实数 $m$ 的取值集合为 $\left\{-\dfrac{13}6\right\}$；}
\stepm{（3）}{由（2）得，当 $p$ 为真时，则有 $m<-2$，}
\step{当 $q$ 为真时，$\Delta=16(m-2)^2-16\geqslant0$，解得 $m\leqslant1$ 或 $m\geqslant3$，}
\step{又因为命题 $p$ 和命题 $q$ 只有一个为真，所以 $\begin{cases}m<-2\\ 1<m<3\end{cases}$ 或 $\begin{cases}m\geqslant-2\\ m\leqslant1\ \text{或}\ m\geqslant3\end{cases}$，}
\step{解得 $-2\leqslant m\leqslant1$ 或 $m\geqslant3$，所以实数 $m$ 的取值范围为 $[-2,1]\cup[3,+\infty)$.}
\solend
\fi

\end{enumerate}

\end{document}
"
% 紧凑版：xelatex "\def\iscompact{1}% 高一95_金汇中学_10月月考.tex
%   —— 高一上数学“2026-2027 年金汇中学高一上 10 月月考”（讲评版）
% 试卷版：xelatex 高一95_金汇中学_10月月考.tex
% 答案版：xelatex "\def\isanswer{1}\input{高一95_金汇中学_10月月考.tex}"
% 紧凑版：xelatex "\def\iscompact{1}\input{高一95_金汇中学_10月月考.tex}"
%
% 原始材料：微信公众号“魔都数学加油站”文章，连载“27学年高一95”；
% 原文标题“27学年高一95——2026-2027年上海市金汇中学高一上10月月考”；
% 原文链接：https://mp.weixin.qq.com/s/H3yH73FTin3V7rxWQ-9fjw。
% 正文为 8 张 PNG 页（`原图/高一95_金汇中学_10月月考/01..08.png`），均 4762×6735，按页序读取。
%
% 篇幅与题号：原卷自第 3 题起编号（第 1、2 题不在原文内）。填空题 3—12（10 题）、
% 选择题 13—16（4 题）、解答题 17—21（5 题），共 19 题。
%
% 讲评版：第 4、5 题只印【答案】（“$m>0$”“$x=0$ 或 $x=1$”），本稿只给 \ans{}；
% 第 13、14 题无【解析】、答案“D”“B”直接印在题干括号内，本稿试卷版按惯例留空、
% 答案版提 \ans{}；其余 15 题印【解析】，其中第 15、16 题的答案在解析末句“故选 $A$／$C$”。
% 第 19 题解析带原卷的分步分值（“（3 分）”“（6 分）”等），照录。全卷无插图。
%
% 副标题：原卷未印副标题。本稿按内容标作“集合与常用逻辑用语、不等式”。
%
% 原卷笔误/存疑（均照录未改，正文对应处另有 % 注）：
%   ① 第 17 题 (1) 解析首句作“**第二个**不等式可化为 $(2x+5)(x-3)\leqslant0$”，而
%      $-2x^2+x+15\geqslant0$ 才是**第一个**不等式；原卷误标，照录；
%   ② 第 19 题 (2) 题干问“$x\in B$ 是 $x\in A$ 的充分非必要条件”，解析却从“$A\cup B=A$，
%      得 $B\subseteq A$”起论（少了“真子集”一层）；原卷照录，结论 $a\leqslant-1$ 未改。
%
% 与原卷的排版差异：
%   ① 原卷每题下直接印【答案】/【解析】，本稿按全稿惯例将答案与解析置于答案版；
%   ② 原卷未印副标题，本稿补出；填空横线统一排 \blankline[4.4em]；
%   ③ 原卷数集记号排斜体 $R$、$N$、$Z$，本稿按全稿惯例改用 $\R$、$\N$、$\Z$；
%   ④ 第 17—21 题的子问用行内“（1）（2）（3）”段，解析中的子问标号用 \stepm；
%   ⑤ 解答题按惯例加 \ansspace；选择题括号排 \mbox{（\quad）}。
%
% 行数核对：逐题数 \step/\stepm 条数，合计 87 条（Q3：2；Q4：0；Q5：0；Q6：1；Q7：2；
% Q8：4；Q9：3；Q10：3；Q11：7；Q12：3；Q13：0；Q14：0；Q15：2；Q16：14；Q17：7；
% Q18：5；Q19：12；Q20：10；Q21：12）。第 4、5 题只印答案，第 13、14 题无解析，
% 均不计入 \step 账。原卷解析的换行按语意合并，未逐行强拆。
%
\documentclass[a4paper,11pt]{article}
\usepackage{common}

\begin{document}

\examtitlest[7]{2026--2027 年金汇中学高一上 10 月月考}{集合与常用逻辑用语、不等式}

\bigsec{一、填空题}

\begin{enumerate}
\setcounter{enumi}{2}

\item 命题：如果 $ab>c$，那么 $a>\dfrac cb$，是\blankline[4.4em]（选填“真”、“假”）命题.

\ifanswer
\solbegin
\step{取 $a=1,b=-2,c=-3$，则 $ab=-2>-3=c$，但 $a=1<\dfrac cb=\dfrac32$，}
\step{所以原命题是假命题.}
\solend
\fi

\item 若“$x=0$”是“$x<m$”的充分条件，则实数 $m$ 的取值范围是\blankline[4.4em].

\ifanswer
\ans{$m>0$}
\fi

\item “$x\neq0$ 且 $x\neq1$”的否定形式为\blankline[4.4em].

\ifanswer
\ans{$x=0$ 或 $x=1$}
\fi

\item 若集合 $M=\{(x,y)\mid x+y=3\}$，$N=\left\{(x,y)\left|y=\dfrac12x\right.\right\}$，则 $M\cap N=\blankline[4.4em]$.

\ifanswer
\solbegin
\step{$x+\dfrac12x=3\Rightarrow x=2,y=1$，所以 $M\cap N=\{(2,1)\}$.}
\solend
\fi

\item 集合 $\{x\in\N\mid0<\lvert x-1\rvert<3\}$ 的真子集个数为\blankline[4.4em].

\ifanswer
\solbegin
\step{因为 $\{x\in\N\mid0<\lvert x-1\rvert<3\}=\{x\in\N\mid-3<x-1<3,x\neq1\}$}
\step{$=\{x\in\N\mid-2<x<4,x\neq1\}=\{0,2,3\}$，所以真子集的个数为 $2^3-1=7$.}
\solend
\fi

\item 设 $A=\{x\mid x^2-3x-4=0\}$，$B=\{x\mid ax-1=0\}$，若 $A\cap B=B$，则实数 $a$ 的值可以为\blankline[4.4em].

\ifanswer
\solbegin
\step{$A=\{-1,4\}$，因为 $A\cap B=B$，所以 $B\subseteq A$，}
\step{①$a=0$ 时，$B=\varnothing$，满足 $B\subseteq A$；}
\step{②$a\neq0$ 时，$B=\left\{\dfrac1a\right\}$，则 $\dfrac1a=-1$ 或 $\dfrac1a=4$，所以 $a=-1$ 或 $\dfrac14$；}
\step{所以 $a$ 的值可以为 $0$、$-1$、$\dfrac14$.}
\solend
\fi

\item 关于 $x$ 的不等式 $ax^2-(a+1)x+1\leqslant0(a<0)$ 的解集为\blankline[4.4em].

\ifanswer
\solbegin
\step{原不等式即 $(ax-1)(x-1)\leqslant0$，因为 $a<0$，所以 $\left(x-\dfrac1a\right)(x-1)\geqslant0$.}
\step{又 $\dfrac1a<0<1$，所以 $x\leqslant\dfrac1a$ 或 $x\geqslant1$，}
\step{故解集为 $\left(-\infty,\dfrac1a\right]\cup[1,+\infty)$.}
\solend
\fi

\item 若不等式 $(a-1)x^2+(2a-2)x-3<0$ 的解集是 $\R$，则 $a$ 的取值范围是\blankline[4.4em].

\ifanswer
\solbegin
\step{当 $a=1$ 时，不等式为 $-3<0$，对任意实数 $x$ 均成立.}
\step{当 $a\neq1$ 时，$\begin{cases}a-1<0\\ \Delta=(2a-2)^2+12(a-1)=4(a-1)(a+2)<0\end{cases}$，解得 $-2<a<1$.}
\step{合并得 $a\in(-2,1]$.}
\solend
\fi

\item 若关于 $x$ 的不等式 $x^2-(m+2)x+2m<0$ 的解集中恰有 $3$ 个整数，则实数 $m$ 的取值范围为\blankline[4.4em].

\ifanswer
\solbegin
\step{将不等式 $x^2-(m+2)x+2m<0$ 因式分解为 $(x-2)(x-m)<0$，}
\step{当 $m<2$ 时，解集为 $(m,2)$.}
\step{要使该区间恰有 $3$ 个整数，这 $3$ 个整数为 $-1,0,1$，故 $m$ 需满足 $-2\leqslant m<-1$.}
\step{当 $m>2$ 时，解集为 $(2,m)$. 要使该区间恰有 $3$ 个整数，}
\step{这 $3$ 个整数为 $3$，$4$，$5$，故 $m$ 需满足 $5<m\leqslant6$.}
\step{当 $m=2$ 时，不等式无解，不满足条件.}
\step{综上，实数 $m$ 的取值范围是 $[-2,-1)\cup(5,6]$.}
\solend
\fi

\item 已知集合 $A=\{x\mid x^2-3x+2\leqslant0\},B=\{y\mid y=x^2-2x+a,x\in\R\}$，若对任何一个 $x\in A$，都有 $x\in B$，则实数 $a$ 的取值范围是\blankline[4.4em].

\ifanswer
\solbegin
\step{$(x-1)(x-2)\leqslant0$，所以 $A=[1,2]$.}
\step{由 $y=(x-1)^2+a-1,x\in\R$，得 $B=[a-1,+\infty)$.}
\step{因此 $A\subseteq B\Leftrightarrow a-1\leqslant1\Leftrightarrow a\leqslant2$，所求范围为 $(-\infty,2]$.}
\solend
\fi

\end{enumerate}

\bigsec{二、选择题}

\begin{enumerate}
\setcounter{enumi}{12}

\item 已知 $a=\sqrt2$，集合 $A=\{x\mid x\leqslant2\}$，则下列表示正确的是\mbox{（\quad）}% 原卷括号内已印“D”，试卷版按惯例留空

\keepwithprev
A. $\{a\}\in A$\qquad B. $a\notin A$\qquad C. $a\subset A$\qquad D. $a\in A$

\ifanswer
\ans{$D$}
\fi

\item 下列四个命题中，真命题是\mbox{（\quad）}% 原卷括号内已印“B”，试卷版按惯例留空

\keepwithprev
A. 若 $a>b$，则 $ac^2>bc^2$\qquad B. 若 $a>b,c>d$，则 $a-d>b-c$

C. 若 $\lvert a\rvert>b$，则 $a^2>b^2$\qquad D. 若 $a>b$，则 $\dfrac1a<\dfrac1b$

\ifanswer
\ans{$B$}
\fi

\item “$x>1$”是“$\dfrac1x<1$”的\mbox{（\quad）}

\keepwithprev
A. 充分不必要条件\qquad B. 必要不充分条件

C. 充要条件\qquad D. 既不充分也不必要条件

\ifanswer
\solbegin
\step{$x>1\Rightarrow0<\dfrac1x<1$，充分性成立.}
\step{当 $x=-1$ 时，$\dfrac1x=-1<1$，但 $x>1$ 不成立，必要性不成立. 故选 $A$.}
\solend
\fi

\item 对于集合 $M=\{a\mid a=x^2-y^2,x\in\Z,y\in\Z\}$，给出如下三个结论：①如果 $P=\{b\mid b=2n+1,n\in\Z\}$，那么 $P\subseteq M$；②如果 $c=4n+2,n\in\Z$，那么 $c\notin M$；③如果 $a_1\in M,a_2\in M$，那么 $a_1a_2\in M$. 其中正确结论的个数是\mbox{（\quad）}

\keepwithprev
A. $1$\qquad B. $2$\qquad C. $3$\qquad D. $0$

\ifanswer
\solbegin
\step{对于①：因为 $2n+1=(n+1)^2-n^2,n\in\Z$，所以 $b=2n+1\in M$，故 $P\subseteq M$，}
\step{故①正确；}
\step{对于②：因为 $c=4n+2=2(2n+1),n\in\Z$，所以 $c$ 为偶数，且不能被 $4$ 整除，}
\step{若 $c\in M$，则存在 $x\in\Z,y\in\Z$，得 $c=x^2-y^2=(x+y)(x-y),x\in\Z,y\in\Z$，}
\step{因为 $x+y$ 和 $x-y$ 同奇或同偶，}
\step{若 $x+y$ 和 $x-y$ 同奇，则 $c=(x+y)(x-y)$ 为奇数，矛盾，不符合，}
\step{若 $x+y$ 和 $x-y$ 同偶，则 $c=(x+y)(x-y)$ 能被 $4$ 整除，矛盾，不符合，}
\step{所以 $c\notin M$，故②正确；}
\step{对于③：因为 $a_1\in M,a_2\in M$，}
\step{所以存在 $x_1,x_2\in\Z,y_1,y_2\in\Z$ 使得 $a_1=x_1^2-y_1^2,a_2=x_2^2-y_2^2$，}
\step{所以 $a_1a_2=(x_1^2-y_1^2)(x_2^2-y_2^2)=(x_1x_2)^2+(y_1y_2)^2-(x_1y_2)^2-(x_2y_1)^2$}
\step{$=(x_1x_2+y_1y_2)^2-(x_1y_2+x_2y_1)^2$，}
\step{因为 $x_1x_2+y_1y_2,x_1y_2+x_2y_1\in\Z$，所以 $a_1a_2\in M$，故③正确.}
\step{故选 $C$.}
\solend
\fi

\end{enumerate}

\bigsec{三、解答题}

\begin{enumerate}
\setcounter{enumi}{16}

\item （1）求不等式组 $\begin{cases}-2x^2+x+15\geqslant0\\ (x+1)(3x-1)>0\end{cases}$ 的解集；

（2）讨论解关于 $x$ 的不等式 $mx+4<m^2+2x$.

\ansspace*[12]

\ifanswer
\solbegin
\stepm{（1）}{第二个不等式可化为 $(2x+5)(x-3)\leqslant0$，解得 $-\dfrac52\leqslant x\leqslant3$.}% 原卷把“第一个”误作“第二个”，照录
\step{第二个不等式 $(x+1)(3x-1)>0$ 的解集为 $(-\infty,-1)\cup\left(\dfrac13,+\infty\right)$.}
\step{原不等式组的解集为 $\left[-\dfrac52,-1\right)\cup\left(\dfrac13,3\right]$.}
\stepm{（2）}{移项并因式分解，得 $(m-2)x<m^2-4\Leftrightarrow(m-2)x<(m-2)(m+2)$.}
\step{当 $m>2$ 时，$m-2>0$，解集为 $(-\infty,m+2)$.}
\step{当 $m<2$ 时，$m-2<0$，解集为 $(m+2,+\infty)$.}
\step{当 $m=2$ 时，不等式化为 $0<0$，不成立，解集为 $\varnothing$.}
\solend
\fi

\item （1）设 $x$ 是实数，比较 $x^2+9$ 与 $6x$ 的值的大小；

（2）已知 $a>0$ 且 $a\neq1$，比较 $\dfrac{a^2+1}{a^2-1}$ 与 $\dfrac{a+1}{a-1}$ 的大小.

\ansspace*[11]

\ifanswer
\solbegin
\stepm{（1）}{作差得 $(x^2+9)-6x=(x-3)^2\geqslant0$，因此 $x^2+9\geqslant6x$；}
\stepm{（2）}{$\dfrac{a^2+1}{a^2-1}-\dfrac{a+1}{a-1}=\dfrac{-2a}{a^2-1}$，}
\step{当 $a>1$ 时，$-2a<0$，$a^2-1>0$，则 $\dfrac{-2a}{a^2-1}<0$，即 $\dfrac{a^2+1}{a^2-1}<\dfrac{a+1}{a-1}$；}
\step{当 $0<a<1$ 时，$-2a<0$，$a^2-1<0$，则 $\dfrac{-2a}{a^2-1}>0$，即 $\dfrac{a^2+1}{a^2-1}>\dfrac{a+1}{a-1}$；}
\step{综上，$a>1$ 时，$\dfrac{a^2+1}{a^2-1}<\dfrac{a+1}{a-1}$；$0<a<1$ 时，$\dfrac{a^2+1}{a^2-1}>\dfrac{a+1}{a-1}$.}
\solend
\fi

\item 已知 $a\in\R$，集合 $A=\left\{x\left|\dfrac{2-x}{x+1}\leqslant0\right.\right\},B=\{x\mid x^2+(3-a)x-3a<0\}$.

（1）当 $a=\dfrac13$ 时，求 $A$ 和 $B$；

（2）“$x\in B$”是“$x\in A$”的充分非必要条件，求实数 $a$ 的取值范围.

\ansspace*[13]

\ifanswer
\solbegin
\stepm{（1）}{由 $\dfrac{2-x}{x+1}\leqslant0$，得 $\begin{cases}(2-x)(x+1)\leqslant0\\ x\neq-1\end{cases}$，解得 $x<-1$ 或 $x\geqslant2$，}
\step{则 $A=\{x\mid x<-1$ 或 $x\geqslant2\}$；$\cdots\cdots\cdots\cdots\cdots\cdots\cdots\cdots\cdots\cdots\cdots\cdots\cdots\cdots\cdots\cdots$（3 分）}
\step{$B=\{x\mid x^2+(3-a)x-3a<0\}=\{x\mid(x+3)(x-a)<0\}$，}
\step{当 $a=\dfrac13$ 时，$B=\left\{x\left|(x+2)\left(x-\dfrac13\right)<0\right.\right\}$，}
\step{解得 $B=\left\{x\left|-3<x<\dfrac13\right.\right\}\cdots\cdots\cdots\cdots\cdots\cdots$（6 分）}
\stepm{（2）}{由 $A\cup B=A$，得 $B\subseteq A$，}% 原卷从 A∪B=A 起论，与题干“充分非必要”差一层，照录
\step{①当 $B=\varnothing$ 时，得 $a=-3$，符合题意；$\cdots\cdots\cdots\cdots\cdots$（2 分）}
\step{②当 $B\neq\varnothing$ 时，若 $a>-3$，则 $B=\{x\mid-3<x<a\}$，}
\step{由 $B\subseteq A$，得 $a\leqslant-1$，此时 $-3<a\leqslant-1$；$\cdots\cdots\cdots\cdots\cdots$（4 分）}
\step{若 $a<-3$，则 $B=\{x\mid a<x<-3\}$，此时 $B\subseteq A$ 恒成立，}
\step{故 $a<-3$ 符合题意；$\cdots\cdots\cdots$（6 分）}
\step{综上所述，实数 $a$ 的取值范围为 $a\leqslant-1$. $\cdots\cdots\cdots\cdots\cdots$（8 分）}
\solend
\fi

\item 已知集合 $A=\{x\mid x<-2$ 或 $x\geqslant3\}$，$U=R$.

（1）求 $\overline A$；

（2）若 $B=\{x\mid mx-1>0\}$，当 $B\subseteq A$ 时，求实数 $m$ 的取值范围；

（3）若 $B=\{x\mid2m-6\leqslant x\leqslant m+4\}$，当 $B\cap\overline A=\varnothing$ 时，求实数 $m$ 的取值范围.

\ansspace*[14]

\ifanswer
\solbegin
\stepm{（1）}{因为集合 $A=\{x\mid x<-2$ 或 $x\geqslant3\}$，$U=R$，所以 $\overline A=\{x\mid-2\leqslant x<3\}$；}
\stepm{（2）}{当 $m=0$ 时，$B=\varnothing$，符合题意，}
\step{当 $m>0$ 时，$B=\left\{x\left|x>\dfrac1m\right.\right\}$，则 $\dfrac1m\geqslant3$，解得 $0<m\leqslant\dfrac13$，}
\step{当 $m<0$ 时，$B=\left\{x\left|x<\dfrac1m\right.\right\}$，则 $\dfrac1m\leqslant-2$，解得 $-\dfrac12\leqslant m<0$，}
\step{综上所述，实数 $m$ 的取值范围 $\left[-\dfrac12,\dfrac13\right]$；}
\stepm{（3）}{由题意得 $B=\{x\mid2m-6\leqslant x\leqslant m+4\}$，$\overline A=\{x\mid-2\leqslant x<3\}$，且 $B\cap\overline A=\varnothing$，}
\step{当 $B=\varnothing$ 时，$2m-6>m+4$，解得 $m>10$，}
\step{当 $B\neq\varnothing$ 时，则 $\begin{cases}2m-6\leqslant m+4\\ m+4<-2\end{cases}$ 或 $\begin{cases}2m-6\leqslant m+4\\ 2m-6\geqslant3\end{cases}$，}
\step{解得 $m<-6$ 或 $\dfrac92\leqslant m\leqslant10$，}
\step{综上所述，实数 $m$ 的取值范围为 $(-\infty,-6)\cup\left[\dfrac92,+\infty\right)$.}
\solend
\fi

\item 已知命题 $p$：关于 $x$ 的方程 $x^2+mx+1=0$ 有两个不同的正根；命题 $q$：关于 $x$ 的方程 $4x^2+4(m-2)x=-1$ 有实根.

（1）若命题 $q$ 为假命题，求整数 $m$ 的解集；

（2）若命题 $p$ 为真命题，且关于 $x$ 的方程 $x^2+mx+1=0$ 的两个实数根 $x_1$、$x_2$ 满足 $3x_1+2x_2=5$，求实数 $m$ 的取值集合；

（3）若命题 $p$ 和命题 $q$ 有且仅有一个成立，求实数 $m$ 的取值范围.

\ansspace*[16]

\ifanswer
\solbegin
\stepm{（1）}{命题 $q$：关于 $x$ 的方程 $4x^2+4(m-2)x=-1$ 有实根为假命题，}
\step{即 $4x^2+4(m-2)x+1=0$ 无实根，}
\step{则 $\Delta=16(m-2)^2-16<0$，则 $1<m<3$，则整数 $m$ 的解集为 $\{2\}$；}
\stepm{（2）}{若命题 $p$ 为真命题，且关于 $x$ 的方程 $x^2+mx+1=0$ 的两个实数根 $x_1$、$x_2$}
\step{满足 $3x_1+2x_2=5$，则方程 $x^2+mx+1=0$ 有两个不同的正根，}
\step{则 $\begin{cases}\Delta=m^2-4>0\\ x_1+x_2=-m>0\end{cases}$，则 $m<-2$，又 $\begin{cases}x_1+x_2=-m\\ x_1x_2=1\\ 3x_1+2x_2=5\end{cases}$，则 $\begin{cases}x_1=5+2m\\ x_2=-3m-5\end{cases}$，}
\step{则 $(5+2m)(-3m-5)=1$，得 $m=-2$ 或 $m=-\dfrac{13}6$，}
\step{又因为 $m<-2$，所以 $m=-\dfrac{13}6$，所以实数 $m$ 的取值集合为 $\left\{-\dfrac{13}6\right\}$；}
\stepm{（3）}{由（2）得，当 $p$ 为真时，则有 $m<-2$，}
\step{当 $q$ 为真时，$\Delta=16(m-2)^2-16\geqslant0$，解得 $m\leqslant1$ 或 $m\geqslant3$，}
\step{又因为命题 $p$ 和命题 $q$ 只有一个为真，所以 $\begin{cases}m<-2\\ 1<m<3\end{cases}$ 或 $\begin{cases}m\geqslant-2\\ m\leqslant1\ \text{或}\ m\geqslant3\end{cases}$，}
\step{解得 $-2\leqslant m\leqslant1$ 或 $m\geqslant3$，所以实数 $m$ 的取值范围为 $[-2,1]\cup[3,+\infty)$.}
\solend
\fi

\end{enumerate}

\end{document}
"
%
% 原始材料：微信公众号“魔都数学加油站”文章，连载“27学年高一95”；
% 原文标题“27学年高一95——2026-2027年上海市金汇中学高一上10月月考”；
% 原文链接：https://mp.weixin.qq.com/s/H3yH73FTin3V7rxWQ-9fjw。
% 正文为 8 张 PNG 页（`原图/高一95_金汇中学_10月月考/01..08.png`），均 4762×6735，按页序读取。
%
% 篇幅与题号：原卷自第 3 题起编号（第 1、2 题不在原文内）。填空题 3—12（10 题）、
% 选择题 13—16（4 题）、解答题 17—21（5 题），共 19 题。
%
% 讲评版：第 4、5 题只印【答案】（“$m>0$”“$x=0$ 或 $x=1$”），本稿只给 \ans{}；
% 第 13、14 题无【解析】、答案“D”“B”直接印在题干括号内，本稿试卷版按惯例留空、
% 答案版提 \ans{}；其余 15 题印【解析】，其中第 15、16 题的答案在解析末句“故选 $A$／$C$”。
% 第 19 题解析带原卷的分步分值（“（3 分）”“（6 分）”等），照录。全卷无插图。
%
% 副标题：原卷未印副标题。本稿按内容标作“集合与常用逻辑用语、不等式”。
%
% 原卷笔误/存疑（均照录未改，正文对应处另有 % 注）：
%   ① 第 17 题 (1) 解析首句作“**第二个**不等式可化为 $(2x+5)(x-3)\leqslant0$”，而
%      $-2x^2+x+15\geqslant0$ 才是**第一个**不等式；原卷误标，照录；
%   ② 第 19 题 (2) 题干问“$x\in B$ 是 $x\in A$ 的充分非必要条件”，解析却从“$A\cup B=A$，
%      得 $B\subseteq A$”起论（少了“真子集”一层）；原卷照录，结论 $a\leqslant-1$ 未改。
%
% 与原卷的排版差异：
%   ① 原卷每题下直接印【答案】/【解析】，本稿按全稿惯例将答案与解析置于答案版；
%   ② 原卷未印副标题，本稿补出；填空横线统一排 \blankline[4.4em]；
%   ③ 原卷数集记号排斜体 $R$、$N$、$Z$，本稿按全稿惯例改用 $\R$、$\N$、$\Z$；
%   ④ 第 17—21 题的子问用行内“（1）（2）（3）”段，解析中的子问标号用 \stepm；
%   ⑤ 解答题按惯例加 \ansspace；选择题括号排 \mbox{（\quad）}。
%
% 行数核对：逐题数 \step/\stepm 条数，合计 87 条（Q3：2；Q4：0；Q5：0；Q6：1；Q7：2；
% Q8：4；Q9：3；Q10：3；Q11：7；Q12：3；Q13：0；Q14：0；Q15：2；Q16：14；Q17：7；
% Q18：5；Q19：12；Q20：10；Q21：12）。第 4、5 题只印答案，第 13、14 题无解析，
% 均不计入 \step 账。原卷解析的换行按语意合并，未逐行强拆。
%
\documentclass[a4paper,11pt]{article}
\usepackage{common}

\begin{document}

\examtitlest[7]{2026--2027 年金汇中学高一上 10 月月考}{集合与常用逻辑用语、不等式}

\bigsec{一、填空题}

\begin{enumerate}
\setcounter{enumi}{2}

\item 命题：如果 $ab>c$，那么 $a>\dfrac cb$，是\blankline[4.4em]（选填“真”、“假”）命题.

\ifanswer
\solbegin
\step{取 $a=1,b=-2,c=-3$，则 $ab=-2>-3=c$，但 $a=1<\dfrac cb=\dfrac32$，}
\step{所以原命题是假命题.}
\solend
\fi

\item 若“$x=0$”是“$x<m$”的充分条件，则实数 $m$ 的取值范围是\blankline[4.4em].

\ifanswer
\ans{$m>0$}
\fi

\item “$x\neq0$ 且 $x\neq1$”的否定形式为\blankline[4.4em].

\ifanswer
\ans{$x=0$ 或 $x=1$}
\fi

\item 若集合 $M=\{(x,y)\mid x+y=3\}$，$N=\left\{(x,y)\left|y=\dfrac12x\right.\right\}$，则 $M\cap N=\blankline[4.4em]$.

\ifanswer
\solbegin
\step{$x+\dfrac12x=3\Rightarrow x=2,y=1$，所以 $M\cap N=\{(2,1)\}$.}
\solend
\fi

\item 集合 $\{x\in\N\mid0<\lvert x-1\rvert<3\}$ 的真子集个数为\blankline[4.4em].

\ifanswer
\solbegin
\step{因为 $\{x\in\N\mid0<\lvert x-1\rvert<3\}=\{x\in\N\mid-3<x-1<3,x\neq1\}$}
\step{$=\{x\in\N\mid-2<x<4,x\neq1\}=\{0,2,3\}$，所以真子集的个数为 $2^3-1=7$.}
\solend
\fi

\item 设 $A=\{x\mid x^2-3x-4=0\}$，$B=\{x\mid ax-1=0\}$，若 $A\cap B=B$，则实数 $a$ 的值可以为\blankline[4.4em].

\ifanswer
\solbegin
\step{$A=\{-1,4\}$，因为 $A\cap B=B$，所以 $B\subseteq A$，}
\step{①$a=0$ 时，$B=\varnothing$，满足 $B\subseteq A$；}
\step{②$a\neq0$ 时，$B=\left\{\dfrac1a\right\}$，则 $\dfrac1a=-1$ 或 $\dfrac1a=4$，所以 $a=-1$ 或 $\dfrac14$；}
\step{所以 $a$ 的值可以为 $0$、$-1$、$\dfrac14$.}
\solend
\fi

\item 关于 $x$ 的不等式 $ax^2-(a+1)x+1\leqslant0(a<0)$ 的解集为\blankline[4.4em].

\ifanswer
\solbegin
\step{原不等式即 $(ax-1)(x-1)\leqslant0$，因为 $a<0$，所以 $\left(x-\dfrac1a\right)(x-1)\geqslant0$.}
\step{又 $\dfrac1a<0<1$，所以 $x\leqslant\dfrac1a$ 或 $x\geqslant1$，}
\step{故解集为 $\left(-\infty,\dfrac1a\right]\cup[1,+\infty)$.}
\solend
\fi

\item 若不等式 $(a-1)x^2+(2a-2)x-3<0$ 的解集是 $\R$，则 $a$ 的取值范围是\blankline[4.4em].

\ifanswer
\solbegin
\step{当 $a=1$ 时，不等式为 $-3<0$，对任意实数 $x$ 均成立.}
\step{当 $a\neq1$ 时，$\begin{cases}a-1<0\\ \Delta=(2a-2)^2+12(a-1)=4(a-1)(a+2)<0\end{cases}$，解得 $-2<a<1$.}
\step{合并得 $a\in(-2,1]$.}
\solend
\fi

\item 若关于 $x$ 的不等式 $x^2-(m+2)x+2m<0$ 的解集中恰有 $3$ 个整数，则实数 $m$ 的取值范围为\blankline[4.4em].

\ifanswer
\solbegin
\step{将不等式 $x^2-(m+2)x+2m<0$ 因式分解为 $(x-2)(x-m)<0$，}
\step{当 $m<2$ 时，解集为 $(m,2)$.}
\step{要使该区间恰有 $3$ 个整数，这 $3$ 个整数为 $-1,0,1$，故 $m$ 需满足 $-2\leqslant m<-1$.}
\step{当 $m>2$ 时，解集为 $(2,m)$. 要使该区间恰有 $3$ 个整数，}
\step{这 $3$ 个整数为 $3$，$4$，$5$，故 $m$ 需满足 $5<m\leqslant6$.}
\step{当 $m=2$ 时，不等式无解，不满足条件.}
\step{综上，实数 $m$ 的取值范围是 $[-2,-1)\cup(5,6]$.}
\solend
\fi

\item 已知集合 $A=\{x\mid x^2-3x+2\leqslant0\},B=\{y\mid y=x^2-2x+a,x\in\R\}$，若对任何一个 $x\in A$，都有 $x\in B$，则实数 $a$ 的取值范围是\blankline[4.4em].

\ifanswer
\solbegin
\step{$(x-1)(x-2)\leqslant0$，所以 $A=[1,2]$.}
\step{由 $y=(x-1)^2+a-1,x\in\R$，得 $B=[a-1,+\infty)$.}
\step{因此 $A\subseteq B\Leftrightarrow a-1\leqslant1\Leftrightarrow a\leqslant2$，所求范围为 $(-\infty,2]$.}
\solend
\fi

\end{enumerate}

\bigsec{二、选择题}

\begin{enumerate}
\setcounter{enumi}{12}

\item 已知 $a=\sqrt2$，集合 $A=\{x\mid x\leqslant2\}$，则下列表示正确的是\mbox{（\quad）}% 原卷括号内已印“D”，试卷版按惯例留空

\keepwithprev
A. $\{a\}\in A$\qquad B. $a\notin A$\qquad C. $a\subset A$\qquad D. $a\in A$

\ifanswer
\ans{$D$}
\fi

\item 下列四个命题中，真命题是\mbox{（\quad）}% 原卷括号内已印“B”，试卷版按惯例留空

\keepwithprev
A. 若 $a>b$，则 $ac^2>bc^2$\qquad B. 若 $a>b,c>d$，则 $a-d>b-c$

C. 若 $\lvert a\rvert>b$，则 $a^2>b^2$\qquad D. 若 $a>b$，则 $\dfrac1a<\dfrac1b$

\ifanswer
\ans{$B$}
\fi

\item “$x>1$”是“$\dfrac1x<1$”的\mbox{（\quad）}

\keepwithprev
A. 充分不必要条件\qquad B. 必要不充分条件

C. 充要条件\qquad D. 既不充分也不必要条件

\ifanswer
\solbegin
\step{$x>1\Rightarrow0<\dfrac1x<1$，充分性成立.}
\step{当 $x=-1$ 时，$\dfrac1x=-1<1$，但 $x>1$ 不成立，必要性不成立. 故选 $A$.}
\solend
\fi

\item 对于集合 $M=\{a\mid a=x^2-y^2,x\in\Z,y\in\Z\}$，给出如下三个结论：①如果 $P=\{b\mid b=2n+1,n\in\Z\}$，那么 $P\subseteq M$；②如果 $c=4n+2,n\in\Z$，那么 $c\notin M$；③如果 $a_1\in M,a_2\in M$，那么 $a_1a_2\in M$. 其中正确结论的个数是\mbox{（\quad）}

\keepwithprev
A. $1$\qquad B. $2$\qquad C. $3$\qquad D. $0$

\ifanswer
\solbegin
\step{对于①：因为 $2n+1=(n+1)^2-n^2,n\in\Z$，所以 $b=2n+1\in M$，故 $P\subseteq M$，}
\step{故①正确；}
\step{对于②：因为 $c=4n+2=2(2n+1),n\in\Z$，所以 $c$ 为偶数，且不能被 $4$ 整除，}
\step{若 $c\in M$，则存在 $x\in\Z,y\in\Z$，得 $c=x^2-y^2=(x+y)(x-y),x\in\Z,y\in\Z$，}
\step{因为 $x+y$ 和 $x-y$ 同奇或同偶，}
\step{若 $x+y$ 和 $x-y$ 同奇，则 $c=(x+y)(x-y)$ 为奇数，矛盾，不符合，}
\step{若 $x+y$ 和 $x-y$ 同偶，则 $c=(x+y)(x-y)$ 能被 $4$ 整除，矛盾，不符合，}
\step{所以 $c\notin M$，故②正确；}
\step{对于③：因为 $a_1\in M,a_2\in M$，}
\step{所以存在 $x_1,x_2\in\Z,y_1,y_2\in\Z$ 使得 $a_1=x_1^2-y_1^2,a_2=x_2^2-y_2^2$，}
\step{所以 $a_1a_2=(x_1^2-y_1^2)(x_2^2-y_2^2)=(x_1x_2)^2+(y_1y_2)^2-(x_1y_2)^2-(x_2y_1)^2$}
\step{$=(x_1x_2+y_1y_2)^2-(x_1y_2+x_2y_1)^2$，}
\step{因为 $x_1x_2+y_1y_2,x_1y_2+x_2y_1\in\Z$，所以 $a_1a_2\in M$，故③正确.}
\step{故选 $C$.}
\solend
\fi

\end{enumerate}

\bigsec{三、解答题}

\begin{enumerate}
\setcounter{enumi}{16}

\item （1）求不等式组 $\begin{cases}-2x^2+x+15\geqslant0\\ (x+1)(3x-1)>0\end{cases}$ 的解集；

（2）讨论解关于 $x$ 的不等式 $mx+4<m^2+2x$.

\ansspace*[12]

\ifanswer
\solbegin
\stepm{（1）}{第二个不等式可化为 $(2x+5)(x-3)\leqslant0$，解得 $-\dfrac52\leqslant x\leqslant3$.}% 原卷把“第一个”误作“第二个”，照录
\step{第二个不等式 $(x+1)(3x-1)>0$ 的解集为 $(-\infty,-1)\cup\left(\dfrac13,+\infty\right)$.}
\step{原不等式组的解集为 $\left[-\dfrac52,-1\right)\cup\left(\dfrac13,3\right]$.}
\stepm{（2）}{移项并因式分解，得 $(m-2)x<m^2-4\Leftrightarrow(m-2)x<(m-2)(m+2)$.}
\step{当 $m>2$ 时，$m-2>0$，解集为 $(-\infty,m+2)$.}
\step{当 $m<2$ 时，$m-2<0$，解集为 $(m+2,+\infty)$.}
\step{当 $m=2$ 时，不等式化为 $0<0$，不成立，解集为 $\varnothing$.}
\solend
\fi

\item （1）设 $x$ 是实数，比较 $x^2+9$ 与 $6x$ 的值的大小；

（2）已知 $a>0$ 且 $a\neq1$，比较 $\dfrac{a^2+1}{a^2-1}$ 与 $\dfrac{a+1}{a-1}$ 的大小.

\ansspace*[11]

\ifanswer
\solbegin
\stepm{（1）}{作差得 $(x^2+9)-6x=(x-3)^2\geqslant0$，因此 $x^2+9\geqslant6x$；}
\stepm{（2）}{$\dfrac{a^2+1}{a^2-1}-\dfrac{a+1}{a-1}=\dfrac{-2a}{a^2-1}$，}
\step{当 $a>1$ 时，$-2a<0$，$a^2-1>0$，则 $\dfrac{-2a}{a^2-1}<0$，即 $\dfrac{a^2+1}{a^2-1}<\dfrac{a+1}{a-1}$；}
\step{当 $0<a<1$ 时，$-2a<0$，$a^2-1<0$，则 $\dfrac{-2a}{a^2-1}>0$，即 $\dfrac{a^2+1}{a^2-1}>\dfrac{a+1}{a-1}$；}
\step{综上，$a>1$ 时，$\dfrac{a^2+1}{a^2-1}<\dfrac{a+1}{a-1}$；$0<a<1$ 时，$\dfrac{a^2+1}{a^2-1}>\dfrac{a+1}{a-1}$.}
\solend
\fi

\item 已知 $a\in\R$，集合 $A=\left\{x\left|\dfrac{2-x}{x+1}\leqslant0\right.\right\},B=\{x\mid x^2+(3-a)x-3a<0\}$.

（1）当 $a=\dfrac13$ 时，求 $A$ 和 $B$；

（2）“$x\in B$”是“$x\in A$”的充分非必要条件，求实数 $a$ 的取值范围.

\ansspace*[13]

\ifanswer
\solbegin
\stepm{（1）}{由 $\dfrac{2-x}{x+1}\leqslant0$，得 $\begin{cases}(2-x)(x+1)\leqslant0\\ x\neq-1\end{cases}$，解得 $x<-1$ 或 $x\geqslant2$，}
\step{则 $A=\{x\mid x<-1$ 或 $x\geqslant2\}$；$\cdots\cdots\cdots\cdots\cdots\cdots\cdots\cdots\cdots\cdots\cdots\cdots\cdots\cdots\cdots\cdots$（3 分）}
\step{$B=\{x\mid x^2+(3-a)x-3a<0\}=\{x\mid(x+3)(x-a)<0\}$，}
\step{当 $a=\dfrac13$ 时，$B=\left\{x\left|(x+2)\left(x-\dfrac13\right)<0\right.\right\}$，}
\step{解得 $B=\left\{x\left|-3<x<\dfrac13\right.\right\}\cdots\cdots\cdots\cdots\cdots\cdots$（6 分）}
\stepm{（2）}{由 $A\cup B=A$，得 $B\subseteq A$，}% 原卷从 A∪B=A 起论，与题干“充分非必要”差一层，照录
\step{①当 $B=\varnothing$ 时，得 $a=-3$，符合题意；$\cdots\cdots\cdots\cdots\cdots$（2 分）}
\step{②当 $B\neq\varnothing$ 时，若 $a>-3$，则 $B=\{x\mid-3<x<a\}$，}
\step{由 $B\subseteq A$，得 $a\leqslant-1$，此时 $-3<a\leqslant-1$；$\cdots\cdots\cdots\cdots\cdots$（4 分）}
\step{若 $a<-3$，则 $B=\{x\mid a<x<-3\}$，此时 $B\subseteq A$ 恒成立，}
\step{故 $a<-3$ 符合题意；$\cdots\cdots\cdots$（6 分）}
\step{综上所述，实数 $a$ 的取值范围为 $a\leqslant-1$. $\cdots\cdots\cdots\cdots\cdots$（8 分）}
\solend
\fi

\item 已知集合 $A=\{x\mid x<-2$ 或 $x\geqslant3\}$，$U=R$.

（1）求 $\overline A$；

（2）若 $B=\{x\mid mx-1>0\}$，当 $B\subseteq A$ 时，求实数 $m$ 的取值范围；

（3）若 $B=\{x\mid2m-6\leqslant x\leqslant m+4\}$，当 $B\cap\overline A=\varnothing$ 时，求实数 $m$ 的取值范围.

\ansspace*[14]

\ifanswer
\solbegin
\stepm{（1）}{因为集合 $A=\{x\mid x<-2$ 或 $x\geqslant3\}$，$U=R$，所以 $\overline A=\{x\mid-2\leqslant x<3\}$；}
\stepm{（2）}{当 $m=0$ 时，$B=\varnothing$，符合题意，}
\step{当 $m>0$ 时，$B=\left\{x\left|x>\dfrac1m\right.\right\}$，则 $\dfrac1m\geqslant3$，解得 $0<m\leqslant\dfrac13$，}
\step{当 $m<0$ 时，$B=\left\{x\left|x<\dfrac1m\right.\right\}$，则 $\dfrac1m\leqslant-2$，解得 $-\dfrac12\leqslant m<0$，}
\step{综上所述，实数 $m$ 的取值范围 $\left[-\dfrac12,\dfrac13\right]$；}
\stepm{（3）}{由题意得 $B=\{x\mid2m-6\leqslant x\leqslant m+4\}$，$\overline A=\{x\mid-2\leqslant x<3\}$，且 $B\cap\overline A=\varnothing$，}
\step{当 $B=\varnothing$ 时，$2m-6>m+4$，解得 $m>10$，}
\step{当 $B\neq\varnothing$ 时，则 $\begin{cases}2m-6\leqslant m+4\\ m+4<-2\end{cases}$ 或 $\begin{cases}2m-6\leqslant m+4\\ 2m-6\geqslant3\end{cases}$，}
\step{解得 $m<-6$ 或 $\dfrac92\leqslant m\leqslant10$，}
\step{综上所述，实数 $m$ 的取值范围为 $(-\infty,-6)\cup\left[\dfrac92,+\infty\right)$.}
\solend
\fi

\item 已知命题 $p$：关于 $x$ 的方程 $x^2+mx+1=0$ 有两个不同的正根；命题 $q$：关于 $x$ 的方程 $4x^2+4(m-2)x=-1$ 有实根.

（1）若命题 $q$ 为假命题，求整数 $m$ 的解集；

（2）若命题 $p$ 为真命题，且关于 $x$ 的方程 $x^2+mx+1=0$ 的两个实数根 $x_1$、$x_2$ 满足 $3x_1+2x_2=5$，求实数 $m$ 的取值集合；

（3）若命题 $p$ 和命题 $q$ 有且仅有一个成立，求实数 $m$ 的取值范围.

\ansspace*[16]

\ifanswer
\solbegin
\stepm{（1）}{命题 $q$：关于 $x$ 的方程 $4x^2+4(m-2)x=-1$ 有实根为假命题，}
\step{即 $4x^2+4(m-2)x+1=0$ 无实根，}
\step{则 $\Delta=16(m-2)^2-16<0$，则 $1<m<3$，则整数 $m$ 的解集为 $\{2\}$；}
\stepm{（2）}{若命题 $p$ 为真命题，且关于 $x$ 的方程 $x^2+mx+1=0$ 的两个实数根 $x_1$、$x_2$}
\step{满足 $3x_1+2x_2=5$，则方程 $x^2+mx+1=0$ 有两个不同的正根，}
\step{则 $\begin{cases}\Delta=m^2-4>0\\ x_1+x_2=-m>0\end{cases}$，则 $m<-2$，又 $\begin{cases}x_1+x_2=-m\\ x_1x_2=1\\ 3x_1+2x_2=5\end{cases}$，则 $\begin{cases}x_1=5+2m\\ x_2=-3m-5\end{cases}$，}
\step{则 $(5+2m)(-3m-5)=1$，得 $m=-2$ 或 $m=-\dfrac{13}6$，}
\step{又因为 $m<-2$，所以 $m=-\dfrac{13}6$，所以实数 $m$ 的取值集合为 $\left\{-\dfrac{13}6\right\}$；}
\stepm{（3）}{由（2）得，当 $p$ 为真时，则有 $m<-2$，}
\step{当 $q$ 为真时，$\Delta=16(m-2)^2-16\geqslant0$，解得 $m\leqslant1$ 或 $m\geqslant3$，}
\step{又因为命题 $p$ 和命题 $q$ 只有一个为真，所以 $\begin{cases}m<-2\\ 1<m<3\end{cases}$ 或 $\begin{cases}m\geqslant-2\\ m\leqslant1\ \text{或}\ m\geqslant3\end{cases}$，}
\step{解得 $-2\leqslant m\leqslant1$ 或 $m\geqslant3$，所以实数 $m$ 的取值范围为 $[-2,1]\cup[3,+\infty)$.}
\solend
\fi

\end{enumerate}

\end{document}
"
% 紧凑版：xelatex "\def\iscompact{1}% 高一95_金汇中学_10月月考.tex
%   —— 高一上数学“2026-2027 年金汇中学高一上 10 月月考”（讲评版）
% 试卷版：xelatex 高一95_金汇中学_10月月考.tex
% 答案版：xelatex "\def\isanswer{1}% 高一95_金汇中学_10月月考.tex
%   —— 高一上数学“2026-2027 年金汇中学高一上 10 月月考”（讲评版）
% 试卷版：xelatex 高一95_金汇中学_10月月考.tex
% 答案版：xelatex "\def\isanswer{1}\input{高一95_金汇中学_10月月考.tex}"
% 紧凑版：xelatex "\def\iscompact{1}\input{高一95_金汇中学_10月月考.tex}"
%
% 原始材料：微信公众号“魔都数学加油站”文章，连载“27学年高一95”；
% 原文标题“27学年高一95——2026-2027年上海市金汇中学高一上10月月考”；
% 原文链接：https://mp.weixin.qq.com/s/H3yH73FTin3V7rxWQ-9fjw。
% 正文为 8 张 PNG 页（`原图/高一95_金汇中学_10月月考/01..08.png`），均 4762×6735，按页序读取。
%
% 篇幅与题号：原卷自第 3 题起编号（第 1、2 题不在原文内）。填空题 3—12（10 题）、
% 选择题 13—16（4 题）、解答题 17—21（5 题），共 19 题。
%
% 讲评版：第 4、5 题只印【答案】（“$m>0$”“$x=0$ 或 $x=1$”），本稿只给 \ans{}；
% 第 13、14 题无【解析】、答案“D”“B”直接印在题干括号内，本稿试卷版按惯例留空、
% 答案版提 \ans{}；其余 15 题印【解析】，其中第 15、16 题的答案在解析末句“故选 $A$／$C$”。
% 第 19 题解析带原卷的分步分值（“（3 分）”“（6 分）”等），照录。全卷无插图。
%
% 副标题：原卷未印副标题。本稿按内容标作“集合与常用逻辑用语、不等式”。
%
% 原卷笔误/存疑（均照录未改，正文对应处另有 % 注）：
%   ① 第 17 题 (1) 解析首句作“**第二个**不等式可化为 $(2x+5)(x-3)\leqslant0$”，而
%      $-2x^2+x+15\geqslant0$ 才是**第一个**不等式；原卷误标，照录；
%   ② 第 19 题 (2) 题干问“$x\in B$ 是 $x\in A$ 的充分非必要条件”，解析却从“$A\cup B=A$，
%      得 $B\subseteq A$”起论（少了“真子集”一层）；原卷照录，结论 $a\leqslant-1$ 未改。
%
% 与原卷的排版差异：
%   ① 原卷每题下直接印【答案】/【解析】，本稿按全稿惯例将答案与解析置于答案版；
%   ② 原卷未印副标题，本稿补出；填空横线统一排 \blankline[4.4em]；
%   ③ 原卷数集记号排斜体 $R$、$N$、$Z$，本稿按全稿惯例改用 $\R$、$\N$、$\Z$；
%   ④ 第 17—21 题的子问用行内“（1）（2）（3）”段，解析中的子问标号用 \stepm；
%   ⑤ 解答题按惯例加 \ansspace；选择题括号排 \mbox{（\quad）}。
%
% 行数核对：逐题数 \step/\stepm 条数，合计 87 条（Q3：2；Q4：0；Q5：0；Q6：1；Q7：2；
% Q8：4；Q9：3；Q10：3；Q11：7；Q12：3；Q13：0；Q14：0；Q15：2；Q16：14；Q17：7；
% Q18：5；Q19：12；Q20：10；Q21：12）。第 4、5 题只印答案，第 13、14 题无解析，
% 均不计入 \step 账。原卷解析的换行按语意合并，未逐行强拆。
%
\documentclass[a4paper,11pt]{article}
\usepackage{common}

\begin{document}

\examtitlest[7]{2026--2027 年金汇中学高一上 10 月月考}{集合与常用逻辑用语、不等式}

\bigsec{一、填空题}

\begin{enumerate}
\setcounter{enumi}{2}

\item 命题：如果 $ab>c$，那么 $a>\dfrac cb$，是\blankline[4.4em]（选填“真”、“假”）命题.

\ifanswer
\solbegin
\step{取 $a=1,b=-2,c=-3$，则 $ab=-2>-3=c$，但 $a=1<\dfrac cb=\dfrac32$，}
\step{所以原命题是假命题.}
\solend
\fi

\item 若“$x=0$”是“$x<m$”的充分条件，则实数 $m$ 的取值范围是\blankline[4.4em].

\ifanswer
\ans{$m>0$}
\fi

\item “$x\neq0$ 且 $x\neq1$”的否定形式为\blankline[4.4em].

\ifanswer
\ans{$x=0$ 或 $x=1$}
\fi

\item 若集合 $M=\{(x,y)\mid x+y=3\}$，$N=\left\{(x,y)\left|y=\dfrac12x\right.\right\}$，则 $M\cap N=\blankline[4.4em]$.

\ifanswer
\solbegin
\step{$x+\dfrac12x=3\Rightarrow x=2,y=1$，所以 $M\cap N=\{(2,1)\}$.}
\solend
\fi

\item 集合 $\{x\in\N\mid0<\lvert x-1\rvert<3\}$ 的真子集个数为\blankline[4.4em].

\ifanswer
\solbegin
\step{因为 $\{x\in\N\mid0<\lvert x-1\rvert<3\}=\{x\in\N\mid-3<x-1<3,x\neq1\}$}
\step{$=\{x\in\N\mid-2<x<4,x\neq1\}=\{0,2,3\}$，所以真子集的个数为 $2^3-1=7$.}
\solend
\fi

\item 设 $A=\{x\mid x^2-3x-4=0\}$，$B=\{x\mid ax-1=0\}$，若 $A\cap B=B$，则实数 $a$ 的值可以为\blankline[4.4em].

\ifanswer
\solbegin
\step{$A=\{-1,4\}$，因为 $A\cap B=B$，所以 $B\subseteq A$，}
\step{①$a=0$ 时，$B=\varnothing$，满足 $B\subseteq A$；}
\step{②$a\neq0$ 时，$B=\left\{\dfrac1a\right\}$，则 $\dfrac1a=-1$ 或 $\dfrac1a=4$，所以 $a=-1$ 或 $\dfrac14$；}
\step{所以 $a$ 的值可以为 $0$、$-1$、$\dfrac14$.}
\solend
\fi

\item 关于 $x$ 的不等式 $ax^2-(a+1)x+1\leqslant0(a<0)$ 的解集为\blankline[4.4em].

\ifanswer
\solbegin
\step{原不等式即 $(ax-1)(x-1)\leqslant0$，因为 $a<0$，所以 $\left(x-\dfrac1a\right)(x-1)\geqslant0$.}
\step{又 $\dfrac1a<0<1$，所以 $x\leqslant\dfrac1a$ 或 $x\geqslant1$，}
\step{故解集为 $\left(-\infty,\dfrac1a\right]\cup[1,+\infty)$.}
\solend
\fi

\item 若不等式 $(a-1)x^2+(2a-2)x-3<0$ 的解集是 $\R$，则 $a$ 的取值范围是\blankline[4.4em].

\ifanswer
\solbegin
\step{当 $a=1$ 时，不等式为 $-3<0$，对任意实数 $x$ 均成立.}
\step{当 $a\neq1$ 时，$\begin{cases}a-1<0\\ \Delta=(2a-2)^2+12(a-1)=4(a-1)(a+2)<0\end{cases}$，解得 $-2<a<1$.}
\step{合并得 $a\in(-2,1]$.}
\solend
\fi

\item 若关于 $x$ 的不等式 $x^2-(m+2)x+2m<0$ 的解集中恰有 $3$ 个整数，则实数 $m$ 的取值范围为\blankline[4.4em].

\ifanswer
\solbegin
\step{将不等式 $x^2-(m+2)x+2m<0$ 因式分解为 $(x-2)(x-m)<0$，}
\step{当 $m<2$ 时，解集为 $(m,2)$.}
\step{要使该区间恰有 $3$ 个整数，这 $3$ 个整数为 $-1,0,1$，故 $m$ 需满足 $-2\leqslant m<-1$.}
\step{当 $m>2$ 时，解集为 $(2,m)$. 要使该区间恰有 $3$ 个整数，}
\step{这 $3$ 个整数为 $3$，$4$，$5$，故 $m$ 需满足 $5<m\leqslant6$.}
\step{当 $m=2$ 时，不等式无解，不满足条件.}
\step{综上，实数 $m$ 的取值范围是 $[-2,-1)\cup(5,6]$.}
\solend
\fi

\item 已知集合 $A=\{x\mid x^2-3x+2\leqslant0\},B=\{y\mid y=x^2-2x+a,x\in\R\}$，若对任何一个 $x\in A$，都有 $x\in B$，则实数 $a$ 的取值范围是\blankline[4.4em].

\ifanswer
\solbegin
\step{$(x-1)(x-2)\leqslant0$，所以 $A=[1,2]$.}
\step{由 $y=(x-1)^2+a-1,x\in\R$，得 $B=[a-1,+\infty)$.}
\step{因此 $A\subseteq B\Leftrightarrow a-1\leqslant1\Leftrightarrow a\leqslant2$，所求范围为 $(-\infty,2]$.}
\solend
\fi

\end{enumerate}

\bigsec{二、选择题}

\begin{enumerate}
\setcounter{enumi}{12}

\item 已知 $a=\sqrt2$，集合 $A=\{x\mid x\leqslant2\}$，则下列表示正确的是\mbox{（\quad）}% 原卷括号内已印“D”，试卷版按惯例留空

\keepwithprev
A. $\{a\}\in A$\qquad B. $a\notin A$\qquad C. $a\subset A$\qquad D. $a\in A$

\ifanswer
\ans{$D$}
\fi

\item 下列四个命题中，真命题是\mbox{（\quad）}% 原卷括号内已印“B”，试卷版按惯例留空

\keepwithprev
A. 若 $a>b$，则 $ac^2>bc^2$\qquad B. 若 $a>b,c>d$，则 $a-d>b-c$

C. 若 $\lvert a\rvert>b$，则 $a^2>b^2$\qquad D. 若 $a>b$，则 $\dfrac1a<\dfrac1b$

\ifanswer
\ans{$B$}
\fi

\item “$x>1$”是“$\dfrac1x<1$”的\mbox{（\quad）}

\keepwithprev
A. 充分不必要条件\qquad B. 必要不充分条件

C. 充要条件\qquad D. 既不充分也不必要条件

\ifanswer
\solbegin
\step{$x>1\Rightarrow0<\dfrac1x<1$，充分性成立.}
\step{当 $x=-1$ 时，$\dfrac1x=-1<1$，但 $x>1$ 不成立，必要性不成立. 故选 $A$.}
\solend
\fi

\item 对于集合 $M=\{a\mid a=x^2-y^2,x\in\Z,y\in\Z\}$，给出如下三个结论：①如果 $P=\{b\mid b=2n+1,n\in\Z\}$，那么 $P\subseteq M$；②如果 $c=4n+2,n\in\Z$，那么 $c\notin M$；③如果 $a_1\in M,a_2\in M$，那么 $a_1a_2\in M$. 其中正确结论的个数是\mbox{（\quad）}

\keepwithprev
A. $1$\qquad B. $2$\qquad C. $3$\qquad D. $0$

\ifanswer
\solbegin
\step{对于①：因为 $2n+1=(n+1)^2-n^2,n\in\Z$，所以 $b=2n+1\in M$，故 $P\subseteq M$，}
\step{故①正确；}
\step{对于②：因为 $c=4n+2=2(2n+1),n\in\Z$，所以 $c$ 为偶数，且不能被 $4$ 整除，}
\step{若 $c\in M$，则存在 $x\in\Z,y\in\Z$，得 $c=x^2-y^2=(x+y)(x-y),x\in\Z,y\in\Z$，}
\step{因为 $x+y$ 和 $x-y$ 同奇或同偶，}
\step{若 $x+y$ 和 $x-y$ 同奇，则 $c=(x+y)(x-y)$ 为奇数，矛盾，不符合，}
\step{若 $x+y$ 和 $x-y$ 同偶，则 $c=(x+y)(x-y)$ 能被 $4$ 整除，矛盾，不符合，}
\step{所以 $c\notin M$，故②正确；}
\step{对于③：因为 $a_1\in M,a_2\in M$，}
\step{所以存在 $x_1,x_2\in\Z,y_1,y_2\in\Z$ 使得 $a_1=x_1^2-y_1^2,a_2=x_2^2-y_2^2$，}
\step{所以 $a_1a_2=(x_1^2-y_1^2)(x_2^2-y_2^2)=(x_1x_2)^2+(y_1y_2)^2-(x_1y_2)^2-(x_2y_1)^2$}
\step{$=(x_1x_2+y_1y_2)^2-(x_1y_2+x_2y_1)^2$，}
\step{因为 $x_1x_2+y_1y_2,x_1y_2+x_2y_1\in\Z$，所以 $a_1a_2\in M$，故③正确.}
\step{故选 $C$.}
\solend
\fi

\end{enumerate}

\bigsec{三、解答题}

\begin{enumerate}
\setcounter{enumi}{16}

\item （1）求不等式组 $\begin{cases}-2x^2+x+15\geqslant0\\ (x+1)(3x-1)>0\end{cases}$ 的解集；

（2）讨论解关于 $x$ 的不等式 $mx+4<m^2+2x$.

\ansspace*[12]

\ifanswer
\solbegin
\stepm{（1）}{第二个不等式可化为 $(2x+5)(x-3)\leqslant0$，解得 $-\dfrac52\leqslant x\leqslant3$.}% 原卷把“第一个”误作“第二个”，照录
\step{第二个不等式 $(x+1)(3x-1)>0$ 的解集为 $(-\infty,-1)\cup\left(\dfrac13,+\infty\right)$.}
\step{原不等式组的解集为 $\left[-\dfrac52,-1\right)\cup\left(\dfrac13,3\right]$.}
\stepm{（2）}{移项并因式分解，得 $(m-2)x<m^2-4\Leftrightarrow(m-2)x<(m-2)(m+2)$.}
\step{当 $m>2$ 时，$m-2>0$，解集为 $(-\infty,m+2)$.}
\step{当 $m<2$ 时，$m-2<0$，解集为 $(m+2,+\infty)$.}
\step{当 $m=2$ 时，不等式化为 $0<0$，不成立，解集为 $\varnothing$.}
\solend
\fi

\item （1）设 $x$ 是实数，比较 $x^2+9$ 与 $6x$ 的值的大小；

（2）已知 $a>0$ 且 $a\neq1$，比较 $\dfrac{a^2+1}{a^2-1}$ 与 $\dfrac{a+1}{a-1}$ 的大小.

\ansspace*[11]

\ifanswer
\solbegin
\stepm{（1）}{作差得 $(x^2+9)-6x=(x-3)^2\geqslant0$，因此 $x^2+9\geqslant6x$；}
\stepm{（2）}{$\dfrac{a^2+1}{a^2-1}-\dfrac{a+1}{a-1}=\dfrac{-2a}{a^2-1}$，}
\step{当 $a>1$ 时，$-2a<0$，$a^2-1>0$，则 $\dfrac{-2a}{a^2-1}<0$，即 $\dfrac{a^2+1}{a^2-1}<\dfrac{a+1}{a-1}$；}
\step{当 $0<a<1$ 时，$-2a<0$，$a^2-1<0$，则 $\dfrac{-2a}{a^2-1}>0$，即 $\dfrac{a^2+1}{a^2-1}>\dfrac{a+1}{a-1}$；}
\step{综上，$a>1$ 时，$\dfrac{a^2+1}{a^2-1}<\dfrac{a+1}{a-1}$；$0<a<1$ 时，$\dfrac{a^2+1}{a^2-1}>\dfrac{a+1}{a-1}$.}
\solend
\fi

\item 已知 $a\in\R$，集合 $A=\left\{x\left|\dfrac{2-x}{x+1}\leqslant0\right.\right\},B=\{x\mid x^2+(3-a)x-3a<0\}$.

（1）当 $a=\dfrac13$ 时，求 $A$ 和 $B$；

（2）“$x\in B$”是“$x\in A$”的充分非必要条件，求实数 $a$ 的取值范围.

\ansspace*[13]

\ifanswer
\solbegin
\stepm{（1）}{由 $\dfrac{2-x}{x+1}\leqslant0$，得 $\begin{cases}(2-x)(x+1)\leqslant0\\ x\neq-1\end{cases}$，解得 $x<-1$ 或 $x\geqslant2$，}
\step{则 $A=\{x\mid x<-1$ 或 $x\geqslant2\}$；$\cdots\cdots\cdots\cdots\cdots\cdots\cdots\cdots\cdots\cdots\cdots\cdots\cdots\cdots\cdots\cdots$（3 分）}
\step{$B=\{x\mid x^2+(3-a)x-3a<0\}=\{x\mid(x+3)(x-a)<0\}$，}
\step{当 $a=\dfrac13$ 时，$B=\left\{x\left|(x+2)\left(x-\dfrac13\right)<0\right.\right\}$，}
\step{解得 $B=\left\{x\left|-3<x<\dfrac13\right.\right\}\cdots\cdots\cdots\cdots\cdots\cdots$（6 分）}
\stepm{（2）}{由 $A\cup B=A$，得 $B\subseteq A$，}% 原卷从 A∪B=A 起论，与题干“充分非必要”差一层，照录
\step{①当 $B=\varnothing$ 时，得 $a=-3$，符合题意；$\cdots\cdots\cdots\cdots\cdots$（2 分）}
\step{②当 $B\neq\varnothing$ 时，若 $a>-3$，则 $B=\{x\mid-3<x<a\}$，}
\step{由 $B\subseteq A$，得 $a\leqslant-1$，此时 $-3<a\leqslant-1$；$\cdots\cdots\cdots\cdots\cdots$（4 分）}
\step{若 $a<-3$，则 $B=\{x\mid a<x<-3\}$，此时 $B\subseteq A$ 恒成立，}
\step{故 $a<-3$ 符合题意；$\cdots\cdots\cdots$（6 分）}
\step{综上所述，实数 $a$ 的取值范围为 $a\leqslant-1$. $\cdots\cdots\cdots\cdots\cdots$（8 分）}
\solend
\fi

\item 已知集合 $A=\{x\mid x<-2$ 或 $x\geqslant3\}$，$U=R$.

（1）求 $\overline A$；

（2）若 $B=\{x\mid mx-1>0\}$，当 $B\subseteq A$ 时，求实数 $m$ 的取值范围；

（3）若 $B=\{x\mid2m-6\leqslant x\leqslant m+4\}$，当 $B\cap\overline A=\varnothing$ 时，求实数 $m$ 的取值范围.

\ansspace*[14]

\ifanswer
\solbegin
\stepm{（1）}{因为集合 $A=\{x\mid x<-2$ 或 $x\geqslant3\}$，$U=R$，所以 $\overline A=\{x\mid-2\leqslant x<3\}$；}
\stepm{（2）}{当 $m=0$ 时，$B=\varnothing$，符合题意，}
\step{当 $m>0$ 时，$B=\left\{x\left|x>\dfrac1m\right.\right\}$，则 $\dfrac1m\geqslant3$，解得 $0<m\leqslant\dfrac13$，}
\step{当 $m<0$ 时，$B=\left\{x\left|x<\dfrac1m\right.\right\}$，则 $\dfrac1m\leqslant-2$，解得 $-\dfrac12\leqslant m<0$，}
\step{综上所述，实数 $m$ 的取值范围 $\left[-\dfrac12,\dfrac13\right]$；}
\stepm{（3）}{由题意得 $B=\{x\mid2m-6\leqslant x\leqslant m+4\}$，$\overline A=\{x\mid-2\leqslant x<3\}$，且 $B\cap\overline A=\varnothing$，}
\step{当 $B=\varnothing$ 时，$2m-6>m+4$，解得 $m>10$，}
\step{当 $B\neq\varnothing$ 时，则 $\begin{cases}2m-6\leqslant m+4\\ m+4<-2\end{cases}$ 或 $\begin{cases}2m-6\leqslant m+4\\ 2m-6\geqslant3\end{cases}$，}
\step{解得 $m<-6$ 或 $\dfrac92\leqslant m\leqslant10$，}
\step{综上所述，实数 $m$ 的取值范围为 $(-\infty,-6)\cup\left[\dfrac92,+\infty\right)$.}
\solend
\fi

\item 已知命题 $p$：关于 $x$ 的方程 $x^2+mx+1=0$ 有两个不同的正根；命题 $q$：关于 $x$ 的方程 $4x^2+4(m-2)x=-1$ 有实根.

（1）若命题 $q$ 为假命题，求整数 $m$ 的解集；

（2）若命题 $p$ 为真命题，且关于 $x$ 的方程 $x^2+mx+1=0$ 的两个实数根 $x_1$、$x_2$ 满足 $3x_1+2x_2=5$，求实数 $m$ 的取值集合；

（3）若命题 $p$ 和命题 $q$ 有且仅有一个成立，求实数 $m$ 的取值范围.

\ansspace*[16]

\ifanswer
\solbegin
\stepm{（1）}{命题 $q$：关于 $x$ 的方程 $4x^2+4(m-2)x=-1$ 有实根为假命题，}
\step{即 $4x^2+4(m-2)x+1=0$ 无实根，}
\step{则 $\Delta=16(m-2)^2-16<0$，则 $1<m<3$，则整数 $m$ 的解集为 $\{2\}$；}
\stepm{（2）}{若命题 $p$ 为真命题，且关于 $x$ 的方程 $x^2+mx+1=0$ 的两个实数根 $x_1$、$x_2$}
\step{满足 $3x_1+2x_2=5$，则方程 $x^2+mx+1=0$ 有两个不同的正根，}
\step{则 $\begin{cases}\Delta=m^2-4>0\\ x_1+x_2=-m>0\end{cases}$，则 $m<-2$，又 $\begin{cases}x_1+x_2=-m\\ x_1x_2=1\\ 3x_1+2x_2=5\end{cases}$，则 $\begin{cases}x_1=5+2m\\ x_2=-3m-5\end{cases}$，}
\step{则 $(5+2m)(-3m-5)=1$，得 $m=-2$ 或 $m=-\dfrac{13}6$，}
\step{又因为 $m<-2$，所以 $m=-\dfrac{13}6$，所以实数 $m$ 的取值集合为 $\left\{-\dfrac{13}6\right\}$；}
\stepm{（3）}{由（2）得，当 $p$ 为真时，则有 $m<-2$，}
\step{当 $q$ 为真时，$\Delta=16(m-2)^2-16\geqslant0$，解得 $m\leqslant1$ 或 $m\geqslant3$，}
\step{又因为命题 $p$ 和命题 $q$ 只有一个为真，所以 $\begin{cases}m<-2\\ 1<m<3\end{cases}$ 或 $\begin{cases}m\geqslant-2\\ m\leqslant1\ \text{或}\ m\geqslant3\end{cases}$，}
\step{解得 $-2\leqslant m\leqslant1$ 或 $m\geqslant3$，所以实数 $m$ 的取值范围为 $[-2,1]\cup[3,+\infty)$.}
\solend
\fi

\end{enumerate}

\end{document}
"
% 紧凑版：xelatex "\def\iscompact{1}% 高一95_金汇中学_10月月考.tex
%   —— 高一上数学“2026-2027 年金汇中学高一上 10 月月考”（讲评版）
% 试卷版：xelatex 高一95_金汇中学_10月月考.tex
% 答案版：xelatex "\def\isanswer{1}\input{高一95_金汇中学_10月月考.tex}"
% 紧凑版：xelatex "\def\iscompact{1}\input{高一95_金汇中学_10月月考.tex}"
%
% 原始材料：微信公众号“魔都数学加油站”文章，连载“27学年高一95”；
% 原文标题“27学年高一95——2026-2027年上海市金汇中学高一上10月月考”；
% 原文链接：https://mp.weixin.qq.com/s/H3yH73FTin3V7rxWQ-9fjw。
% 正文为 8 张 PNG 页（`原图/高一95_金汇中学_10月月考/01..08.png`），均 4762×6735，按页序读取。
%
% 篇幅与题号：原卷自第 3 题起编号（第 1、2 题不在原文内）。填空题 3—12（10 题）、
% 选择题 13—16（4 题）、解答题 17—21（5 题），共 19 题。
%
% 讲评版：第 4、5 题只印【答案】（“$m>0$”“$x=0$ 或 $x=1$”），本稿只给 \ans{}；
% 第 13、14 题无【解析】、答案“D”“B”直接印在题干括号内，本稿试卷版按惯例留空、
% 答案版提 \ans{}；其余 15 题印【解析】，其中第 15、16 题的答案在解析末句“故选 $A$／$C$”。
% 第 19 题解析带原卷的分步分值（“（3 分）”“（6 分）”等），照录。全卷无插图。
%
% 副标题：原卷未印副标题。本稿按内容标作“集合与常用逻辑用语、不等式”。
%
% 原卷笔误/存疑（均照录未改，正文对应处另有 % 注）：
%   ① 第 17 题 (1) 解析首句作“**第二个**不等式可化为 $(2x+5)(x-3)\leqslant0$”，而
%      $-2x^2+x+15\geqslant0$ 才是**第一个**不等式；原卷误标，照录；
%   ② 第 19 题 (2) 题干问“$x\in B$ 是 $x\in A$ 的充分非必要条件”，解析却从“$A\cup B=A$，
%      得 $B\subseteq A$”起论（少了“真子集”一层）；原卷照录，结论 $a\leqslant-1$ 未改。
%
% 与原卷的排版差异：
%   ① 原卷每题下直接印【答案】/【解析】，本稿按全稿惯例将答案与解析置于答案版；
%   ② 原卷未印副标题，本稿补出；填空横线统一排 \blankline[4.4em]；
%   ③ 原卷数集记号排斜体 $R$、$N$、$Z$，本稿按全稿惯例改用 $\R$、$\N$、$\Z$；
%   ④ 第 17—21 题的子问用行内“（1）（2）（3）”段，解析中的子问标号用 \stepm；
%   ⑤ 解答题按惯例加 \ansspace；选择题括号排 \mbox{（\quad）}。
%
% 行数核对：逐题数 \step/\stepm 条数，合计 87 条（Q3：2；Q4：0；Q5：0；Q6：1；Q7：2；
% Q8：4；Q9：3；Q10：3；Q11：7；Q12：3；Q13：0；Q14：0；Q15：2；Q16：14；Q17：7；
% Q18：5；Q19：12；Q20：10；Q21：12）。第 4、5 题只印答案，第 13、14 题无解析，
% 均不计入 \step 账。原卷解析的换行按语意合并，未逐行强拆。
%
\documentclass[a4paper,11pt]{article}
\usepackage{common}

\begin{document}

\examtitlest[7]{2026--2027 年金汇中学高一上 10 月月考}{集合与常用逻辑用语、不等式}

\bigsec{一、填空题}

\begin{enumerate}
\setcounter{enumi}{2}

\item 命题：如果 $ab>c$，那么 $a>\dfrac cb$，是\blankline[4.4em]（选填“真”、“假”）命题.

\ifanswer
\solbegin
\step{取 $a=1,b=-2,c=-3$，则 $ab=-2>-3=c$，但 $a=1<\dfrac cb=\dfrac32$，}
\step{所以原命题是假命题.}
\solend
\fi

\item 若“$x=0$”是“$x<m$”的充分条件，则实数 $m$ 的取值范围是\blankline[4.4em].

\ifanswer
\ans{$m>0$}
\fi

\item “$x\neq0$ 且 $x\neq1$”的否定形式为\blankline[4.4em].

\ifanswer
\ans{$x=0$ 或 $x=1$}
\fi

\item 若集合 $M=\{(x,y)\mid x+y=3\}$，$N=\left\{(x,y)\left|y=\dfrac12x\right.\right\}$，则 $M\cap N=\blankline[4.4em]$.

\ifanswer
\solbegin
\step{$x+\dfrac12x=3\Rightarrow x=2,y=1$，所以 $M\cap N=\{(2,1)\}$.}
\solend
\fi

\item 集合 $\{x\in\N\mid0<\lvert x-1\rvert<3\}$ 的真子集个数为\blankline[4.4em].

\ifanswer
\solbegin
\step{因为 $\{x\in\N\mid0<\lvert x-1\rvert<3\}=\{x\in\N\mid-3<x-1<3,x\neq1\}$}
\step{$=\{x\in\N\mid-2<x<4,x\neq1\}=\{0,2,3\}$，所以真子集的个数为 $2^3-1=7$.}
\solend
\fi

\item 设 $A=\{x\mid x^2-3x-4=0\}$，$B=\{x\mid ax-1=0\}$，若 $A\cap B=B$，则实数 $a$ 的值可以为\blankline[4.4em].

\ifanswer
\solbegin
\step{$A=\{-1,4\}$，因为 $A\cap B=B$，所以 $B\subseteq A$，}
\step{①$a=0$ 时，$B=\varnothing$，满足 $B\subseteq A$；}
\step{②$a\neq0$ 时，$B=\left\{\dfrac1a\right\}$，则 $\dfrac1a=-1$ 或 $\dfrac1a=4$，所以 $a=-1$ 或 $\dfrac14$；}
\step{所以 $a$ 的值可以为 $0$、$-1$、$\dfrac14$.}
\solend
\fi

\item 关于 $x$ 的不等式 $ax^2-(a+1)x+1\leqslant0(a<0)$ 的解集为\blankline[4.4em].

\ifanswer
\solbegin
\step{原不等式即 $(ax-1)(x-1)\leqslant0$，因为 $a<0$，所以 $\left(x-\dfrac1a\right)(x-1)\geqslant0$.}
\step{又 $\dfrac1a<0<1$，所以 $x\leqslant\dfrac1a$ 或 $x\geqslant1$，}
\step{故解集为 $\left(-\infty,\dfrac1a\right]\cup[1,+\infty)$.}
\solend
\fi

\item 若不等式 $(a-1)x^2+(2a-2)x-3<0$ 的解集是 $\R$，则 $a$ 的取值范围是\blankline[4.4em].

\ifanswer
\solbegin
\step{当 $a=1$ 时，不等式为 $-3<0$，对任意实数 $x$ 均成立.}
\step{当 $a\neq1$ 时，$\begin{cases}a-1<0\\ \Delta=(2a-2)^2+12(a-1)=4(a-1)(a+2)<0\end{cases}$，解得 $-2<a<1$.}
\step{合并得 $a\in(-2,1]$.}
\solend
\fi

\item 若关于 $x$ 的不等式 $x^2-(m+2)x+2m<0$ 的解集中恰有 $3$ 个整数，则实数 $m$ 的取值范围为\blankline[4.4em].

\ifanswer
\solbegin
\step{将不等式 $x^2-(m+2)x+2m<0$ 因式分解为 $(x-2)(x-m)<0$，}
\step{当 $m<2$ 时，解集为 $(m,2)$.}
\step{要使该区间恰有 $3$ 个整数，这 $3$ 个整数为 $-1,0,1$，故 $m$ 需满足 $-2\leqslant m<-1$.}
\step{当 $m>2$ 时，解集为 $(2,m)$. 要使该区间恰有 $3$ 个整数，}
\step{这 $3$ 个整数为 $3$，$4$，$5$，故 $m$ 需满足 $5<m\leqslant6$.}
\step{当 $m=2$ 时，不等式无解，不满足条件.}
\step{综上，实数 $m$ 的取值范围是 $[-2,-1)\cup(5,6]$.}
\solend
\fi

\item 已知集合 $A=\{x\mid x^2-3x+2\leqslant0\},B=\{y\mid y=x^2-2x+a,x\in\R\}$，若对任何一个 $x\in A$，都有 $x\in B$，则实数 $a$ 的取值范围是\blankline[4.4em].

\ifanswer
\solbegin
\step{$(x-1)(x-2)\leqslant0$，所以 $A=[1,2]$.}
\step{由 $y=(x-1)^2+a-1,x\in\R$，得 $B=[a-1,+\infty)$.}
\step{因此 $A\subseteq B\Leftrightarrow a-1\leqslant1\Leftrightarrow a\leqslant2$，所求范围为 $(-\infty,2]$.}
\solend
\fi

\end{enumerate}

\bigsec{二、选择题}

\begin{enumerate}
\setcounter{enumi}{12}

\item 已知 $a=\sqrt2$，集合 $A=\{x\mid x\leqslant2\}$，则下列表示正确的是\mbox{（\quad）}% 原卷括号内已印“D”，试卷版按惯例留空

\keepwithprev
A. $\{a\}\in A$\qquad B. $a\notin A$\qquad C. $a\subset A$\qquad D. $a\in A$

\ifanswer
\ans{$D$}
\fi

\item 下列四个命题中，真命题是\mbox{（\quad）}% 原卷括号内已印“B”，试卷版按惯例留空

\keepwithprev
A. 若 $a>b$，则 $ac^2>bc^2$\qquad B. 若 $a>b,c>d$，则 $a-d>b-c$

C. 若 $\lvert a\rvert>b$，则 $a^2>b^2$\qquad D. 若 $a>b$，则 $\dfrac1a<\dfrac1b$

\ifanswer
\ans{$B$}
\fi

\item “$x>1$”是“$\dfrac1x<1$”的\mbox{（\quad）}

\keepwithprev
A. 充分不必要条件\qquad B. 必要不充分条件

C. 充要条件\qquad D. 既不充分也不必要条件

\ifanswer
\solbegin
\step{$x>1\Rightarrow0<\dfrac1x<1$，充分性成立.}
\step{当 $x=-1$ 时，$\dfrac1x=-1<1$，但 $x>1$ 不成立，必要性不成立. 故选 $A$.}
\solend
\fi

\item 对于集合 $M=\{a\mid a=x^2-y^2,x\in\Z,y\in\Z\}$，给出如下三个结论：①如果 $P=\{b\mid b=2n+1,n\in\Z\}$，那么 $P\subseteq M$；②如果 $c=4n+2,n\in\Z$，那么 $c\notin M$；③如果 $a_1\in M,a_2\in M$，那么 $a_1a_2\in M$. 其中正确结论的个数是\mbox{（\quad）}

\keepwithprev
A. $1$\qquad B. $2$\qquad C. $3$\qquad D. $0$

\ifanswer
\solbegin
\step{对于①：因为 $2n+1=(n+1)^2-n^2,n\in\Z$，所以 $b=2n+1\in M$，故 $P\subseteq M$，}
\step{故①正确；}
\step{对于②：因为 $c=4n+2=2(2n+1),n\in\Z$，所以 $c$ 为偶数，且不能被 $4$ 整除，}
\step{若 $c\in M$，则存在 $x\in\Z,y\in\Z$，得 $c=x^2-y^2=(x+y)(x-y),x\in\Z,y\in\Z$，}
\step{因为 $x+y$ 和 $x-y$ 同奇或同偶，}
\step{若 $x+y$ 和 $x-y$ 同奇，则 $c=(x+y)(x-y)$ 为奇数，矛盾，不符合，}
\step{若 $x+y$ 和 $x-y$ 同偶，则 $c=(x+y)(x-y)$ 能被 $4$ 整除，矛盾，不符合，}
\step{所以 $c\notin M$，故②正确；}
\step{对于③：因为 $a_1\in M,a_2\in M$，}
\step{所以存在 $x_1,x_2\in\Z,y_1,y_2\in\Z$ 使得 $a_1=x_1^2-y_1^2,a_2=x_2^2-y_2^2$，}
\step{所以 $a_1a_2=(x_1^2-y_1^2)(x_2^2-y_2^2)=(x_1x_2)^2+(y_1y_2)^2-(x_1y_2)^2-(x_2y_1)^2$}
\step{$=(x_1x_2+y_1y_2)^2-(x_1y_2+x_2y_1)^2$，}
\step{因为 $x_1x_2+y_1y_2,x_1y_2+x_2y_1\in\Z$，所以 $a_1a_2\in M$，故③正确.}
\step{故选 $C$.}
\solend
\fi

\end{enumerate}

\bigsec{三、解答题}

\begin{enumerate}
\setcounter{enumi}{16}

\item （1）求不等式组 $\begin{cases}-2x^2+x+15\geqslant0\\ (x+1)(3x-1)>0\end{cases}$ 的解集；

（2）讨论解关于 $x$ 的不等式 $mx+4<m^2+2x$.

\ansspace*[12]

\ifanswer
\solbegin
\stepm{（1）}{第二个不等式可化为 $(2x+5)(x-3)\leqslant0$，解得 $-\dfrac52\leqslant x\leqslant3$.}% 原卷把“第一个”误作“第二个”，照录
\step{第二个不等式 $(x+1)(3x-1)>0$ 的解集为 $(-\infty,-1)\cup\left(\dfrac13,+\infty\right)$.}
\step{原不等式组的解集为 $\left[-\dfrac52,-1\right)\cup\left(\dfrac13,3\right]$.}
\stepm{（2）}{移项并因式分解，得 $(m-2)x<m^2-4\Leftrightarrow(m-2)x<(m-2)(m+2)$.}
\step{当 $m>2$ 时，$m-2>0$，解集为 $(-\infty,m+2)$.}
\step{当 $m<2$ 时，$m-2<0$，解集为 $(m+2,+\infty)$.}
\step{当 $m=2$ 时，不等式化为 $0<0$，不成立，解集为 $\varnothing$.}
\solend
\fi

\item （1）设 $x$ 是实数，比较 $x^2+9$ 与 $6x$ 的值的大小；

（2）已知 $a>0$ 且 $a\neq1$，比较 $\dfrac{a^2+1}{a^2-1}$ 与 $\dfrac{a+1}{a-1}$ 的大小.

\ansspace*[11]

\ifanswer
\solbegin
\stepm{（1）}{作差得 $(x^2+9)-6x=(x-3)^2\geqslant0$，因此 $x^2+9\geqslant6x$；}
\stepm{（2）}{$\dfrac{a^2+1}{a^2-1}-\dfrac{a+1}{a-1}=\dfrac{-2a}{a^2-1}$，}
\step{当 $a>1$ 时，$-2a<0$，$a^2-1>0$，则 $\dfrac{-2a}{a^2-1}<0$，即 $\dfrac{a^2+1}{a^2-1}<\dfrac{a+1}{a-1}$；}
\step{当 $0<a<1$ 时，$-2a<0$，$a^2-1<0$，则 $\dfrac{-2a}{a^2-1}>0$，即 $\dfrac{a^2+1}{a^2-1}>\dfrac{a+1}{a-1}$；}
\step{综上，$a>1$ 时，$\dfrac{a^2+1}{a^2-1}<\dfrac{a+1}{a-1}$；$0<a<1$ 时，$\dfrac{a^2+1}{a^2-1}>\dfrac{a+1}{a-1}$.}
\solend
\fi

\item 已知 $a\in\R$，集合 $A=\left\{x\left|\dfrac{2-x}{x+1}\leqslant0\right.\right\},B=\{x\mid x^2+(3-a)x-3a<0\}$.

（1）当 $a=\dfrac13$ 时，求 $A$ 和 $B$；

（2）“$x\in B$”是“$x\in A$”的充分非必要条件，求实数 $a$ 的取值范围.

\ansspace*[13]

\ifanswer
\solbegin
\stepm{（1）}{由 $\dfrac{2-x}{x+1}\leqslant0$，得 $\begin{cases}(2-x)(x+1)\leqslant0\\ x\neq-1\end{cases}$，解得 $x<-1$ 或 $x\geqslant2$，}
\step{则 $A=\{x\mid x<-1$ 或 $x\geqslant2\}$；$\cdots\cdots\cdots\cdots\cdots\cdots\cdots\cdots\cdots\cdots\cdots\cdots\cdots\cdots\cdots\cdots$（3 分）}
\step{$B=\{x\mid x^2+(3-a)x-3a<0\}=\{x\mid(x+3)(x-a)<0\}$，}
\step{当 $a=\dfrac13$ 时，$B=\left\{x\left|(x+2)\left(x-\dfrac13\right)<0\right.\right\}$，}
\step{解得 $B=\left\{x\left|-3<x<\dfrac13\right.\right\}\cdots\cdots\cdots\cdots\cdots\cdots$（6 分）}
\stepm{（2）}{由 $A\cup B=A$，得 $B\subseteq A$，}% 原卷从 A∪B=A 起论，与题干“充分非必要”差一层，照录
\step{①当 $B=\varnothing$ 时，得 $a=-3$，符合题意；$\cdots\cdots\cdots\cdots\cdots$（2 分）}
\step{②当 $B\neq\varnothing$ 时，若 $a>-3$，则 $B=\{x\mid-3<x<a\}$，}
\step{由 $B\subseteq A$，得 $a\leqslant-1$，此时 $-3<a\leqslant-1$；$\cdots\cdots\cdots\cdots\cdots$（4 分）}
\step{若 $a<-3$，则 $B=\{x\mid a<x<-3\}$，此时 $B\subseteq A$ 恒成立，}
\step{故 $a<-3$ 符合题意；$\cdots\cdots\cdots$（6 分）}
\step{综上所述，实数 $a$ 的取值范围为 $a\leqslant-1$. $\cdots\cdots\cdots\cdots\cdots$（8 分）}
\solend
\fi

\item 已知集合 $A=\{x\mid x<-2$ 或 $x\geqslant3\}$，$U=R$.

（1）求 $\overline A$；

（2）若 $B=\{x\mid mx-1>0\}$，当 $B\subseteq A$ 时，求实数 $m$ 的取值范围；

（3）若 $B=\{x\mid2m-6\leqslant x\leqslant m+4\}$，当 $B\cap\overline A=\varnothing$ 时，求实数 $m$ 的取值范围.

\ansspace*[14]

\ifanswer
\solbegin
\stepm{（1）}{因为集合 $A=\{x\mid x<-2$ 或 $x\geqslant3\}$，$U=R$，所以 $\overline A=\{x\mid-2\leqslant x<3\}$；}
\stepm{（2）}{当 $m=0$ 时，$B=\varnothing$，符合题意，}
\step{当 $m>0$ 时，$B=\left\{x\left|x>\dfrac1m\right.\right\}$，则 $\dfrac1m\geqslant3$，解得 $0<m\leqslant\dfrac13$，}
\step{当 $m<0$ 时，$B=\left\{x\left|x<\dfrac1m\right.\right\}$，则 $\dfrac1m\leqslant-2$，解得 $-\dfrac12\leqslant m<0$，}
\step{综上所述，实数 $m$ 的取值范围 $\left[-\dfrac12,\dfrac13\right]$；}
\stepm{（3）}{由题意得 $B=\{x\mid2m-6\leqslant x\leqslant m+4\}$，$\overline A=\{x\mid-2\leqslant x<3\}$，且 $B\cap\overline A=\varnothing$，}
\step{当 $B=\varnothing$ 时，$2m-6>m+4$，解得 $m>10$，}
\step{当 $B\neq\varnothing$ 时，则 $\begin{cases}2m-6\leqslant m+4\\ m+4<-2\end{cases}$ 或 $\begin{cases}2m-6\leqslant m+4\\ 2m-6\geqslant3\end{cases}$，}
\step{解得 $m<-6$ 或 $\dfrac92\leqslant m\leqslant10$，}
\step{综上所述，实数 $m$ 的取值范围为 $(-\infty,-6)\cup\left[\dfrac92,+\infty\right)$.}
\solend
\fi

\item 已知命题 $p$：关于 $x$ 的方程 $x^2+mx+1=0$ 有两个不同的正根；命题 $q$：关于 $x$ 的方程 $4x^2+4(m-2)x=-1$ 有实根.

（1）若命题 $q$ 为假命题，求整数 $m$ 的解集；

（2）若命题 $p$ 为真命题，且关于 $x$ 的方程 $x^2+mx+1=0$ 的两个实数根 $x_1$、$x_2$ 满足 $3x_1+2x_2=5$，求实数 $m$ 的取值集合；

（3）若命题 $p$ 和命题 $q$ 有且仅有一个成立，求实数 $m$ 的取值范围.

\ansspace*[16]

\ifanswer
\solbegin
\stepm{（1）}{命题 $q$：关于 $x$ 的方程 $4x^2+4(m-2)x=-1$ 有实根为假命题，}
\step{即 $4x^2+4(m-2)x+1=0$ 无实根，}
\step{则 $\Delta=16(m-2)^2-16<0$，则 $1<m<3$，则整数 $m$ 的解集为 $\{2\}$；}
\stepm{（2）}{若命题 $p$ 为真命题，且关于 $x$ 的方程 $x^2+mx+1=0$ 的两个实数根 $x_1$、$x_2$}
\step{满足 $3x_1+2x_2=5$，则方程 $x^2+mx+1=0$ 有两个不同的正根，}
\step{则 $\begin{cases}\Delta=m^2-4>0\\ x_1+x_2=-m>0\end{cases}$，则 $m<-2$，又 $\begin{cases}x_1+x_2=-m\\ x_1x_2=1\\ 3x_1+2x_2=5\end{cases}$，则 $\begin{cases}x_1=5+2m\\ x_2=-3m-5\end{cases}$，}
\step{则 $(5+2m)(-3m-5)=1$，得 $m=-2$ 或 $m=-\dfrac{13}6$，}
\step{又因为 $m<-2$，所以 $m=-\dfrac{13}6$，所以实数 $m$ 的取值集合为 $\left\{-\dfrac{13}6\right\}$；}
\stepm{（3）}{由（2）得，当 $p$ 为真时，则有 $m<-2$，}
\step{当 $q$ 为真时，$\Delta=16(m-2)^2-16\geqslant0$，解得 $m\leqslant1$ 或 $m\geqslant3$，}
\step{又因为命题 $p$ 和命题 $q$ 只有一个为真，所以 $\begin{cases}m<-2\\ 1<m<3\end{cases}$ 或 $\begin{cases}m\geqslant-2\\ m\leqslant1\ \text{或}\ m\geqslant3\end{cases}$，}
\step{解得 $-2\leqslant m\leqslant1$ 或 $m\geqslant3$，所以实数 $m$ 的取值范围为 $[-2,1]\cup[3,+\infty)$.}
\solend
\fi

\end{enumerate}

\end{document}
"
%
% 原始材料：微信公众号“魔都数学加油站”文章，连载“27学年高一95”；
% 原文标题“27学年高一95——2026-2027年上海市金汇中学高一上10月月考”；
% 原文链接：https://mp.weixin.qq.com/s/H3yH73FTin3V7rxWQ-9fjw。
% 正文为 8 张 PNG 页（`原图/高一95_金汇中学_10月月考/01..08.png`），均 4762×6735，按页序读取。
%
% 篇幅与题号：原卷自第 3 题起编号（第 1、2 题不在原文内）。填空题 3—12（10 题）、
% 选择题 13—16（4 题）、解答题 17—21（5 题），共 19 题。
%
% 讲评版：第 4、5 题只印【答案】（“$m>0$”“$x=0$ 或 $x=1$”），本稿只给 \ans{}；
% 第 13、14 题无【解析】、答案“D”“B”直接印在题干括号内，本稿试卷版按惯例留空、
% 答案版提 \ans{}；其余 15 题印【解析】，其中第 15、16 题的答案在解析末句“故选 $A$／$C$”。
% 第 19 题解析带原卷的分步分值（“（3 分）”“（6 分）”等），照录。全卷无插图。
%
% 副标题：原卷未印副标题。本稿按内容标作“集合与常用逻辑用语、不等式”。
%
% 原卷笔误/存疑（均照录未改，正文对应处另有 % 注）：
%   ① 第 17 题 (1) 解析首句作“**第二个**不等式可化为 $(2x+5)(x-3)\leqslant0$”，而
%      $-2x^2+x+15\geqslant0$ 才是**第一个**不等式；原卷误标，照录；
%   ② 第 19 题 (2) 题干问“$x\in B$ 是 $x\in A$ 的充分非必要条件”，解析却从“$A\cup B=A$，
%      得 $B\subseteq A$”起论（少了“真子集”一层）；原卷照录，结论 $a\leqslant-1$ 未改。
%
% 与原卷的排版差异：
%   ① 原卷每题下直接印【答案】/【解析】，本稿按全稿惯例将答案与解析置于答案版；
%   ② 原卷未印副标题，本稿补出；填空横线统一排 \blankline[4.4em]；
%   ③ 原卷数集记号排斜体 $R$、$N$、$Z$，本稿按全稿惯例改用 $\R$、$\N$、$\Z$；
%   ④ 第 17—21 题的子问用行内“（1）（2）（3）”段，解析中的子问标号用 \stepm；
%   ⑤ 解答题按惯例加 \ansspace；选择题括号排 \mbox{（\quad）}。
%
% 行数核对：逐题数 \step/\stepm 条数，合计 87 条（Q3：2；Q4：0；Q5：0；Q6：1；Q7：2；
% Q8：4；Q9：3；Q10：3；Q11：7；Q12：3；Q13：0；Q14：0；Q15：2；Q16：14；Q17：7；
% Q18：5；Q19：12；Q20：10；Q21：12）。第 4、5 题只印答案，第 13、14 题无解析，
% 均不计入 \step 账。原卷解析的换行按语意合并，未逐行强拆。
%
\documentclass[a4paper,11pt]{article}
\usepackage{common}

\begin{document}

\examtitlest[7]{2026--2027 年金汇中学高一上 10 月月考}{集合与常用逻辑用语、不等式}

\bigsec{一、填空题}

\begin{enumerate}
\setcounter{enumi}{2}

\item 命题：如果 $ab>c$，那么 $a>\dfrac cb$，是\blankline[4.4em]（选填“真”、“假”）命题.

\ifanswer
\solbegin
\step{取 $a=1,b=-2,c=-3$，则 $ab=-2>-3=c$，但 $a=1<\dfrac cb=\dfrac32$，}
\step{所以原命题是假命题.}
\solend
\fi

\item 若“$x=0$”是“$x<m$”的充分条件，则实数 $m$ 的取值范围是\blankline[4.4em].

\ifanswer
\ans{$m>0$}
\fi

\item “$x\neq0$ 且 $x\neq1$”的否定形式为\blankline[4.4em].

\ifanswer
\ans{$x=0$ 或 $x=1$}
\fi

\item 若集合 $M=\{(x,y)\mid x+y=3\}$，$N=\left\{(x,y)\left|y=\dfrac12x\right.\right\}$，则 $M\cap N=\blankline[4.4em]$.

\ifanswer
\solbegin
\step{$x+\dfrac12x=3\Rightarrow x=2,y=1$，所以 $M\cap N=\{(2,1)\}$.}
\solend
\fi

\item 集合 $\{x\in\N\mid0<\lvert x-1\rvert<3\}$ 的真子集个数为\blankline[4.4em].

\ifanswer
\solbegin
\step{因为 $\{x\in\N\mid0<\lvert x-1\rvert<3\}=\{x\in\N\mid-3<x-1<3,x\neq1\}$}
\step{$=\{x\in\N\mid-2<x<4,x\neq1\}=\{0,2,3\}$，所以真子集的个数为 $2^3-1=7$.}
\solend
\fi

\item 设 $A=\{x\mid x^2-3x-4=0\}$，$B=\{x\mid ax-1=0\}$，若 $A\cap B=B$，则实数 $a$ 的值可以为\blankline[4.4em].

\ifanswer
\solbegin
\step{$A=\{-1,4\}$，因为 $A\cap B=B$，所以 $B\subseteq A$，}
\step{①$a=0$ 时，$B=\varnothing$，满足 $B\subseteq A$；}
\step{②$a\neq0$ 时，$B=\left\{\dfrac1a\right\}$，则 $\dfrac1a=-1$ 或 $\dfrac1a=4$，所以 $a=-1$ 或 $\dfrac14$；}
\step{所以 $a$ 的值可以为 $0$、$-1$、$\dfrac14$.}
\solend
\fi

\item 关于 $x$ 的不等式 $ax^2-(a+1)x+1\leqslant0(a<0)$ 的解集为\blankline[4.4em].

\ifanswer
\solbegin
\step{原不等式即 $(ax-1)(x-1)\leqslant0$，因为 $a<0$，所以 $\left(x-\dfrac1a\right)(x-1)\geqslant0$.}
\step{又 $\dfrac1a<0<1$，所以 $x\leqslant\dfrac1a$ 或 $x\geqslant1$，}
\step{故解集为 $\left(-\infty,\dfrac1a\right]\cup[1,+\infty)$.}
\solend
\fi

\item 若不等式 $(a-1)x^2+(2a-2)x-3<0$ 的解集是 $\R$，则 $a$ 的取值范围是\blankline[4.4em].

\ifanswer
\solbegin
\step{当 $a=1$ 时，不等式为 $-3<0$，对任意实数 $x$ 均成立.}
\step{当 $a\neq1$ 时，$\begin{cases}a-1<0\\ \Delta=(2a-2)^2+12(a-1)=4(a-1)(a+2)<0\end{cases}$，解得 $-2<a<1$.}
\step{合并得 $a\in(-2,1]$.}
\solend
\fi

\item 若关于 $x$ 的不等式 $x^2-(m+2)x+2m<0$ 的解集中恰有 $3$ 个整数，则实数 $m$ 的取值范围为\blankline[4.4em].

\ifanswer
\solbegin
\step{将不等式 $x^2-(m+2)x+2m<0$ 因式分解为 $(x-2)(x-m)<0$，}
\step{当 $m<2$ 时，解集为 $(m,2)$.}
\step{要使该区间恰有 $3$ 个整数，这 $3$ 个整数为 $-1,0,1$，故 $m$ 需满足 $-2\leqslant m<-1$.}
\step{当 $m>2$ 时，解集为 $(2,m)$. 要使该区间恰有 $3$ 个整数，}
\step{这 $3$ 个整数为 $3$，$4$，$5$，故 $m$ 需满足 $5<m\leqslant6$.}
\step{当 $m=2$ 时，不等式无解，不满足条件.}
\step{综上，实数 $m$ 的取值范围是 $[-2,-1)\cup(5,6]$.}
\solend
\fi

\item 已知集合 $A=\{x\mid x^2-3x+2\leqslant0\},B=\{y\mid y=x^2-2x+a,x\in\R\}$，若对任何一个 $x\in A$，都有 $x\in B$，则实数 $a$ 的取值范围是\blankline[4.4em].

\ifanswer
\solbegin
\step{$(x-1)(x-2)\leqslant0$，所以 $A=[1,2]$.}
\step{由 $y=(x-1)^2+a-1,x\in\R$，得 $B=[a-1,+\infty)$.}
\step{因此 $A\subseteq B\Leftrightarrow a-1\leqslant1\Leftrightarrow a\leqslant2$，所求范围为 $(-\infty,2]$.}
\solend
\fi

\end{enumerate}

\bigsec{二、选择题}

\begin{enumerate}
\setcounter{enumi}{12}

\item 已知 $a=\sqrt2$，集合 $A=\{x\mid x\leqslant2\}$，则下列表示正确的是\mbox{（\quad）}% 原卷括号内已印“D”，试卷版按惯例留空

\keepwithprev
A. $\{a\}\in A$\qquad B. $a\notin A$\qquad C. $a\subset A$\qquad D. $a\in A$

\ifanswer
\ans{$D$}
\fi

\item 下列四个命题中，真命题是\mbox{（\quad）}% 原卷括号内已印“B”，试卷版按惯例留空

\keepwithprev
A. 若 $a>b$，则 $ac^2>bc^2$\qquad B. 若 $a>b,c>d$，则 $a-d>b-c$

C. 若 $\lvert a\rvert>b$，则 $a^2>b^2$\qquad D. 若 $a>b$，则 $\dfrac1a<\dfrac1b$

\ifanswer
\ans{$B$}
\fi

\item “$x>1$”是“$\dfrac1x<1$”的\mbox{（\quad）}

\keepwithprev
A. 充分不必要条件\qquad B. 必要不充分条件

C. 充要条件\qquad D. 既不充分也不必要条件

\ifanswer
\solbegin
\step{$x>1\Rightarrow0<\dfrac1x<1$，充分性成立.}
\step{当 $x=-1$ 时，$\dfrac1x=-1<1$，但 $x>1$ 不成立，必要性不成立. 故选 $A$.}
\solend
\fi

\item 对于集合 $M=\{a\mid a=x^2-y^2,x\in\Z,y\in\Z\}$，给出如下三个结论：①如果 $P=\{b\mid b=2n+1,n\in\Z\}$，那么 $P\subseteq M$；②如果 $c=4n+2,n\in\Z$，那么 $c\notin M$；③如果 $a_1\in M,a_2\in M$，那么 $a_1a_2\in M$. 其中正确结论的个数是\mbox{（\quad）}

\keepwithprev
A. $1$\qquad B. $2$\qquad C. $3$\qquad D. $0$

\ifanswer
\solbegin
\step{对于①：因为 $2n+1=(n+1)^2-n^2,n\in\Z$，所以 $b=2n+1\in M$，故 $P\subseteq M$，}
\step{故①正确；}
\step{对于②：因为 $c=4n+2=2(2n+1),n\in\Z$，所以 $c$ 为偶数，且不能被 $4$ 整除，}
\step{若 $c\in M$，则存在 $x\in\Z,y\in\Z$，得 $c=x^2-y^2=(x+y)(x-y),x\in\Z,y\in\Z$，}
\step{因为 $x+y$ 和 $x-y$ 同奇或同偶，}
\step{若 $x+y$ 和 $x-y$ 同奇，则 $c=(x+y)(x-y)$ 为奇数，矛盾，不符合，}
\step{若 $x+y$ 和 $x-y$ 同偶，则 $c=(x+y)(x-y)$ 能被 $4$ 整除，矛盾，不符合，}
\step{所以 $c\notin M$，故②正确；}
\step{对于③：因为 $a_1\in M,a_2\in M$，}
\step{所以存在 $x_1,x_2\in\Z,y_1,y_2\in\Z$ 使得 $a_1=x_1^2-y_1^2,a_2=x_2^2-y_2^2$，}
\step{所以 $a_1a_2=(x_1^2-y_1^2)(x_2^2-y_2^2)=(x_1x_2)^2+(y_1y_2)^2-(x_1y_2)^2-(x_2y_1)^2$}
\step{$=(x_1x_2+y_1y_2)^2-(x_1y_2+x_2y_1)^2$，}
\step{因为 $x_1x_2+y_1y_2,x_1y_2+x_2y_1\in\Z$，所以 $a_1a_2\in M$，故③正确.}
\step{故选 $C$.}
\solend
\fi

\end{enumerate}

\bigsec{三、解答题}

\begin{enumerate}
\setcounter{enumi}{16}

\item （1）求不等式组 $\begin{cases}-2x^2+x+15\geqslant0\\ (x+1)(3x-1)>0\end{cases}$ 的解集；

（2）讨论解关于 $x$ 的不等式 $mx+4<m^2+2x$.

\ansspace*[12]

\ifanswer
\solbegin
\stepm{（1）}{第二个不等式可化为 $(2x+5)(x-3)\leqslant0$，解得 $-\dfrac52\leqslant x\leqslant3$.}% 原卷把“第一个”误作“第二个”，照录
\step{第二个不等式 $(x+1)(3x-1)>0$ 的解集为 $(-\infty,-1)\cup\left(\dfrac13,+\infty\right)$.}
\step{原不等式组的解集为 $\left[-\dfrac52,-1\right)\cup\left(\dfrac13,3\right]$.}
\stepm{（2）}{移项并因式分解，得 $(m-2)x<m^2-4\Leftrightarrow(m-2)x<(m-2)(m+2)$.}
\step{当 $m>2$ 时，$m-2>0$，解集为 $(-\infty,m+2)$.}
\step{当 $m<2$ 时，$m-2<0$，解集为 $(m+2,+\infty)$.}
\step{当 $m=2$ 时，不等式化为 $0<0$，不成立，解集为 $\varnothing$.}
\solend
\fi

\item （1）设 $x$ 是实数，比较 $x^2+9$ 与 $6x$ 的值的大小；

（2）已知 $a>0$ 且 $a\neq1$，比较 $\dfrac{a^2+1}{a^2-1}$ 与 $\dfrac{a+1}{a-1}$ 的大小.

\ansspace*[11]

\ifanswer
\solbegin
\stepm{（1）}{作差得 $(x^2+9)-6x=(x-3)^2\geqslant0$，因此 $x^2+9\geqslant6x$；}
\stepm{（2）}{$\dfrac{a^2+1}{a^2-1}-\dfrac{a+1}{a-1}=\dfrac{-2a}{a^2-1}$，}
\step{当 $a>1$ 时，$-2a<0$，$a^2-1>0$，则 $\dfrac{-2a}{a^2-1}<0$，即 $\dfrac{a^2+1}{a^2-1}<\dfrac{a+1}{a-1}$；}
\step{当 $0<a<1$ 时，$-2a<0$，$a^2-1<0$，则 $\dfrac{-2a}{a^2-1}>0$，即 $\dfrac{a^2+1}{a^2-1}>\dfrac{a+1}{a-1}$；}
\step{综上，$a>1$ 时，$\dfrac{a^2+1}{a^2-1}<\dfrac{a+1}{a-1}$；$0<a<1$ 时，$\dfrac{a^2+1}{a^2-1}>\dfrac{a+1}{a-1}$.}
\solend
\fi

\item 已知 $a\in\R$，集合 $A=\left\{x\left|\dfrac{2-x}{x+1}\leqslant0\right.\right\},B=\{x\mid x^2+(3-a)x-3a<0\}$.

（1）当 $a=\dfrac13$ 时，求 $A$ 和 $B$；

（2）“$x\in B$”是“$x\in A$”的充分非必要条件，求实数 $a$ 的取值范围.

\ansspace*[13]

\ifanswer
\solbegin
\stepm{（1）}{由 $\dfrac{2-x}{x+1}\leqslant0$，得 $\begin{cases}(2-x)(x+1)\leqslant0\\ x\neq-1\end{cases}$，解得 $x<-1$ 或 $x\geqslant2$，}
\step{则 $A=\{x\mid x<-1$ 或 $x\geqslant2\}$；$\cdots\cdots\cdots\cdots\cdots\cdots\cdots\cdots\cdots\cdots\cdots\cdots\cdots\cdots\cdots\cdots$（3 分）}
\step{$B=\{x\mid x^2+(3-a)x-3a<0\}=\{x\mid(x+3)(x-a)<0\}$，}
\step{当 $a=\dfrac13$ 时，$B=\left\{x\left|(x+2)\left(x-\dfrac13\right)<0\right.\right\}$，}
\step{解得 $B=\left\{x\left|-3<x<\dfrac13\right.\right\}\cdots\cdots\cdots\cdots\cdots\cdots$（6 分）}
\stepm{（2）}{由 $A\cup B=A$，得 $B\subseteq A$，}% 原卷从 A∪B=A 起论，与题干“充分非必要”差一层，照录
\step{①当 $B=\varnothing$ 时，得 $a=-3$，符合题意；$\cdots\cdots\cdots\cdots\cdots$（2 分）}
\step{②当 $B\neq\varnothing$ 时，若 $a>-3$，则 $B=\{x\mid-3<x<a\}$，}
\step{由 $B\subseteq A$，得 $a\leqslant-1$，此时 $-3<a\leqslant-1$；$\cdots\cdots\cdots\cdots\cdots$（4 分）}
\step{若 $a<-3$，则 $B=\{x\mid a<x<-3\}$，此时 $B\subseteq A$ 恒成立，}
\step{故 $a<-3$ 符合题意；$\cdots\cdots\cdots$（6 分）}
\step{综上所述，实数 $a$ 的取值范围为 $a\leqslant-1$. $\cdots\cdots\cdots\cdots\cdots$（8 分）}
\solend
\fi

\item 已知集合 $A=\{x\mid x<-2$ 或 $x\geqslant3\}$，$U=R$.

（1）求 $\overline A$；

（2）若 $B=\{x\mid mx-1>0\}$，当 $B\subseteq A$ 时，求实数 $m$ 的取值范围；

（3）若 $B=\{x\mid2m-6\leqslant x\leqslant m+4\}$，当 $B\cap\overline A=\varnothing$ 时，求实数 $m$ 的取值范围.

\ansspace*[14]

\ifanswer
\solbegin
\stepm{（1）}{因为集合 $A=\{x\mid x<-2$ 或 $x\geqslant3\}$，$U=R$，所以 $\overline A=\{x\mid-2\leqslant x<3\}$；}
\stepm{（2）}{当 $m=0$ 时，$B=\varnothing$，符合题意，}
\step{当 $m>0$ 时，$B=\left\{x\left|x>\dfrac1m\right.\right\}$，则 $\dfrac1m\geqslant3$，解得 $0<m\leqslant\dfrac13$，}
\step{当 $m<0$ 时，$B=\left\{x\left|x<\dfrac1m\right.\right\}$，则 $\dfrac1m\leqslant-2$，解得 $-\dfrac12\leqslant m<0$，}
\step{综上所述，实数 $m$ 的取值范围 $\left[-\dfrac12,\dfrac13\right]$；}
\stepm{（3）}{由题意得 $B=\{x\mid2m-6\leqslant x\leqslant m+4\}$，$\overline A=\{x\mid-2\leqslant x<3\}$，且 $B\cap\overline A=\varnothing$，}
\step{当 $B=\varnothing$ 时，$2m-6>m+4$，解得 $m>10$，}
\step{当 $B\neq\varnothing$ 时，则 $\begin{cases}2m-6\leqslant m+4\\ m+4<-2\end{cases}$ 或 $\begin{cases}2m-6\leqslant m+4\\ 2m-6\geqslant3\end{cases}$，}
\step{解得 $m<-6$ 或 $\dfrac92\leqslant m\leqslant10$，}
\step{综上所述，实数 $m$ 的取值范围为 $(-\infty,-6)\cup\left[\dfrac92,+\infty\right)$.}
\solend
\fi

\item 已知命题 $p$：关于 $x$ 的方程 $x^2+mx+1=0$ 有两个不同的正根；命题 $q$：关于 $x$ 的方程 $4x^2+4(m-2)x=-1$ 有实根.

（1）若命题 $q$ 为假命题，求整数 $m$ 的解集；

（2）若命题 $p$ 为真命题，且关于 $x$ 的方程 $x^2+mx+1=0$ 的两个实数根 $x_1$、$x_2$ 满足 $3x_1+2x_2=5$，求实数 $m$ 的取值集合；

（3）若命题 $p$ 和命题 $q$ 有且仅有一个成立，求实数 $m$ 的取值范围.

\ansspace*[16]

\ifanswer
\solbegin
\stepm{（1）}{命题 $q$：关于 $x$ 的方程 $4x^2+4(m-2)x=-1$ 有实根为假命题，}
\step{即 $4x^2+4(m-2)x+1=0$ 无实根，}
\step{则 $\Delta=16(m-2)^2-16<0$，则 $1<m<3$，则整数 $m$ 的解集为 $\{2\}$；}
\stepm{（2）}{若命题 $p$ 为真命题，且关于 $x$ 的方程 $x^2+mx+1=0$ 的两个实数根 $x_1$、$x_2$}
\step{满足 $3x_1+2x_2=5$，则方程 $x^2+mx+1=0$ 有两个不同的正根，}
\step{则 $\begin{cases}\Delta=m^2-4>0\\ x_1+x_2=-m>0\end{cases}$，则 $m<-2$，又 $\begin{cases}x_1+x_2=-m\\ x_1x_2=1\\ 3x_1+2x_2=5\end{cases}$，则 $\begin{cases}x_1=5+2m\\ x_2=-3m-5\end{cases}$，}
\step{则 $(5+2m)(-3m-5)=1$，得 $m=-2$ 或 $m=-\dfrac{13}6$，}
\step{又因为 $m<-2$，所以 $m=-\dfrac{13}6$，所以实数 $m$ 的取值集合为 $\left\{-\dfrac{13}6\right\}$；}
\stepm{（3）}{由（2）得，当 $p$ 为真时，则有 $m<-2$，}
\step{当 $q$ 为真时，$\Delta=16(m-2)^2-16\geqslant0$，解得 $m\leqslant1$ 或 $m\geqslant3$，}
\step{又因为命题 $p$ 和命题 $q$ 只有一个为真，所以 $\begin{cases}m<-2\\ 1<m<3\end{cases}$ 或 $\begin{cases}m\geqslant-2\\ m\leqslant1\ \text{或}\ m\geqslant3\end{cases}$，}
\step{解得 $-2\leqslant m\leqslant1$ 或 $m\geqslant3$，所以实数 $m$ 的取值范围为 $[-2,1]\cup[3,+\infty)$.}
\solend
\fi

\end{enumerate}

\end{document}
"
%
% 原始材料：微信公众号“魔都数学加油站”文章，连载“27学年高一95”；
% 原文标题“27学年高一95——2026-2027年上海市金汇中学高一上10月月考”；
% 原文链接：https://mp.weixin.qq.com/s/H3yH73FTin3V7rxWQ-9fjw。
% 正文为 8 张 PNG 页（`原图/高一95_金汇中学_10月月考/01..08.png`），均 4762×6735，按页序读取。
%
% 篇幅与题号：原卷自第 3 题起编号（第 1、2 题不在原文内）。填空题 3—12（10 题）、
% 选择题 13—16（4 题）、解答题 17—21（5 题），共 19 题。
%
% 讲评版：第 4、5 题只印【答案】（“$m>0$”“$x=0$ 或 $x=1$”），本稿只给 \ans{}；
% 第 13、14 题无【解析】、答案“D”“B”直接印在题干括号内，本稿试卷版按惯例留空、
% 答案版提 \ans{}；其余 15 题印【解析】，其中第 15、16 题的答案在解析末句“故选 $A$／$C$”。
% 第 19 题解析带原卷的分步分值（“（3 分）”“（6 分）”等），照录。全卷无插图。
%
% 副标题：原卷未印副标题。本稿按内容标作“集合与常用逻辑用语、不等式”。
%
% 原卷笔误/存疑（均照录未改，正文对应处另有 % 注）：
%   ① 第 17 题 (1) 解析首句作“**第二个**不等式可化为 $(2x+5)(x-3)\leqslant0$”，而
%      $-2x^2+x+15\geqslant0$ 才是**第一个**不等式；原卷误标，照录；
%   ② 第 19 题 (2) 题干问“$x\in B$ 是 $x\in A$ 的充分非必要条件”，解析却从“$A\cup B=A$，
%      得 $B\subseteq A$”起论（少了“真子集”一层）；原卷照录，结论 $a\leqslant-1$ 未改。
%
% 与原卷的排版差异：
%   ① 原卷每题下直接印【答案】/【解析】，本稿按全稿惯例将答案与解析置于答案版；
%   ② 原卷未印副标题，本稿补出；填空横线统一排 \blankline[4.4em]；
%   ③ 原卷数集记号排斜体 $R$、$N$、$Z$，本稿按全稿惯例改用 $\R$、$\N$、$\Z$；
%   ④ 第 17—21 题的子问用行内“（1）（2）（3）”段，解析中的子问标号用 \stepm；
%   ⑤ 解答题按惯例加 \ansspace；选择题括号排 \mbox{（\quad）}。
%
% 行数核对：逐题数 \step/\stepm 条数，合计 87 条（Q3：2；Q4：0；Q5：0；Q6：1；Q7：2；
% Q8：4；Q9：3；Q10：3；Q11：7；Q12：3；Q13：0；Q14：0；Q15：2；Q16：14；Q17：7；
% Q18：5；Q19：12；Q20：10；Q21：12）。第 4、5 题只印答案，第 13、14 题无解析，
% 均不计入 \step 账。原卷解析的换行按语意合并，未逐行强拆。
%
\documentclass[a4paper,11pt]{article}
\usepackage{common}

\begin{document}

\examtitlest[7]{2026--2027 年金汇中学高一上 10 月月考}{集合与常用逻辑用语、不等式}

\bigsec{一、填空题}

\begin{enumerate}
\setcounter{enumi}{2}

\item 命题：如果 $ab>c$，那么 $a>\dfrac cb$，是\blankline[4.4em]（选填“真”、“假”）命题.

\ifanswer
\solbegin
\step{取 $a=1,b=-2,c=-3$，则 $ab=-2>-3=c$，但 $a=1<\dfrac cb=\dfrac32$，}
\step{所以原命题是假命题.}
\solend
\fi

\item 若“$x=0$”是“$x<m$”的充分条件，则实数 $m$ 的取值范围是\blankline[4.4em].

\ifanswer
\ans{$m>0$}
\fi

\item “$x\neq0$ 且 $x\neq1$”的否定形式为\blankline[4.4em].

\ifanswer
\ans{$x=0$ 或 $x=1$}
\fi

\item 若集合 $M=\{(x,y)\mid x+y=3\}$，$N=\left\{(x,y)\left|y=\dfrac12x\right.\right\}$，则 $M\cap N=\blankline[4.4em]$.

\ifanswer
\solbegin
\step{$x+\dfrac12x=3\Rightarrow x=2,y=1$，所以 $M\cap N=\{(2,1)\}$.}
\solend
\fi

\item 集合 $\{x\in\N\mid0<\lvert x-1\rvert<3\}$ 的真子集个数为\blankline[4.4em].

\ifanswer
\solbegin
\step{因为 $\{x\in\N\mid0<\lvert x-1\rvert<3\}=\{x\in\N\mid-3<x-1<3,x\neq1\}$}
\step{$=\{x\in\N\mid-2<x<4,x\neq1\}=\{0,2,3\}$，所以真子集的个数为 $2^3-1=7$.}
\solend
\fi

\item 设 $A=\{x\mid x^2-3x-4=0\}$，$B=\{x\mid ax-1=0\}$，若 $A\cap B=B$，则实数 $a$ 的值可以为\blankline[4.4em].

\ifanswer
\solbegin
\step{$A=\{-1,4\}$，因为 $A\cap B=B$，所以 $B\subseteq A$，}
\step{①$a=0$ 时，$B=\varnothing$，满足 $B\subseteq A$；}
\step{②$a\neq0$ 时，$B=\left\{\dfrac1a\right\}$，则 $\dfrac1a=-1$ 或 $\dfrac1a=4$，所以 $a=-1$ 或 $\dfrac14$；}
\step{所以 $a$ 的值可以为 $0$、$-1$、$\dfrac14$.}
\solend
\fi

\item 关于 $x$ 的不等式 $ax^2-(a+1)x+1\leqslant0(a<0)$ 的解集为\blankline[4.4em].

\ifanswer
\solbegin
\step{原不等式即 $(ax-1)(x-1)\leqslant0$，因为 $a<0$，所以 $\left(x-\dfrac1a\right)(x-1)\geqslant0$.}
\step{又 $\dfrac1a<0<1$，所以 $x\leqslant\dfrac1a$ 或 $x\geqslant1$，}
\step{故解集为 $\left(-\infty,\dfrac1a\right]\cup[1,+\infty)$.}
\solend
\fi

\item 若不等式 $(a-1)x^2+(2a-2)x-3<0$ 的解集是 $\R$，则 $a$ 的取值范围是\blankline[4.4em].

\ifanswer
\solbegin
\step{当 $a=1$ 时，不等式为 $-3<0$，对任意实数 $x$ 均成立.}
\step{当 $a\neq1$ 时，$\begin{cases}a-1<0\\ \Delta=(2a-2)^2+12(a-1)=4(a-1)(a+2)<0\end{cases}$，解得 $-2<a<1$.}
\step{合并得 $a\in(-2,1]$.}
\solend
\fi

\item 若关于 $x$ 的不等式 $x^2-(m+2)x+2m<0$ 的解集中恰有 $3$ 个整数，则实数 $m$ 的取值范围为\blankline[4.4em].

\ifanswer
\solbegin
\step{将不等式 $x^2-(m+2)x+2m<0$ 因式分解为 $(x-2)(x-m)<0$，}
\step{当 $m<2$ 时，解集为 $(m,2)$.}
\step{要使该区间恰有 $3$ 个整数，这 $3$ 个整数为 $-1,0,1$，故 $m$ 需满足 $-2\leqslant m<-1$.}
\step{当 $m>2$ 时，解集为 $(2,m)$. 要使该区间恰有 $3$ 个整数，}
\step{这 $3$ 个整数为 $3$，$4$，$5$，故 $m$ 需满足 $5<m\leqslant6$.}
\step{当 $m=2$ 时，不等式无解，不满足条件.}
\step{综上，实数 $m$ 的取值范围是 $[-2,-1)\cup(5,6]$.}
\solend
\fi

\item 已知集合 $A=\{x\mid x^2-3x+2\leqslant0\},B=\{y\mid y=x^2-2x+a,x\in\R\}$，若对任何一个 $x\in A$，都有 $x\in B$，则实数 $a$ 的取值范围是\blankline[4.4em].

\ifanswer
\solbegin
\step{$(x-1)(x-2)\leqslant0$，所以 $A=[1,2]$.}
\step{由 $y=(x-1)^2+a-1,x\in\R$，得 $B=[a-1,+\infty)$.}
\step{因此 $A\subseteq B\Leftrightarrow a-1\leqslant1\Leftrightarrow a\leqslant2$，所求范围为 $(-\infty,2]$.}
\solend
\fi

\end{enumerate}

\bigsec{二、选择题}

\begin{enumerate}
\setcounter{enumi}{12}

\item 已知 $a=\sqrt2$，集合 $A=\{x\mid x\leqslant2\}$，则下列表示正确的是\mbox{（\quad）}% 原卷括号内已印“D”，试卷版按惯例留空

\keepwithprev
A. $\{a\}\in A$\qquad B. $a\notin A$\qquad C. $a\subset A$\qquad D. $a\in A$

\ifanswer
\ans{$D$}
\fi

\item 下列四个命题中，真命题是\mbox{（\quad）}% 原卷括号内已印“B”，试卷版按惯例留空

\keepwithprev
A. 若 $a>b$，则 $ac^2>bc^2$\qquad B. 若 $a>b,c>d$，则 $a-d>b-c$

C. 若 $\lvert a\rvert>b$，则 $a^2>b^2$\qquad D. 若 $a>b$，则 $\dfrac1a<\dfrac1b$

\ifanswer
\ans{$B$}
\fi

\item “$x>1$”是“$\dfrac1x<1$”的\mbox{（\quad）}

\keepwithprev
A. 充分不必要条件\qquad B. 必要不充分条件

C. 充要条件\qquad D. 既不充分也不必要条件

\ifanswer
\solbegin
\step{$x>1\Rightarrow0<\dfrac1x<1$，充分性成立.}
\step{当 $x=-1$ 时，$\dfrac1x=-1<1$，但 $x>1$ 不成立，必要性不成立. 故选 $A$.}
\solend
\fi

\item 对于集合 $M=\{a\mid a=x^2-y^2,x\in\Z,y\in\Z\}$，给出如下三个结论：①如果 $P=\{b\mid b=2n+1,n\in\Z\}$，那么 $P\subseteq M$；②如果 $c=4n+2,n\in\Z$，那么 $c\notin M$；③如果 $a_1\in M,a_2\in M$，那么 $a_1a_2\in M$. 其中正确结论的个数是\mbox{（\quad）}

\keepwithprev
A. $1$\qquad B. $2$\qquad C. $3$\qquad D. $0$

\ifanswer
\solbegin
\step{对于①：因为 $2n+1=(n+1)^2-n^2,n\in\Z$，所以 $b=2n+1\in M$，故 $P\subseteq M$，}
\step{故①正确；}
\step{对于②：因为 $c=4n+2=2(2n+1),n\in\Z$，所以 $c$ 为偶数，且不能被 $4$ 整除，}
\step{若 $c\in M$，则存在 $x\in\Z,y\in\Z$，得 $c=x^2-y^2=(x+y)(x-y),x\in\Z,y\in\Z$，}
\step{因为 $x+y$ 和 $x-y$ 同奇或同偶，}
\step{若 $x+y$ 和 $x-y$ 同奇，则 $c=(x+y)(x-y)$ 为奇数，矛盾，不符合，}
\step{若 $x+y$ 和 $x-y$ 同偶，则 $c=(x+y)(x-y)$ 能被 $4$ 整除，矛盾，不符合，}
\step{所以 $c\notin M$，故②正确；}
\step{对于③：因为 $a_1\in M,a_2\in M$，}
\step{所以存在 $x_1,x_2\in\Z,y_1,y_2\in\Z$ 使得 $a_1=x_1^2-y_1^2,a_2=x_2^2-y_2^2$，}
\step{所以 $a_1a_2=(x_1^2-y_1^2)(x_2^2-y_2^2)=(x_1x_2)^2+(y_1y_2)^2-(x_1y_2)^2-(x_2y_1)^2$}
\step{$=(x_1x_2+y_1y_2)^2-(x_1y_2+x_2y_1)^2$，}
\step{因为 $x_1x_2+y_1y_2,x_1y_2+x_2y_1\in\Z$，所以 $a_1a_2\in M$，故③正确.}
\step{故选 $C$.}
\solend
\fi

\end{enumerate}

\bigsec{三、解答题}

\begin{enumerate}
\setcounter{enumi}{16}

\item （1）求不等式组 $\begin{cases}-2x^2+x+15\geqslant0\\ (x+1)(3x-1)>0\end{cases}$ 的解集；

（2）讨论解关于 $x$ 的不等式 $mx+4<m^2+2x$.

\ansspace*[12]

\ifanswer
\solbegin
\stepm{（1）}{第二个不等式可化为 $(2x+5)(x-3)\leqslant0$，解得 $-\dfrac52\leqslant x\leqslant3$.}% 原卷把“第一个”误作“第二个”，照录
\step{第二个不等式 $(x+1)(3x-1)>0$ 的解集为 $(-\infty,-1)\cup\left(\dfrac13,+\infty\right)$.}
\step{原不等式组的解集为 $\left[-\dfrac52,-1\right)\cup\left(\dfrac13,3\right]$.}
\stepm{（2）}{移项并因式分解，得 $(m-2)x<m^2-4\Leftrightarrow(m-2)x<(m-2)(m+2)$.}
\step{当 $m>2$ 时，$m-2>0$，解集为 $(-\infty,m+2)$.}
\step{当 $m<2$ 时，$m-2<0$，解集为 $(m+2,+\infty)$.}
\step{当 $m=2$ 时，不等式化为 $0<0$，不成立，解集为 $\varnothing$.}
\solend
\fi

\item （1）设 $x$ 是实数，比较 $x^2+9$ 与 $6x$ 的值的大小；

（2）已知 $a>0$ 且 $a\neq1$，比较 $\dfrac{a^2+1}{a^2-1}$ 与 $\dfrac{a+1}{a-1}$ 的大小.

\ansspace*[11]

\ifanswer
\solbegin
\stepm{（1）}{作差得 $(x^2+9)-6x=(x-3)^2\geqslant0$，因此 $x^2+9\geqslant6x$；}
\stepm{（2）}{$\dfrac{a^2+1}{a^2-1}-\dfrac{a+1}{a-1}=\dfrac{-2a}{a^2-1}$，}
\step{当 $a>1$ 时，$-2a<0$，$a^2-1>0$，则 $\dfrac{-2a}{a^2-1}<0$，即 $\dfrac{a^2+1}{a^2-1}<\dfrac{a+1}{a-1}$；}
\step{当 $0<a<1$ 时，$-2a<0$，$a^2-1<0$，则 $\dfrac{-2a}{a^2-1}>0$，即 $\dfrac{a^2+1}{a^2-1}>\dfrac{a+1}{a-1}$；}
\step{综上，$a>1$ 时，$\dfrac{a^2+1}{a^2-1}<\dfrac{a+1}{a-1}$；$0<a<1$ 时，$\dfrac{a^2+1}{a^2-1}>\dfrac{a+1}{a-1}$.}
\solend
\fi

\item 已知 $a\in\R$，集合 $A=\left\{x\left|\dfrac{2-x}{x+1}\leqslant0\right.\right\},B=\{x\mid x^2+(3-a)x-3a<0\}$.

（1）当 $a=\dfrac13$ 时，求 $A$ 和 $B$；

（2）“$x\in B$”是“$x\in A$”的充分非必要条件，求实数 $a$ 的取值范围.

\ansspace*[13]

\ifanswer
\solbegin
\stepm{（1）}{由 $\dfrac{2-x}{x+1}\leqslant0$，得 $\begin{cases}(2-x)(x+1)\leqslant0\\ x\neq-1\end{cases}$，解得 $x<-1$ 或 $x\geqslant2$，}
\step{则 $A=\{x\mid x<-1$ 或 $x\geqslant2\}$；$\cdots\cdots\cdots\cdots\cdots\cdots\cdots\cdots\cdots\cdots\cdots\cdots\cdots\cdots\cdots\cdots$（3 分）}
\step{$B=\{x\mid x^2+(3-a)x-3a<0\}=\{x\mid(x+3)(x-a)<0\}$，}
\step{当 $a=\dfrac13$ 时，$B=\left\{x\left|(x+2)\left(x-\dfrac13\right)<0\right.\right\}$，}
\step{解得 $B=\left\{x\left|-3<x<\dfrac13\right.\right\}\cdots\cdots\cdots\cdots\cdots\cdots$（6 分）}
\stepm{（2）}{由 $A\cup B=A$，得 $B\subseteq A$，}% 原卷从 A∪B=A 起论，与题干“充分非必要”差一层，照录
\step{①当 $B=\varnothing$ 时，得 $a=-3$，符合题意；$\cdots\cdots\cdots\cdots\cdots$（2 分）}
\step{②当 $B\neq\varnothing$ 时，若 $a>-3$，则 $B=\{x\mid-3<x<a\}$，}
\step{由 $B\subseteq A$，得 $a\leqslant-1$，此时 $-3<a\leqslant-1$；$\cdots\cdots\cdots\cdots\cdots$（4 分）}
\step{若 $a<-3$，则 $B=\{x\mid a<x<-3\}$，此时 $B\subseteq A$ 恒成立，}
\step{故 $a<-3$ 符合题意；$\cdots\cdots\cdots$（6 分）}
\step{综上所述，实数 $a$ 的取值范围为 $a\leqslant-1$. $\cdots\cdots\cdots\cdots\cdots$（8 分）}
\solend
\fi

\item 已知集合 $A=\{x\mid x<-2$ 或 $x\geqslant3\}$，$U=R$.

（1）求 $\overline A$；

（2）若 $B=\{x\mid mx-1>0\}$，当 $B\subseteq A$ 时，求实数 $m$ 的取值范围；

（3）若 $B=\{x\mid2m-6\leqslant x\leqslant m+4\}$，当 $B\cap\overline A=\varnothing$ 时，求实数 $m$ 的取值范围.

\ansspace*[14]

\ifanswer
\solbegin
\stepm{（1）}{因为集合 $A=\{x\mid x<-2$ 或 $x\geqslant3\}$，$U=R$，所以 $\overline A=\{x\mid-2\leqslant x<3\}$；}
\stepm{（2）}{当 $m=0$ 时，$B=\varnothing$，符合题意，}
\step{当 $m>0$ 时，$B=\left\{x\left|x>\dfrac1m\right.\right\}$，则 $\dfrac1m\geqslant3$，解得 $0<m\leqslant\dfrac13$，}
\step{当 $m<0$ 时，$B=\left\{x\left|x<\dfrac1m\right.\right\}$，则 $\dfrac1m\leqslant-2$，解得 $-\dfrac12\leqslant m<0$，}
\step{综上所述，实数 $m$ 的取值范围 $\left[-\dfrac12,\dfrac13\right]$；}
\stepm{（3）}{由题意得 $B=\{x\mid2m-6\leqslant x\leqslant m+4\}$，$\overline A=\{x\mid-2\leqslant x<3\}$，且 $B\cap\overline A=\varnothing$，}
\step{当 $B=\varnothing$ 时，$2m-6>m+4$，解得 $m>10$，}
\step{当 $B\neq\varnothing$ 时，则 $\begin{cases}2m-6\leqslant m+4\\ m+4<-2\end{cases}$ 或 $\begin{cases}2m-6\leqslant m+4\\ 2m-6\geqslant3\end{cases}$，}
\step{解得 $m<-6$ 或 $\dfrac92\leqslant m\leqslant10$，}
\step{综上所述，实数 $m$ 的取值范围为 $(-\infty,-6)\cup\left[\dfrac92,+\infty\right)$.}
\solend
\fi

\item 已知命题 $p$：关于 $x$ 的方程 $x^2+mx+1=0$ 有两个不同的正根；命题 $q$：关于 $x$ 的方程 $4x^2+4(m-2)x=-1$ 有实根.

（1）若命题 $q$ 为假命题，求整数 $m$ 的解集；

（2）若命题 $p$ 为真命题，且关于 $x$ 的方程 $x^2+mx+1=0$ 的两个实数根 $x_1$、$x_2$ 满足 $3x_1+2x_2=5$，求实数 $m$ 的取值集合；

（3）若命题 $p$ 和命题 $q$ 有且仅有一个成立，求实数 $m$ 的取值范围.

\ansspace*[16]

\ifanswer
\solbegin
\stepm{（1）}{命题 $q$：关于 $x$ 的方程 $4x^2+4(m-2)x=-1$ 有实根为假命题，}
\step{即 $4x^2+4(m-2)x+1=0$ 无实根，}
\step{则 $\Delta=16(m-2)^2-16<0$，则 $1<m<3$，则整数 $m$ 的解集为 $\{2\}$；}
\stepm{（2）}{若命题 $p$ 为真命题，且关于 $x$ 的方程 $x^2+mx+1=0$ 的两个实数根 $x_1$、$x_2$}
\step{满足 $3x_1+2x_2=5$，则方程 $x^2+mx+1=0$ 有两个不同的正根，}
\step{则 $\begin{cases}\Delta=m^2-4>0\\ x_1+x_2=-m>0\end{cases}$，则 $m<-2$，又 $\begin{cases}x_1+x_2=-m\\ x_1x_2=1\\ 3x_1+2x_2=5\end{cases}$，则 $\begin{cases}x_1=5+2m\\ x_2=-3m-5\end{cases}$，}
\step{则 $(5+2m)(-3m-5)=1$，得 $m=-2$ 或 $m=-\dfrac{13}6$，}
\step{又因为 $m<-2$，所以 $m=-\dfrac{13}6$，所以实数 $m$ 的取值集合为 $\left\{-\dfrac{13}6\right\}$；}
\stepm{（3）}{由（2）得，当 $p$ 为真时，则有 $m<-2$，}
\step{当 $q$ 为真时，$\Delta=16(m-2)^2-16\geqslant0$，解得 $m\leqslant1$ 或 $m\geqslant3$，}
\step{又因为命题 $p$ 和命题 $q$ 只有一个为真，所以 $\begin{cases}m<-2\\ 1<m<3\end{cases}$ 或 $\begin{cases}m\geqslant-2\\ m\leqslant1\ \text{或}\ m\geqslant3\end{cases}$，}
\step{解得 $-2\leqslant m\leqslant1$ 或 $m\geqslant3$，所以实数 $m$ 的取值范围为 $[-2,1]\cup[3,+\infty)$.}
\solend
\fi

\end{enumerate}

\end{document}
"
% 紧凑版：xelatex "\def\iscompact{1}% 高一95_金汇中学_10月月考.tex
%   —— 高一上数学“2026-2027 年金汇中学高一上 10 月月考”（讲评版）
% 试卷版：xelatex 高一95_金汇中学_10月月考.tex
% 答案版：xelatex "\def\isanswer{1}% 高一95_金汇中学_10月月考.tex
%   —— 高一上数学“2026-2027 年金汇中学高一上 10 月月考”（讲评版）
% 试卷版：xelatex 高一95_金汇中学_10月月考.tex
% 答案版：xelatex "\def\isanswer{1}% 高一95_金汇中学_10月月考.tex
%   —— 高一上数学“2026-2027 年金汇中学高一上 10 月月考”（讲评版）
% 试卷版：xelatex 高一95_金汇中学_10月月考.tex
% 答案版：xelatex "\def\isanswer{1}\input{高一95_金汇中学_10月月考.tex}"
% 紧凑版：xelatex "\def\iscompact{1}\input{高一95_金汇中学_10月月考.tex}"
%
% 原始材料：微信公众号“魔都数学加油站”文章，连载“27学年高一95”；
% 原文标题“27学年高一95——2026-2027年上海市金汇中学高一上10月月考”；
% 原文链接：https://mp.weixin.qq.com/s/H3yH73FTin3V7rxWQ-9fjw。
% 正文为 8 张 PNG 页（`原图/高一95_金汇中学_10月月考/01..08.png`），均 4762×6735，按页序读取。
%
% 篇幅与题号：原卷自第 3 题起编号（第 1、2 题不在原文内）。填空题 3—12（10 题）、
% 选择题 13—16（4 题）、解答题 17—21（5 题），共 19 题。
%
% 讲评版：第 4、5 题只印【答案】（“$m>0$”“$x=0$ 或 $x=1$”），本稿只给 \ans{}；
% 第 13、14 题无【解析】、答案“D”“B”直接印在题干括号内，本稿试卷版按惯例留空、
% 答案版提 \ans{}；其余 15 题印【解析】，其中第 15、16 题的答案在解析末句“故选 $A$／$C$”。
% 第 19 题解析带原卷的分步分值（“（3 分）”“（6 分）”等），照录。全卷无插图。
%
% 副标题：原卷未印副标题。本稿按内容标作“集合与常用逻辑用语、不等式”。
%
% 原卷笔误/存疑（均照录未改，正文对应处另有 % 注）：
%   ① 第 17 题 (1) 解析首句作“**第二个**不等式可化为 $(2x+5)(x-3)\leqslant0$”，而
%      $-2x^2+x+15\geqslant0$ 才是**第一个**不等式；原卷误标，照录；
%   ② 第 19 题 (2) 题干问“$x\in B$ 是 $x\in A$ 的充分非必要条件”，解析却从“$A\cup B=A$，
%      得 $B\subseteq A$”起论（少了“真子集”一层）；原卷照录，结论 $a\leqslant-1$ 未改。
%
% 与原卷的排版差异：
%   ① 原卷每题下直接印【答案】/【解析】，本稿按全稿惯例将答案与解析置于答案版；
%   ② 原卷未印副标题，本稿补出；填空横线统一排 \blankline[4.4em]；
%   ③ 原卷数集记号排斜体 $R$、$N$、$Z$，本稿按全稿惯例改用 $\R$、$\N$、$\Z$；
%   ④ 第 17—21 题的子问用行内“（1）（2）（3）”段，解析中的子问标号用 \stepm；
%   ⑤ 解答题按惯例加 \ansspace；选择题括号排 \mbox{（\quad）}。
%
% 行数核对：逐题数 \step/\stepm 条数，合计 87 条（Q3：2；Q4：0；Q5：0；Q6：1；Q7：2；
% Q8：4；Q9：3；Q10：3；Q11：7；Q12：3；Q13：0；Q14：0；Q15：2；Q16：14；Q17：7；
% Q18：5；Q19：12；Q20：10；Q21：12）。第 4、5 题只印答案，第 13、14 题无解析，
% 均不计入 \step 账。原卷解析的换行按语意合并，未逐行强拆。
%
\documentclass[a4paper,11pt]{article}
\usepackage{common}

\begin{document}

\examtitlest[7]{2026--2027 年金汇中学高一上 10 月月考}{集合与常用逻辑用语、不等式}

\bigsec{一、填空题}

\begin{enumerate}
\setcounter{enumi}{2}

\item 命题：如果 $ab>c$，那么 $a>\dfrac cb$，是\blankline[4.4em]（选填“真”、“假”）命题.

\ifanswer
\solbegin
\step{取 $a=1,b=-2,c=-3$，则 $ab=-2>-3=c$，但 $a=1<\dfrac cb=\dfrac32$，}
\step{所以原命题是假命题.}
\solend
\fi

\item 若“$x=0$”是“$x<m$”的充分条件，则实数 $m$ 的取值范围是\blankline[4.4em].

\ifanswer
\ans{$m>0$}
\fi

\item “$x\neq0$ 且 $x\neq1$”的否定形式为\blankline[4.4em].

\ifanswer
\ans{$x=0$ 或 $x=1$}
\fi

\item 若集合 $M=\{(x,y)\mid x+y=3\}$，$N=\left\{(x,y)\left|y=\dfrac12x\right.\right\}$，则 $M\cap N=\blankline[4.4em]$.

\ifanswer
\solbegin
\step{$x+\dfrac12x=3\Rightarrow x=2,y=1$，所以 $M\cap N=\{(2,1)\}$.}
\solend
\fi

\item 集合 $\{x\in\N\mid0<\lvert x-1\rvert<3\}$ 的真子集个数为\blankline[4.4em].

\ifanswer
\solbegin
\step{因为 $\{x\in\N\mid0<\lvert x-1\rvert<3\}=\{x\in\N\mid-3<x-1<3,x\neq1\}$}
\step{$=\{x\in\N\mid-2<x<4,x\neq1\}=\{0,2,3\}$，所以真子集的个数为 $2^3-1=7$.}
\solend
\fi

\item 设 $A=\{x\mid x^2-3x-4=0\}$，$B=\{x\mid ax-1=0\}$，若 $A\cap B=B$，则实数 $a$ 的值可以为\blankline[4.4em].

\ifanswer
\solbegin
\step{$A=\{-1,4\}$，因为 $A\cap B=B$，所以 $B\subseteq A$，}
\step{①$a=0$ 时，$B=\varnothing$，满足 $B\subseteq A$；}
\step{②$a\neq0$ 时，$B=\left\{\dfrac1a\right\}$，则 $\dfrac1a=-1$ 或 $\dfrac1a=4$，所以 $a=-1$ 或 $\dfrac14$；}
\step{所以 $a$ 的值可以为 $0$、$-1$、$\dfrac14$.}
\solend
\fi

\item 关于 $x$ 的不等式 $ax^2-(a+1)x+1\leqslant0(a<0)$ 的解集为\blankline[4.4em].

\ifanswer
\solbegin
\step{原不等式即 $(ax-1)(x-1)\leqslant0$，因为 $a<0$，所以 $\left(x-\dfrac1a\right)(x-1)\geqslant0$.}
\step{又 $\dfrac1a<0<1$，所以 $x\leqslant\dfrac1a$ 或 $x\geqslant1$，}
\step{故解集为 $\left(-\infty,\dfrac1a\right]\cup[1,+\infty)$.}
\solend
\fi

\item 若不等式 $(a-1)x^2+(2a-2)x-3<0$ 的解集是 $\R$，则 $a$ 的取值范围是\blankline[4.4em].

\ifanswer
\solbegin
\step{当 $a=1$ 时，不等式为 $-3<0$，对任意实数 $x$ 均成立.}
\step{当 $a\neq1$ 时，$\begin{cases}a-1<0\\ \Delta=(2a-2)^2+12(a-1)=4(a-1)(a+2)<0\end{cases}$，解得 $-2<a<1$.}
\step{合并得 $a\in(-2,1]$.}
\solend
\fi

\item 若关于 $x$ 的不等式 $x^2-(m+2)x+2m<0$ 的解集中恰有 $3$ 个整数，则实数 $m$ 的取值范围为\blankline[4.4em].

\ifanswer
\solbegin
\step{将不等式 $x^2-(m+2)x+2m<0$ 因式分解为 $(x-2)(x-m)<0$，}
\step{当 $m<2$ 时，解集为 $(m,2)$.}
\step{要使该区间恰有 $3$ 个整数，这 $3$ 个整数为 $-1,0,1$，故 $m$ 需满足 $-2\leqslant m<-1$.}
\step{当 $m>2$ 时，解集为 $(2,m)$. 要使该区间恰有 $3$ 个整数，}
\step{这 $3$ 个整数为 $3$，$4$，$5$，故 $m$ 需满足 $5<m\leqslant6$.}
\step{当 $m=2$ 时，不等式无解，不满足条件.}
\step{综上，实数 $m$ 的取值范围是 $[-2,-1)\cup(5,6]$.}
\solend
\fi

\item 已知集合 $A=\{x\mid x^2-3x+2\leqslant0\},B=\{y\mid y=x^2-2x+a,x\in\R\}$，若对任何一个 $x\in A$，都有 $x\in B$，则实数 $a$ 的取值范围是\blankline[4.4em].

\ifanswer
\solbegin
\step{$(x-1)(x-2)\leqslant0$，所以 $A=[1,2]$.}
\step{由 $y=(x-1)^2+a-1,x\in\R$，得 $B=[a-1,+\infty)$.}
\step{因此 $A\subseteq B\Leftrightarrow a-1\leqslant1\Leftrightarrow a\leqslant2$，所求范围为 $(-\infty,2]$.}
\solend
\fi

\end{enumerate}

\bigsec{二、选择题}

\begin{enumerate}
\setcounter{enumi}{12}

\item 已知 $a=\sqrt2$，集合 $A=\{x\mid x\leqslant2\}$，则下列表示正确的是\mbox{（\quad）}% 原卷括号内已印“D”，试卷版按惯例留空

\keepwithprev
A. $\{a\}\in A$\qquad B. $a\notin A$\qquad C. $a\subset A$\qquad D. $a\in A$

\ifanswer
\ans{$D$}
\fi

\item 下列四个命题中，真命题是\mbox{（\quad）}% 原卷括号内已印“B”，试卷版按惯例留空

\keepwithprev
A. 若 $a>b$，则 $ac^2>bc^2$\qquad B. 若 $a>b,c>d$，则 $a-d>b-c$

C. 若 $\lvert a\rvert>b$，则 $a^2>b^2$\qquad D. 若 $a>b$，则 $\dfrac1a<\dfrac1b$

\ifanswer
\ans{$B$}
\fi

\item “$x>1$”是“$\dfrac1x<1$”的\mbox{（\quad）}

\keepwithprev
A. 充分不必要条件\qquad B. 必要不充分条件

C. 充要条件\qquad D. 既不充分也不必要条件

\ifanswer
\solbegin
\step{$x>1\Rightarrow0<\dfrac1x<1$，充分性成立.}
\step{当 $x=-1$ 时，$\dfrac1x=-1<1$，但 $x>1$ 不成立，必要性不成立. 故选 $A$.}
\solend
\fi

\item 对于集合 $M=\{a\mid a=x^2-y^2,x\in\Z,y\in\Z\}$，给出如下三个结论：①如果 $P=\{b\mid b=2n+1,n\in\Z\}$，那么 $P\subseteq M$；②如果 $c=4n+2,n\in\Z$，那么 $c\notin M$；③如果 $a_1\in M,a_2\in M$，那么 $a_1a_2\in M$. 其中正确结论的个数是\mbox{（\quad）}

\keepwithprev
A. $1$\qquad B. $2$\qquad C. $3$\qquad D. $0$

\ifanswer
\solbegin
\step{对于①：因为 $2n+1=(n+1)^2-n^2,n\in\Z$，所以 $b=2n+1\in M$，故 $P\subseteq M$，}
\step{故①正确；}
\step{对于②：因为 $c=4n+2=2(2n+1),n\in\Z$，所以 $c$ 为偶数，且不能被 $4$ 整除，}
\step{若 $c\in M$，则存在 $x\in\Z,y\in\Z$，得 $c=x^2-y^2=(x+y)(x-y),x\in\Z,y\in\Z$，}
\step{因为 $x+y$ 和 $x-y$ 同奇或同偶，}
\step{若 $x+y$ 和 $x-y$ 同奇，则 $c=(x+y)(x-y)$ 为奇数，矛盾，不符合，}
\step{若 $x+y$ 和 $x-y$ 同偶，则 $c=(x+y)(x-y)$ 能被 $4$ 整除，矛盾，不符合，}
\step{所以 $c\notin M$，故②正确；}
\step{对于③：因为 $a_1\in M,a_2\in M$，}
\step{所以存在 $x_1,x_2\in\Z,y_1,y_2\in\Z$ 使得 $a_1=x_1^2-y_1^2,a_2=x_2^2-y_2^2$，}
\step{所以 $a_1a_2=(x_1^2-y_1^2)(x_2^2-y_2^2)=(x_1x_2)^2+(y_1y_2)^2-(x_1y_2)^2-(x_2y_1)^2$}
\step{$=(x_1x_2+y_1y_2)^2-(x_1y_2+x_2y_1)^2$，}
\step{因为 $x_1x_2+y_1y_2,x_1y_2+x_2y_1\in\Z$，所以 $a_1a_2\in M$，故③正确.}
\step{故选 $C$.}
\solend
\fi

\end{enumerate}

\bigsec{三、解答题}

\begin{enumerate}
\setcounter{enumi}{16}

\item （1）求不等式组 $\begin{cases}-2x^2+x+15\geqslant0\\ (x+1)(3x-1)>0\end{cases}$ 的解集；

（2）讨论解关于 $x$ 的不等式 $mx+4<m^2+2x$.

\ansspace*[12]

\ifanswer
\solbegin
\stepm{（1）}{第二个不等式可化为 $(2x+5)(x-3)\leqslant0$，解得 $-\dfrac52\leqslant x\leqslant3$.}% 原卷把“第一个”误作“第二个”，照录
\step{第二个不等式 $(x+1)(3x-1)>0$ 的解集为 $(-\infty,-1)\cup\left(\dfrac13,+\infty\right)$.}
\step{原不等式组的解集为 $\left[-\dfrac52,-1\right)\cup\left(\dfrac13,3\right]$.}
\stepm{（2）}{移项并因式分解，得 $(m-2)x<m^2-4\Leftrightarrow(m-2)x<(m-2)(m+2)$.}
\step{当 $m>2$ 时，$m-2>0$，解集为 $(-\infty,m+2)$.}
\step{当 $m<2$ 时，$m-2<0$，解集为 $(m+2,+\infty)$.}
\step{当 $m=2$ 时，不等式化为 $0<0$，不成立，解集为 $\varnothing$.}
\solend
\fi

\item （1）设 $x$ 是实数，比较 $x^2+9$ 与 $6x$ 的值的大小；

（2）已知 $a>0$ 且 $a\neq1$，比较 $\dfrac{a^2+1}{a^2-1}$ 与 $\dfrac{a+1}{a-1}$ 的大小.

\ansspace*[11]

\ifanswer
\solbegin
\stepm{（1）}{作差得 $(x^2+9)-6x=(x-3)^2\geqslant0$，因此 $x^2+9\geqslant6x$；}
\stepm{（2）}{$\dfrac{a^2+1}{a^2-1}-\dfrac{a+1}{a-1}=\dfrac{-2a}{a^2-1}$，}
\step{当 $a>1$ 时，$-2a<0$，$a^2-1>0$，则 $\dfrac{-2a}{a^2-1}<0$，即 $\dfrac{a^2+1}{a^2-1}<\dfrac{a+1}{a-1}$；}
\step{当 $0<a<1$ 时，$-2a<0$，$a^2-1<0$，则 $\dfrac{-2a}{a^2-1}>0$，即 $\dfrac{a^2+1}{a^2-1}>\dfrac{a+1}{a-1}$；}
\step{综上，$a>1$ 时，$\dfrac{a^2+1}{a^2-1}<\dfrac{a+1}{a-1}$；$0<a<1$ 时，$\dfrac{a^2+1}{a^2-1}>\dfrac{a+1}{a-1}$.}
\solend
\fi

\item 已知 $a\in\R$，集合 $A=\left\{x\left|\dfrac{2-x}{x+1}\leqslant0\right.\right\},B=\{x\mid x^2+(3-a)x-3a<0\}$.

（1）当 $a=\dfrac13$ 时，求 $A$ 和 $B$；

（2）“$x\in B$”是“$x\in A$”的充分非必要条件，求实数 $a$ 的取值范围.

\ansspace*[13]

\ifanswer
\solbegin
\stepm{（1）}{由 $\dfrac{2-x}{x+1}\leqslant0$，得 $\begin{cases}(2-x)(x+1)\leqslant0\\ x\neq-1\end{cases}$，解得 $x<-1$ 或 $x\geqslant2$，}
\step{则 $A=\{x\mid x<-1$ 或 $x\geqslant2\}$；$\cdots\cdots\cdots\cdots\cdots\cdots\cdots\cdots\cdots\cdots\cdots\cdots\cdots\cdots\cdots\cdots$（3 分）}
\step{$B=\{x\mid x^2+(3-a)x-3a<0\}=\{x\mid(x+3)(x-a)<0\}$，}
\step{当 $a=\dfrac13$ 时，$B=\left\{x\left|(x+2)\left(x-\dfrac13\right)<0\right.\right\}$，}
\step{解得 $B=\left\{x\left|-3<x<\dfrac13\right.\right\}\cdots\cdots\cdots\cdots\cdots\cdots$（6 分）}
\stepm{（2）}{由 $A\cup B=A$，得 $B\subseteq A$，}% 原卷从 A∪B=A 起论，与题干“充分非必要”差一层，照录
\step{①当 $B=\varnothing$ 时，得 $a=-3$，符合题意；$\cdots\cdots\cdots\cdots\cdots$（2 分）}
\step{②当 $B\neq\varnothing$ 时，若 $a>-3$，则 $B=\{x\mid-3<x<a\}$，}
\step{由 $B\subseteq A$，得 $a\leqslant-1$，此时 $-3<a\leqslant-1$；$\cdots\cdots\cdots\cdots\cdots$（4 分）}
\step{若 $a<-3$，则 $B=\{x\mid a<x<-3\}$，此时 $B\subseteq A$ 恒成立，}
\step{故 $a<-3$ 符合题意；$\cdots\cdots\cdots$（6 分）}
\step{综上所述，实数 $a$ 的取值范围为 $a\leqslant-1$. $\cdots\cdots\cdots\cdots\cdots$（8 分）}
\solend
\fi

\item 已知集合 $A=\{x\mid x<-2$ 或 $x\geqslant3\}$，$U=R$.

（1）求 $\overline A$；

（2）若 $B=\{x\mid mx-1>0\}$，当 $B\subseteq A$ 时，求实数 $m$ 的取值范围；

（3）若 $B=\{x\mid2m-6\leqslant x\leqslant m+4\}$，当 $B\cap\overline A=\varnothing$ 时，求实数 $m$ 的取值范围.

\ansspace*[14]

\ifanswer
\solbegin
\stepm{（1）}{因为集合 $A=\{x\mid x<-2$ 或 $x\geqslant3\}$，$U=R$，所以 $\overline A=\{x\mid-2\leqslant x<3\}$；}
\stepm{（2）}{当 $m=0$ 时，$B=\varnothing$，符合题意，}
\step{当 $m>0$ 时，$B=\left\{x\left|x>\dfrac1m\right.\right\}$，则 $\dfrac1m\geqslant3$，解得 $0<m\leqslant\dfrac13$，}
\step{当 $m<0$ 时，$B=\left\{x\left|x<\dfrac1m\right.\right\}$，则 $\dfrac1m\leqslant-2$，解得 $-\dfrac12\leqslant m<0$，}
\step{综上所述，实数 $m$ 的取值范围 $\left[-\dfrac12,\dfrac13\right]$；}
\stepm{（3）}{由题意得 $B=\{x\mid2m-6\leqslant x\leqslant m+4\}$，$\overline A=\{x\mid-2\leqslant x<3\}$，且 $B\cap\overline A=\varnothing$，}
\step{当 $B=\varnothing$ 时，$2m-6>m+4$，解得 $m>10$，}
\step{当 $B\neq\varnothing$ 时，则 $\begin{cases}2m-6\leqslant m+4\\ m+4<-2\end{cases}$ 或 $\begin{cases}2m-6\leqslant m+4\\ 2m-6\geqslant3\end{cases}$，}
\step{解得 $m<-6$ 或 $\dfrac92\leqslant m\leqslant10$，}
\step{综上所述，实数 $m$ 的取值范围为 $(-\infty,-6)\cup\left[\dfrac92,+\infty\right)$.}
\solend
\fi

\item 已知命题 $p$：关于 $x$ 的方程 $x^2+mx+1=0$ 有两个不同的正根；命题 $q$：关于 $x$ 的方程 $4x^2+4(m-2)x=-1$ 有实根.

（1）若命题 $q$ 为假命题，求整数 $m$ 的解集；

（2）若命题 $p$ 为真命题，且关于 $x$ 的方程 $x^2+mx+1=0$ 的两个实数根 $x_1$、$x_2$ 满足 $3x_1+2x_2=5$，求实数 $m$ 的取值集合；

（3）若命题 $p$ 和命题 $q$ 有且仅有一个成立，求实数 $m$ 的取值范围.

\ansspace*[16]

\ifanswer
\solbegin
\stepm{（1）}{命题 $q$：关于 $x$ 的方程 $4x^2+4(m-2)x=-1$ 有实根为假命题，}
\step{即 $4x^2+4(m-2)x+1=0$ 无实根，}
\step{则 $\Delta=16(m-2)^2-16<0$，则 $1<m<3$，则整数 $m$ 的解集为 $\{2\}$；}
\stepm{（2）}{若命题 $p$ 为真命题，且关于 $x$ 的方程 $x^2+mx+1=0$ 的两个实数根 $x_1$、$x_2$}
\step{满足 $3x_1+2x_2=5$，则方程 $x^2+mx+1=0$ 有两个不同的正根，}
\step{则 $\begin{cases}\Delta=m^2-4>0\\ x_1+x_2=-m>0\end{cases}$，则 $m<-2$，又 $\begin{cases}x_1+x_2=-m\\ x_1x_2=1\\ 3x_1+2x_2=5\end{cases}$，则 $\begin{cases}x_1=5+2m\\ x_2=-3m-5\end{cases}$，}
\step{则 $(5+2m)(-3m-5)=1$，得 $m=-2$ 或 $m=-\dfrac{13}6$，}
\step{又因为 $m<-2$，所以 $m=-\dfrac{13}6$，所以实数 $m$ 的取值集合为 $\left\{-\dfrac{13}6\right\}$；}
\stepm{（3）}{由（2）得，当 $p$ 为真时，则有 $m<-2$，}
\step{当 $q$ 为真时，$\Delta=16(m-2)^2-16\geqslant0$，解得 $m\leqslant1$ 或 $m\geqslant3$，}
\step{又因为命题 $p$ 和命题 $q$ 只有一个为真，所以 $\begin{cases}m<-2\\ 1<m<3\end{cases}$ 或 $\begin{cases}m\geqslant-2\\ m\leqslant1\ \text{或}\ m\geqslant3\end{cases}$，}
\step{解得 $-2\leqslant m\leqslant1$ 或 $m\geqslant3$，所以实数 $m$ 的取值范围为 $[-2,1]\cup[3,+\infty)$.}
\solend
\fi

\end{enumerate}

\end{document}
"
% 紧凑版：xelatex "\def\iscompact{1}% 高一95_金汇中学_10月月考.tex
%   —— 高一上数学“2026-2027 年金汇中学高一上 10 月月考”（讲评版）
% 试卷版：xelatex 高一95_金汇中学_10月月考.tex
% 答案版：xelatex "\def\isanswer{1}\input{高一95_金汇中学_10月月考.tex}"
% 紧凑版：xelatex "\def\iscompact{1}\input{高一95_金汇中学_10月月考.tex}"
%
% 原始材料：微信公众号“魔都数学加油站”文章，连载“27学年高一95”；
% 原文标题“27学年高一95——2026-2027年上海市金汇中学高一上10月月考”；
% 原文链接：https://mp.weixin.qq.com/s/H3yH73FTin3V7rxWQ-9fjw。
% 正文为 8 张 PNG 页（`原图/高一95_金汇中学_10月月考/01..08.png`），均 4762×6735，按页序读取。
%
% 篇幅与题号：原卷自第 3 题起编号（第 1、2 题不在原文内）。填空题 3—12（10 题）、
% 选择题 13—16（4 题）、解答题 17—21（5 题），共 19 题。
%
% 讲评版：第 4、5 题只印【答案】（“$m>0$”“$x=0$ 或 $x=1$”），本稿只给 \ans{}；
% 第 13、14 题无【解析】、答案“D”“B”直接印在题干括号内，本稿试卷版按惯例留空、
% 答案版提 \ans{}；其余 15 题印【解析】，其中第 15、16 题的答案在解析末句“故选 $A$／$C$”。
% 第 19 题解析带原卷的分步分值（“（3 分）”“（6 分）”等），照录。全卷无插图。
%
% 副标题：原卷未印副标题。本稿按内容标作“集合与常用逻辑用语、不等式”。
%
% 原卷笔误/存疑（均照录未改，正文对应处另有 % 注）：
%   ① 第 17 题 (1) 解析首句作“**第二个**不等式可化为 $(2x+5)(x-3)\leqslant0$”，而
%      $-2x^2+x+15\geqslant0$ 才是**第一个**不等式；原卷误标，照录；
%   ② 第 19 题 (2) 题干问“$x\in B$ 是 $x\in A$ 的充分非必要条件”，解析却从“$A\cup B=A$，
%      得 $B\subseteq A$”起论（少了“真子集”一层）；原卷照录，结论 $a\leqslant-1$ 未改。
%
% 与原卷的排版差异：
%   ① 原卷每题下直接印【答案】/【解析】，本稿按全稿惯例将答案与解析置于答案版；
%   ② 原卷未印副标题，本稿补出；填空横线统一排 \blankline[4.4em]；
%   ③ 原卷数集记号排斜体 $R$、$N$、$Z$，本稿按全稿惯例改用 $\R$、$\N$、$\Z$；
%   ④ 第 17—21 题的子问用行内“（1）（2）（3）”段，解析中的子问标号用 \stepm；
%   ⑤ 解答题按惯例加 \ansspace；选择题括号排 \mbox{（\quad）}。
%
% 行数核对：逐题数 \step/\stepm 条数，合计 87 条（Q3：2；Q4：0；Q5：0；Q6：1；Q7：2；
% Q8：4；Q9：3；Q10：3；Q11：7；Q12：3；Q13：0；Q14：0；Q15：2；Q16：14；Q17：7；
% Q18：5；Q19：12；Q20：10；Q21：12）。第 4、5 题只印答案，第 13、14 题无解析，
% 均不计入 \step 账。原卷解析的换行按语意合并，未逐行强拆。
%
\documentclass[a4paper,11pt]{article}
\usepackage{common}

\begin{document}

\examtitlest[7]{2026--2027 年金汇中学高一上 10 月月考}{集合与常用逻辑用语、不等式}

\bigsec{一、填空题}

\begin{enumerate}
\setcounter{enumi}{2}

\item 命题：如果 $ab>c$，那么 $a>\dfrac cb$，是\blankline[4.4em]（选填“真”、“假”）命题.

\ifanswer
\solbegin
\step{取 $a=1,b=-2,c=-3$，则 $ab=-2>-3=c$，但 $a=1<\dfrac cb=\dfrac32$，}
\step{所以原命题是假命题.}
\solend
\fi

\item 若“$x=0$”是“$x<m$”的充分条件，则实数 $m$ 的取值范围是\blankline[4.4em].

\ifanswer
\ans{$m>0$}
\fi

\item “$x\neq0$ 且 $x\neq1$”的否定形式为\blankline[4.4em].

\ifanswer
\ans{$x=0$ 或 $x=1$}
\fi

\item 若集合 $M=\{(x,y)\mid x+y=3\}$，$N=\left\{(x,y)\left|y=\dfrac12x\right.\right\}$，则 $M\cap N=\blankline[4.4em]$.

\ifanswer
\solbegin
\step{$x+\dfrac12x=3\Rightarrow x=2,y=1$，所以 $M\cap N=\{(2,1)\}$.}
\solend
\fi

\item 集合 $\{x\in\N\mid0<\lvert x-1\rvert<3\}$ 的真子集个数为\blankline[4.4em].

\ifanswer
\solbegin
\step{因为 $\{x\in\N\mid0<\lvert x-1\rvert<3\}=\{x\in\N\mid-3<x-1<3,x\neq1\}$}
\step{$=\{x\in\N\mid-2<x<4,x\neq1\}=\{0,2,3\}$，所以真子集的个数为 $2^3-1=7$.}
\solend
\fi

\item 设 $A=\{x\mid x^2-3x-4=0\}$，$B=\{x\mid ax-1=0\}$，若 $A\cap B=B$，则实数 $a$ 的值可以为\blankline[4.4em].

\ifanswer
\solbegin
\step{$A=\{-1,4\}$，因为 $A\cap B=B$，所以 $B\subseteq A$，}
\step{①$a=0$ 时，$B=\varnothing$，满足 $B\subseteq A$；}
\step{②$a\neq0$ 时，$B=\left\{\dfrac1a\right\}$，则 $\dfrac1a=-1$ 或 $\dfrac1a=4$，所以 $a=-1$ 或 $\dfrac14$；}
\step{所以 $a$ 的值可以为 $0$、$-1$、$\dfrac14$.}
\solend
\fi

\item 关于 $x$ 的不等式 $ax^2-(a+1)x+1\leqslant0(a<0)$ 的解集为\blankline[4.4em].

\ifanswer
\solbegin
\step{原不等式即 $(ax-1)(x-1)\leqslant0$，因为 $a<0$，所以 $\left(x-\dfrac1a\right)(x-1)\geqslant0$.}
\step{又 $\dfrac1a<0<1$，所以 $x\leqslant\dfrac1a$ 或 $x\geqslant1$，}
\step{故解集为 $\left(-\infty,\dfrac1a\right]\cup[1,+\infty)$.}
\solend
\fi

\item 若不等式 $(a-1)x^2+(2a-2)x-3<0$ 的解集是 $\R$，则 $a$ 的取值范围是\blankline[4.4em].

\ifanswer
\solbegin
\step{当 $a=1$ 时，不等式为 $-3<0$，对任意实数 $x$ 均成立.}
\step{当 $a\neq1$ 时，$\begin{cases}a-1<0\\ \Delta=(2a-2)^2+12(a-1)=4(a-1)(a+2)<0\end{cases}$，解得 $-2<a<1$.}
\step{合并得 $a\in(-2,1]$.}
\solend
\fi

\item 若关于 $x$ 的不等式 $x^2-(m+2)x+2m<0$ 的解集中恰有 $3$ 个整数，则实数 $m$ 的取值范围为\blankline[4.4em].

\ifanswer
\solbegin
\step{将不等式 $x^2-(m+2)x+2m<0$ 因式分解为 $(x-2)(x-m)<0$，}
\step{当 $m<2$ 时，解集为 $(m,2)$.}
\step{要使该区间恰有 $3$ 个整数，这 $3$ 个整数为 $-1,0,1$，故 $m$ 需满足 $-2\leqslant m<-1$.}
\step{当 $m>2$ 时，解集为 $(2,m)$. 要使该区间恰有 $3$ 个整数，}
\step{这 $3$ 个整数为 $3$，$4$，$5$，故 $m$ 需满足 $5<m\leqslant6$.}
\step{当 $m=2$ 时，不等式无解，不满足条件.}
\step{综上，实数 $m$ 的取值范围是 $[-2,-1)\cup(5,6]$.}
\solend
\fi

\item 已知集合 $A=\{x\mid x^2-3x+2\leqslant0\},B=\{y\mid y=x^2-2x+a,x\in\R\}$，若对任何一个 $x\in A$，都有 $x\in B$，则实数 $a$ 的取值范围是\blankline[4.4em].

\ifanswer
\solbegin
\step{$(x-1)(x-2)\leqslant0$，所以 $A=[1,2]$.}
\step{由 $y=(x-1)^2+a-1,x\in\R$，得 $B=[a-1,+\infty)$.}
\step{因此 $A\subseteq B\Leftrightarrow a-1\leqslant1\Leftrightarrow a\leqslant2$，所求范围为 $(-\infty,2]$.}
\solend
\fi

\end{enumerate}

\bigsec{二、选择题}

\begin{enumerate}
\setcounter{enumi}{12}

\item 已知 $a=\sqrt2$，集合 $A=\{x\mid x\leqslant2\}$，则下列表示正确的是\mbox{（\quad）}% 原卷括号内已印“D”，试卷版按惯例留空

\keepwithprev
A. $\{a\}\in A$\qquad B. $a\notin A$\qquad C. $a\subset A$\qquad D. $a\in A$

\ifanswer
\ans{$D$}
\fi

\item 下列四个命题中，真命题是\mbox{（\quad）}% 原卷括号内已印“B”，试卷版按惯例留空

\keepwithprev
A. 若 $a>b$，则 $ac^2>bc^2$\qquad B. 若 $a>b,c>d$，则 $a-d>b-c$

C. 若 $\lvert a\rvert>b$，则 $a^2>b^2$\qquad D. 若 $a>b$，则 $\dfrac1a<\dfrac1b$

\ifanswer
\ans{$B$}
\fi

\item “$x>1$”是“$\dfrac1x<1$”的\mbox{（\quad）}

\keepwithprev
A. 充分不必要条件\qquad B. 必要不充分条件

C. 充要条件\qquad D. 既不充分也不必要条件

\ifanswer
\solbegin
\step{$x>1\Rightarrow0<\dfrac1x<1$，充分性成立.}
\step{当 $x=-1$ 时，$\dfrac1x=-1<1$，但 $x>1$ 不成立，必要性不成立. 故选 $A$.}
\solend
\fi

\item 对于集合 $M=\{a\mid a=x^2-y^2,x\in\Z,y\in\Z\}$，给出如下三个结论：①如果 $P=\{b\mid b=2n+1,n\in\Z\}$，那么 $P\subseteq M$；②如果 $c=4n+2,n\in\Z$，那么 $c\notin M$；③如果 $a_1\in M,a_2\in M$，那么 $a_1a_2\in M$. 其中正确结论的个数是\mbox{（\quad）}

\keepwithprev
A. $1$\qquad B. $2$\qquad C. $3$\qquad D. $0$

\ifanswer
\solbegin
\step{对于①：因为 $2n+1=(n+1)^2-n^2,n\in\Z$，所以 $b=2n+1\in M$，故 $P\subseteq M$，}
\step{故①正确；}
\step{对于②：因为 $c=4n+2=2(2n+1),n\in\Z$，所以 $c$ 为偶数，且不能被 $4$ 整除，}
\step{若 $c\in M$，则存在 $x\in\Z,y\in\Z$，得 $c=x^2-y^2=(x+y)(x-y),x\in\Z,y\in\Z$，}
\step{因为 $x+y$ 和 $x-y$ 同奇或同偶，}
\step{若 $x+y$ 和 $x-y$ 同奇，则 $c=(x+y)(x-y)$ 为奇数，矛盾，不符合，}
\step{若 $x+y$ 和 $x-y$ 同偶，则 $c=(x+y)(x-y)$ 能被 $4$ 整除，矛盾，不符合，}
\step{所以 $c\notin M$，故②正确；}
\step{对于③：因为 $a_1\in M,a_2\in M$，}
\step{所以存在 $x_1,x_2\in\Z,y_1,y_2\in\Z$ 使得 $a_1=x_1^2-y_1^2,a_2=x_2^2-y_2^2$，}
\step{所以 $a_1a_2=(x_1^2-y_1^2)(x_2^2-y_2^2)=(x_1x_2)^2+(y_1y_2)^2-(x_1y_2)^2-(x_2y_1)^2$}
\step{$=(x_1x_2+y_1y_2)^2-(x_1y_2+x_2y_1)^2$，}
\step{因为 $x_1x_2+y_1y_2,x_1y_2+x_2y_1\in\Z$，所以 $a_1a_2\in M$，故③正确.}
\step{故选 $C$.}
\solend
\fi

\end{enumerate}

\bigsec{三、解答题}

\begin{enumerate}
\setcounter{enumi}{16}

\item （1）求不等式组 $\begin{cases}-2x^2+x+15\geqslant0\\ (x+1)(3x-1)>0\end{cases}$ 的解集；

（2）讨论解关于 $x$ 的不等式 $mx+4<m^2+2x$.

\ansspace*[12]

\ifanswer
\solbegin
\stepm{（1）}{第二个不等式可化为 $(2x+5)(x-3)\leqslant0$，解得 $-\dfrac52\leqslant x\leqslant3$.}% 原卷把“第一个”误作“第二个”，照录
\step{第二个不等式 $(x+1)(3x-1)>0$ 的解集为 $(-\infty,-1)\cup\left(\dfrac13,+\infty\right)$.}
\step{原不等式组的解集为 $\left[-\dfrac52,-1\right)\cup\left(\dfrac13,3\right]$.}
\stepm{（2）}{移项并因式分解，得 $(m-2)x<m^2-4\Leftrightarrow(m-2)x<(m-2)(m+2)$.}
\step{当 $m>2$ 时，$m-2>0$，解集为 $(-\infty,m+2)$.}
\step{当 $m<2$ 时，$m-2<0$，解集为 $(m+2,+\infty)$.}
\step{当 $m=2$ 时，不等式化为 $0<0$，不成立，解集为 $\varnothing$.}
\solend
\fi

\item （1）设 $x$ 是实数，比较 $x^2+9$ 与 $6x$ 的值的大小；

（2）已知 $a>0$ 且 $a\neq1$，比较 $\dfrac{a^2+1}{a^2-1}$ 与 $\dfrac{a+1}{a-1}$ 的大小.

\ansspace*[11]

\ifanswer
\solbegin
\stepm{（1）}{作差得 $(x^2+9)-6x=(x-3)^2\geqslant0$，因此 $x^2+9\geqslant6x$；}
\stepm{（2）}{$\dfrac{a^2+1}{a^2-1}-\dfrac{a+1}{a-1}=\dfrac{-2a}{a^2-1}$，}
\step{当 $a>1$ 时，$-2a<0$，$a^2-1>0$，则 $\dfrac{-2a}{a^2-1}<0$，即 $\dfrac{a^2+1}{a^2-1}<\dfrac{a+1}{a-1}$；}
\step{当 $0<a<1$ 时，$-2a<0$，$a^2-1<0$，则 $\dfrac{-2a}{a^2-1}>0$，即 $\dfrac{a^2+1}{a^2-1}>\dfrac{a+1}{a-1}$；}
\step{综上，$a>1$ 时，$\dfrac{a^2+1}{a^2-1}<\dfrac{a+1}{a-1}$；$0<a<1$ 时，$\dfrac{a^2+1}{a^2-1}>\dfrac{a+1}{a-1}$.}
\solend
\fi

\item 已知 $a\in\R$，集合 $A=\left\{x\left|\dfrac{2-x}{x+1}\leqslant0\right.\right\},B=\{x\mid x^2+(3-a)x-3a<0\}$.

（1）当 $a=\dfrac13$ 时，求 $A$ 和 $B$；

（2）“$x\in B$”是“$x\in A$”的充分非必要条件，求实数 $a$ 的取值范围.

\ansspace*[13]

\ifanswer
\solbegin
\stepm{（1）}{由 $\dfrac{2-x}{x+1}\leqslant0$，得 $\begin{cases}(2-x)(x+1)\leqslant0\\ x\neq-1\end{cases}$，解得 $x<-1$ 或 $x\geqslant2$，}
\step{则 $A=\{x\mid x<-1$ 或 $x\geqslant2\}$；$\cdots\cdots\cdots\cdots\cdots\cdots\cdots\cdots\cdots\cdots\cdots\cdots\cdots\cdots\cdots\cdots$（3 分）}
\step{$B=\{x\mid x^2+(3-a)x-3a<0\}=\{x\mid(x+3)(x-a)<0\}$，}
\step{当 $a=\dfrac13$ 时，$B=\left\{x\left|(x+2)\left(x-\dfrac13\right)<0\right.\right\}$，}
\step{解得 $B=\left\{x\left|-3<x<\dfrac13\right.\right\}\cdots\cdots\cdots\cdots\cdots\cdots$（6 分）}
\stepm{（2）}{由 $A\cup B=A$，得 $B\subseteq A$，}% 原卷从 A∪B=A 起论，与题干“充分非必要”差一层，照录
\step{①当 $B=\varnothing$ 时，得 $a=-3$，符合题意；$\cdots\cdots\cdots\cdots\cdots$（2 分）}
\step{②当 $B\neq\varnothing$ 时，若 $a>-3$，则 $B=\{x\mid-3<x<a\}$，}
\step{由 $B\subseteq A$，得 $a\leqslant-1$，此时 $-3<a\leqslant-1$；$\cdots\cdots\cdots\cdots\cdots$（4 分）}
\step{若 $a<-3$，则 $B=\{x\mid a<x<-3\}$，此时 $B\subseteq A$ 恒成立，}
\step{故 $a<-3$ 符合题意；$\cdots\cdots\cdots$（6 分）}
\step{综上所述，实数 $a$ 的取值范围为 $a\leqslant-1$. $\cdots\cdots\cdots\cdots\cdots$（8 分）}
\solend
\fi

\item 已知集合 $A=\{x\mid x<-2$ 或 $x\geqslant3\}$，$U=R$.

（1）求 $\overline A$；

（2）若 $B=\{x\mid mx-1>0\}$，当 $B\subseteq A$ 时，求实数 $m$ 的取值范围；

（3）若 $B=\{x\mid2m-6\leqslant x\leqslant m+4\}$，当 $B\cap\overline A=\varnothing$ 时，求实数 $m$ 的取值范围.

\ansspace*[14]

\ifanswer
\solbegin
\stepm{（1）}{因为集合 $A=\{x\mid x<-2$ 或 $x\geqslant3\}$，$U=R$，所以 $\overline A=\{x\mid-2\leqslant x<3\}$；}
\stepm{（2）}{当 $m=0$ 时，$B=\varnothing$，符合题意，}
\step{当 $m>0$ 时，$B=\left\{x\left|x>\dfrac1m\right.\right\}$，则 $\dfrac1m\geqslant3$，解得 $0<m\leqslant\dfrac13$，}
\step{当 $m<0$ 时，$B=\left\{x\left|x<\dfrac1m\right.\right\}$，则 $\dfrac1m\leqslant-2$，解得 $-\dfrac12\leqslant m<0$，}
\step{综上所述，实数 $m$ 的取值范围 $\left[-\dfrac12,\dfrac13\right]$；}
\stepm{（3）}{由题意得 $B=\{x\mid2m-6\leqslant x\leqslant m+4\}$，$\overline A=\{x\mid-2\leqslant x<3\}$，且 $B\cap\overline A=\varnothing$，}
\step{当 $B=\varnothing$ 时，$2m-6>m+4$，解得 $m>10$，}
\step{当 $B\neq\varnothing$ 时，则 $\begin{cases}2m-6\leqslant m+4\\ m+4<-2\end{cases}$ 或 $\begin{cases}2m-6\leqslant m+4\\ 2m-6\geqslant3\end{cases}$，}
\step{解得 $m<-6$ 或 $\dfrac92\leqslant m\leqslant10$，}
\step{综上所述，实数 $m$ 的取值范围为 $(-\infty,-6)\cup\left[\dfrac92,+\infty\right)$.}
\solend
\fi

\item 已知命题 $p$：关于 $x$ 的方程 $x^2+mx+1=0$ 有两个不同的正根；命题 $q$：关于 $x$ 的方程 $4x^2+4(m-2)x=-1$ 有实根.

（1）若命题 $q$ 为假命题，求整数 $m$ 的解集；

（2）若命题 $p$ 为真命题，且关于 $x$ 的方程 $x^2+mx+1=0$ 的两个实数根 $x_1$、$x_2$ 满足 $3x_1+2x_2=5$，求实数 $m$ 的取值集合；

（3）若命题 $p$ 和命题 $q$ 有且仅有一个成立，求实数 $m$ 的取值范围.

\ansspace*[16]

\ifanswer
\solbegin
\stepm{（1）}{命题 $q$：关于 $x$ 的方程 $4x^2+4(m-2)x=-1$ 有实根为假命题，}
\step{即 $4x^2+4(m-2)x+1=0$ 无实根，}
\step{则 $\Delta=16(m-2)^2-16<0$，则 $1<m<3$，则整数 $m$ 的解集为 $\{2\}$；}
\stepm{（2）}{若命题 $p$ 为真命题，且关于 $x$ 的方程 $x^2+mx+1=0$ 的两个实数根 $x_1$、$x_2$}
\step{满足 $3x_1+2x_2=5$，则方程 $x^2+mx+1=0$ 有两个不同的正根，}
\step{则 $\begin{cases}\Delta=m^2-4>0\\ x_1+x_2=-m>0\end{cases}$，则 $m<-2$，又 $\begin{cases}x_1+x_2=-m\\ x_1x_2=1\\ 3x_1+2x_2=5\end{cases}$，则 $\begin{cases}x_1=5+2m\\ x_2=-3m-5\end{cases}$，}
\step{则 $(5+2m)(-3m-5)=1$，得 $m=-2$ 或 $m=-\dfrac{13}6$，}
\step{又因为 $m<-2$，所以 $m=-\dfrac{13}6$，所以实数 $m$ 的取值集合为 $\left\{-\dfrac{13}6\right\}$；}
\stepm{（3）}{由（2）得，当 $p$ 为真时，则有 $m<-2$，}
\step{当 $q$ 为真时，$\Delta=16(m-2)^2-16\geqslant0$，解得 $m\leqslant1$ 或 $m\geqslant3$，}
\step{又因为命题 $p$ 和命题 $q$ 只有一个为真，所以 $\begin{cases}m<-2\\ 1<m<3\end{cases}$ 或 $\begin{cases}m\geqslant-2\\ m\leqslant1\ \text{或}\ m\geqslant3\end{cases}$，}
\step{解得 $-2\leqslant m\leqslant1$ 或 $m\geqslant3$，所以实数 $m$ 的取值范围为 $[-2,1]\cup[3,+\infty)$.}
\solend
\fi

\end{enumerate}

\end{document}
"
%
% 原始材料：微信公众号“魔都数学加油站”文章，连载“27学年高一95”；
% 原文标题“27学年高一95——2026-2027年上海市金汇中学高一上10月月考”；
% 原文链接：https://mp.weixin.qq.com/s/H3yH73FTin3V7rxWQ-9fjw。
% 正文为 8 张 PNG 页（`原图/高一95_金汇中学_10月月考/01..08.png`），均 4762×6735，按页序读取。
%
% 篇幅与题号：原卷自第 3 题起编号（第 1、2 题不在原文内）。填空题 3—12（10 题）、
% 选择题 13—16（4 题）、解答题 17—21（5 题），共 19 题。
%
% 讲评版：第 4、5 题只印【答案】（“$m>0$”“$x=0$ 或 $x=1$”），本稿只给 \ans{}；
% 第 13、14 题无【解析】、答案“D”“B”直接印在题干括号内，本稿试卷版按惯例留空、
% 答案版提 \ans{}；其余 15 题印【解析】，其中第 15、16 题的答案在解析末句“故选 $A$／$C$”。
% 第 19 题解析带原卷的分步分值（“（3 分）”“（6 分）”等），照录。全卷无插图。
%
% 副标题：原卷未印副标题。本稿按内容标作“集合与常用逻辑用语、不等式”。
%
% 原卷笔误/存疑（均照录未改，正文对应处另有 % 注）：
%   ① 第 17 题 (1) 解析首句作“**第二个**不等式可化为 $(2x+5)(x-3)\leqslant0$”，而
%      $-2x^2+x+15\geqslant0$ 才是**第一个**不等式；原卷误标，照录；
%   ② 第 19 题 (2) 题干问“$x\in B$ 是 $x\in A$ 的充分非必要条件”，解析却从“$A\cup B=A$，
%      得 $B\subseteq A$”起论（少了“真子集”一层）；原卷照录，结论 $a\leqslant-1$ 未改。
%
% 与原卷的排版差异：
%   ① 原卷每题下直接印【答案】/【解析】，本稿按全稿惯例将答案与解析置于答案版；
%   ② 原卷未印副标题，本稿补出；填空横线统一排 \blankline[4.4em]；
%   ③ 原卷数集记号排斜体 $R$、$N$、$Z$，本稿按全稿惯例改用 $\R$、$\N$、$\Z$；
%   ④ 第 17—21 题的子问用行内“（1）（2）（3）”段，解析中的子问标号用 \stepm；
%   ⑤ 解答题按惯例加 \ansspace；选择题括号排 \mbox{（\quad）}。
%
% 行数核对：逐题数 \step/\stepm 条数，合计 87 条（Q3：2；Q4：0；Q5：0；Q6：1；Q7：2；
% Q8：4；Q9：3；Q10：3；Q11：7；Q12：3；Q13：0；Q14：0；Q15：2；Q16：14；Q17：7；
% Q18：5；Q19：12；Q20：10；Q21：12）。第 4、5 题只印答案，第 13、14 题无解析，
% 均不计入 \step 账。原卷解析的换行按语意合并，未逐行强拆。
%
\documentclass[a4paper,11pt]{article}
\usepackage{common}

\begin{document}

\examtitlest[7]{2026--2027 年金汇中学高一上 10 月月考}{集合与常用逻辑用语、不等式}

\bigsec{一、填空题}

\begin{enumerate}
\setcounter{enumi}{2}

\item 命题：如果 $ab>c$，那么 $a>\dfrac cb$，是\blankline[4.4em]（选填“真”、“假”）命题.

\ifanswer
\solbegin
\step{取 $a=1,b=-2,c=-3$，则 $ab=-2>-3=c$，但 $a=1<\dfrac cb=\dfrac32$，}
\step{所以原命题是假命题.}
\solend
\fi

\item 若“$x=0$”是“$x<m$”的充分条件，则实数 $m$ 的取值范围是\blankline[4.4em].

\ifanswer
\ans{$m>0$}
\fi

\item “$x\neq0$ 且 $x\neq1$”的否定形式为\blankline[4.4em].

\ifanswer
\ans{$x=0$ 或 $x=1$}
\fi

\item 若集合 $M=\{(x,y)\mid x+y=3\}$，$N=\left\{(x,y)\left|y=\dfrac12x\right.\right\}$，则 $M\cap N=\blankline[4.4em]$.

\ifanswer
\solbegin
\step{$x+\dfrac12x=3\Rightarrow x=2,y=1$，所以 $M\cap N=\{(2,1)\}$.}
\solend
\fi

\item 集合 $\{x\in\N\mid0<\lvert x-1\rvert<3\}$ 的真子集个数为\blankline[4.4em].

\ifanswer
\solbegin
\step{因为 $\{x\in\N\mid0<\lvert x-1\rvert<3\}=\{x\in\N\mid-3<x-1<3,x\neq1\}$}
\step{$=\{x\in\N\mid-2<x<4,x\neq1\}=\{0,2,3\}$，所以真子集的个数为 $2^3-1=7$.}
\solend
\fi

\item 设 $A=\{x\mid x^2-3x-4=0\}$，$B=\{x\mid ax-1=0\}$，若 $A\cap B=B$，则实数 $a$ 的值可以为\blankline[4.4em].

\ifanswer
\solbegin
\step{$A=\{-1,4\}$，因为 $A\cap B=B$，所以 $B\subseteq A$，}
\step{①$a=0$ 时，$B=\varnothing$，满足 $B\subseteq A$；}
\step{②$a\neq0$ 时，$B=\left\{\dfrac1a\right\}$，则 $\dfrac1a=-1$ 或 $\dfrac1a=4$，所以 $a=-1$ 或 $\dfrac14$；}
\step{所以 $a$ 的值可以为 $0$、$-1$、$\dfrac14$.}
\solend
\fi

\item 关于 $x$ 的不等式 $ax^2-(a+1)x+1\leqslant0(a<0)$ 的解集为\blankline[4.4em].

\ifanswer
\solbegin
\step{原不等式即 $(ax-1)(x-1)\leqslant0$，因为 $a<0$，所以 $\left(x-\dfrac1a\right)(x-1)\geqslant0$.}
\step{又 $\dfrac1a<0<1$，所以 $x\leqslant\dfrac1a$ 或 $x\geqslant1$，}
\step{故解集为 $\left(-\infty,\dfrac1a\right]\cup[1,+\infty)$.}
\solend
\fi

\item 若不等式 $(a-1)x^2+(2a-2)x-3<0$ 的解集是 $\R$，则 $a$ 的取值范围是\blankline[4.4em].

\ifanswer
\solbegin
\step{当 $a=1$ 时，不等式为 $-3<0$，对任意实数 $x$ 均成立.}
\step{当 $a\neq1$ 时，$\begin{cases}a-1<0\\ \Delta=(2a-2)^2+12(a-1)=4(a-1)(a+2)<0\end{cases}$，解得 $-2<a<1$.}
\step{合并得 $a\in(-2,1]$.}
\solend
\fi

\item 若关于 $x$ 的不等式 $x^2-(m+2)x+2m<0$ 的解集中恰有 $3$ 个整数，则实数 $m$ 的取值范围为\blankline[4.4em].

\ifanswer
\solbegin
\step{将不等式 $x^2-(m+2)x+2m<0$ 因式分解为 $(x-2)(x-m)<0$，}
\step{当 $m<2$ 时，解集为 $(m,2)$.}
\step{要使该区间恰有 $3$ 个整数，这 $3$ 个整数为 $-1,0,1$，故 $m$ 需满足 $-2\leqslant m<-1$.}
\step{当 $m>2$ 时，解集为 $(2,m)$. 要使该区间恰有 $3$ 个整数，}
\step{这 $3$ 个整数为 $3$，$4$，$5$，故 $m$ 需满足 $5<m\leqslant6$.}
\step{当 $m=2$ 时，不等式无解，不满足条件.}
\step{综上，实数 $m$ 的取值范围是 $[-2,-1)\cup(5,6]$.}
\solend
\fi

\item 已知集合 $A=\{x\mid x^2-3x+2\leqslant0\},B=\{y\mid y=x^2-2x+a,x\in\R\}$，若对任何一个 $x\in A$，都有 $x\in B$，则实数 $a$ 的取值范围是\blankline[4.4em].

\ifanswer
\solbegin
\step{$(x-1)(x-2)\leqslant0$，所以 $A=[1,2]$.}
\step{由 $y=(x-1)^2+a-1,x\in\R$，得 $B=[a-1,+\infty)$.}
\step{因此 $A\subseteq B\Leftrightarrow a-1\leqslant1\Leftrightarrow a\leqslant2$，所求范围为 $(-\infty,2]$.}
\solend
\fi

\end{enumerate}

\bigsec{二、选择题}

\begin{enumerate}
\setcounter{enumi}{12}

\item 已知 $a=\sqrt2$，集合 $A=\{x\mid x\leqslant2\}$，则下列表示正确的是\mbox{（\quad）}% 原卷括号内已印“D”，试卷版按惯例留空

\keepwithprev
A. $\{a\}\in A$\qquad B. $a\notin A$\qquad C. $a\subset A$\qquad D. $a\in A$

\ifanswer
\ans{$D$}
\fi

\item 下列四个命题中，真命题是\mbox{（\quad）}% 原卷括号内已印“B”，试卷版按惯例留空

\keepwithprev
A. 若 $a>b$，则 $ac^2>bc^2$\qquad B. 若 $a>b,c>d$，则 $a-d>b-c$

C. 若 $\lvert a\rvert>b$，则 $a^2>b^2$\qquad D. 若 $a>b$，则 $\dfrac1a<\dfrac1b$

\ifanswer
\ans{$B$}
\fi

\item “$x>1$”是“$\dfrac1x<1$”的\mbox{（\quad）}

\keepwithprev
A. 充分不必要条件\qquad B. 必要不充分条件

C. 充要条件\qquad D. 既不充分也不必要条件

\ifanswer
\solbegin
\step{$x>1\Rightarrow0<\dfrac1x<1$，充分性成立.}
\step{当 $x=-1$ 时，$\dfrac1x=-1<1$，但 $x>1$ 不成立，必要性不成立. 故选 $A$.}
\solend
\fi

\item 对于集合 $M=\{a\mid a=x^2-y^2,x\in\Z,y\in\Z\}$，给出如下三个结论：①如果 $P=\{b\mid b=2n+1,n\in\Z\}$，那么 $P\subseteq M$；②如果 $c=4n+2,n\in\Z$，那么 $c\notin M$；③如果 $a_1\in M,a_2\in M$，那么 $a_1a_2\in M$. 其中正确结论的个数是\mbox{（\quad）}

\keepwithprev
A. $1$\qquad B. $2$\qquad C. $3$\qquad D. $0$

\ifanswer
\solbegin
\step{对于①：因为 $2n+1=(n+1)^2-n^2,n\in\Z$，所以 $b=2n+1\in M$，故 $P\subseteq M$，}
\step{故①正确；}
\step{对于②：因为 $c=4n+2=2(2n+1),n\in\Z$，所以 $c$ 为偶数，且不能被 $4$ 整除，}
\step{若 $c\in M$，则存在 $x\in\Z,y\in\Z$，得 $c=x^2-y^2=(x+y)(x-y),x\in\Z,y\in\Z$，}
\step{因为 $x+y$ 和 $x-y$ 同奇或同偶，}
\step{若 $x+y$ 和 $x-y$ 同奇，则 $c=(x+y)(x-y)$ 为奇数，矛盾，不符合，}
\step{若 $x+y$ 和 $x-y$ 同偶，则 $c=(x+y)(x-y)$ 能被 $4$ 整除，矛盾，不符合，}
\step{所以 $c\notin M$，故②正确；}
\step{对于③：因为 $a_1\in M,a_2\in M$，}
\step{所以存在 $x_1,x_2\in\Z,y_1,y_2\in\Z$ 使得 $a_1=x_1^2-y_1^2,a_2=x_2^2-y_2^2$，}
\step{所以 $a_1a_2=(x_1^2-y_1^2)(x_2^2-y_2^2)=(x_1x_2)^2+(y_1y_2)^2-(x_1y_2)^2-(x_2y_1)^2$}
\step{$=(x_1x_2+y_1y_2)^2-(x_1y_2+x_2y_1)^2$，}
\step{因为 $x_1x_2+y_1y_2,x_1y_2+x_2y_1\in\Z$，所以 $a_1a_2\in M$，故③正确.}
\step{故选 $C$.}
\solend
\fi

\end{enumerate}

\bigsec{三、解答题}

\begin{enumerate}
\setcounter{enumi}{16}

\item （1）求不等式组 $\begin{cases}-2x^2+x+15\geqslant0\\ (x+1)(3x-1)>0\end{cases}$ 的解集；

（2）讨论解关于 $x$ 的不等式 $mx+4<m^2+2x$.

\ansspace*[12]

\ifanswer
\solbegin
\stepm{（1）}{第二个不等式可化为 $(2x+5)(x-3)\leqslant0$，解得 $-\dfrac52\leqslant x\leqslant3$.}% 原卷把“第一个”误作“第二个”，照录
\step{第二个不等式 $(x+1)(3x-1)>0$ 的解集为 $(-\infty,-1)\cup\left(\dfrac13,+\infty\right)$.}
\step{原不等式组的解集为 $\left[-\dfrac52,-1\right)\cup\left(\dfrac13,3\right]$.}
\stepm{（2）}{移项并因式分解，得 $(m-2)x<m^2-4\Leftrightarrow(m-2)x<(m-2)(m+2)$.}
\step{当 $m>2$ 时，$m-2>0$，解集为 $(-\infty,m+2)$.}
\step{当 $m<2$ 时，$m-2<0$，解集为 $(m+2,+\infty)$.}
\step{当 $m=2$ 时，不等式化为 $0<0$，不成立，解集为 $\varnothing$.}
\solend
\fi

\item （1）设 $x$ 是实数，比较 $x^2+9$ 与 $6x$ 的值的大小；

（2）已知 $a>0$ 且 $a\neq1$，比较 $\dfrac{a^2+1}{a^2-1}$ 与 $\dfrac{a+1}{a-1}$ 的大小.

\ansspace*[11]

\ifanswer
\solbegin
\stepm{（1）}{作差得 $(x^2+9)-6x=(x-3)^2\geqslant0$，因此 $x^2+9\geqslant6x$；}
\stepm{（2）}{$\dfrac{a^2+1}{a^2-1}-\dfrac{a+1}{a-1}=\dfrac{-2a}{a^2-1}$，}
\step{当 $a>1$ 时，$-2a<0$，$a^2-1>0$，则 $\dfrac{-2a}{a^2-1}<0$，即 $\dfrac{a^2+1}{a^2-1}<\dfrac{a+1}{a-1}$；}
\step{当 $0<a<1$ 时，$-2a<0$，$a^2-1<0$，则 $\dfrac{-2a}{a^2-1}>0$，即 $\dfrac{a^2+1}{a^2-1}>\dfrac{a+1}{a-1}$；}
\step{综上，$a>1$ 时，$\dfrac{a^2+1}{a^2-1}<\dfrac{a+1}{a-1}$；$0<a<1$ 时，$\dfrac{a^2+1}{a^2-1}>\dfrac{a+1}{a-1}$.}
\solend
\fi

\item 已知 $a\in\R$，集合 $A=\left\{x\left|\dfrac{2-x}{x+1}\leqslant0\right.\right\},B=\{x\mid x^2+(3-a)x-3a<0\}$.

（1）当 $a=\dfrac13$ 时，求 $A$ 和 $B$；

（2）“$x\in B$”是“$x\in A$”的充分非必要条件，求实数 $a$ 的取值范围.

\ansspace*[13]

\ifanswer
\solbegin
\stepm{（1）}{由 $\dfrac{2-x}{x+1}\leqslant0$，得 $\begin{cases}(2-x)(x+1)\leqslant0\\ x\neq-1\end{cases}$，解得 $x<-1$ 或 $x\geqslant2$，}
\step{则 $A=\{x\mid x<-1$ 或 $x\geqslant2\}$；$\cdots\cdots\cdots\cdots\cdots\cdots\cdots\cdots\cdots\cdots\cdots\cdots\cdots\cdots\cdots\cdots$（3 分）}
\step{$B=\{x\mid x^2+(3-a)x-3a<0\}=\{x\mid(x+3)(x-a)<0\}$，}
\step{当 $a=\dfrac13$ 时，$B=\left\{x\left|(x+2)\left(x-\dfrac13\right)<0\right.\right\}$，}
\step{解得 $B=\left\{x\left|-3<x<\dfrac13\right.\right\}\cdots\cdots\cdots\cdots\cdots\cdots$（6 分）}
\stepm{（2）}{由 $A\cup B=A$，得 $B\subseteq A$，}% 原卷从 A∪B=A 起论，与题干“充分非必要”差一层，照录
\step{①当 $B=\varnothing$ 时，得 $a=-3$，符合题意；$\cdots\cdots\cdots\cdots\cdots$（2 分）}
\step{②当 $B\neq\varnothing$ 时，若 $a>-3$，则 $B=\{x\mid-3<x<a\}$，}
\step{由 $B\subseteq A$，得 $a\leqslant-1$，此时 $-3<a\leqslant-1$；$\cdots\cdots\cdots\cdots\cdots$（4 分）}
\step{若 $a<-3$，则 $B=\{x\mid a<x<-3\}$，此时 $B\subseteq A$ 恒成立，}
\step{故 $a<-3$ 符合题意；$\cdots\cdots\cdots$（6 分）}
\step{综上所述，实数 $a$ 的取值范围为 $a\leqslant-1$. $\cdots\cdots\cdots\cdots\cdots$（8 分）}
\solend
\fi

\item 已知集合 $A=\{x\mid x<-2$ 或 $x\geqslant3\}$，$U=R$.

（1）求 $\overline A$；

（2）若 $B=\{x\mid mx-1>0\}$，当 $B\subseteq A$ 时，求实数 $m$ 的取值范围；

（3）若 $B=\{x\mid2m-6\leqslant x\leqslant m+4\}$，当 $B\cap\overline A=\varnothing$ 时，求实数 $m$ 的取值范围.

\ansspace*[14]

\ifanswer
\solbegin
\stepm{（1）}{因为集合 $A=\{x\mid x<-2$ 或 $x\geqslant3\}$，$U=R$，所以 $\overline A=\{x\mid-2\leqslant x<3\}$；}
\stepm{（2）}{当 $m=0$ 时，$B=\varnothing$，符合题意，}
\step{当 $m>0$ 时，$B=\left\{x\left|x>\dfrac1m\right.\right\}$，则 $\dfrac1m\geqslant3$，解得 $0<m\leqslant\dfrac13$，}
\step{当 $m<0$ 时，$B=\left\{x\left|x<\dfrac1m\right.\right\}$，则 $\dfrac1m\leqslant-2$，解得 $-\dfrac12\leqslant m<0$，}
\step{综上所述，实数 $m$ 的取值范围 $\left[-\dfrac12,\dfrac13\right]$；}
\stepm{（3）}{由题意得 $B=\{x\mid2m-6\leqslant x\leqslant m+4\}$，$\overline A=\{x\mid-2\leqslant x<3\}$，且 $B\cap\overline A=\varnothing$，}
\step{当 $B=\varnothing$ 时，$2m-6>m+4$，解得 $m>10$，}
\step{当 $B\neq\varnothing$ 时，则 $\begin{cases}2m-6\leqslant m+4\\ m+4<-2\end{cases}$ 或 $\begin{cases}2m-6\leqslant m+4\\ 2m-6\geqslant3\end{cases}$，}
\step{解得 $m<-6$ 或 $\dfrac92\leqslant m\leqslant10$，}
\step{综上所述，实数 $m$ 的取值范围为 $(-\infty,-6)\cup\left[\dfrac92,+\infty\right)$.}
\solend
\fi

\item 已知命题 $p$：关于 $x$ 的方程 $x^2+mx+1=0$ 有两个不同的正根；命题 $q$：关于 $x$ 的方程 $4x^2+4(m-2)x=-1$ 有实根.

（1）若命题 $q$ 为假命题，求整数 $m$ 的解集；

（2）若命题 $p$ 为真命题，且关于 $x$ 的方程 $x^2+mx+1=0$ 的两个实数根 $x_1$、$x_2$ 满足 $3x_1+2x_2=5$，求实数 $m$ 的取值集合；

（3）若命题 $p$ 和命题 $q$ 有且仅有一个成立，求实数 $m$ 的取值范围.

\ansspace*[16]

\ifanswer
\solbegin
\stepm{（1）}{命题 $q$：关于 $x$ 的方程 $4x^2+4(m-2)x=-1$ 有实根为假命题，}
\step{即 $4x^2+4(m-2)x+1=0$ 无实根，}
\step{则 $\Delta=16(m-2)^2-16<0$，则 $1<m<3$，则整数 $m$ 的解集为 $\{2\}$；}
\stepm{（2）}{若命题 $p$ 为真命题，且关于 $x$ 的方程 $x^2+mx+1=0$ 的两个实数根 $x_1$、$x_2$}
\step{满足 $3x_1+2x_2=5$，则方程 $x^2+mx+1=0$ 有两个不同的正根，}
\step{则 $\begin{cases}\Delta=m^2-4>0\\ x_1+x_2=-m>0\end{cases}$，则 $m<-2$，又 $\begin{cases}x_1+x_2=-m\\ x_1x_2=1\\ 3x_1+2x_2=5\end{cases}$，则 $\begin{cases}x_1=5+2m\\ x_2=-3m-5\end{cases}$，}
\step{则 $(5+2m)(-3m-5)=1$，得 $m=-2$ 或 $m=-\dfrac{13}6$，}
\step{又因为 $m<-2$，所以 $m=-\dfrac{13}6$，所以实数 $m$ 的取值集合为 $\left\{-\dfrac{13}6\right\}$；}
\stepm{（3）}{由（2）得，当 $p$ 为真时，则有 $m<-2$，}
\step{当 $q$ 为真时，$\Delta=16(m-2)^2-16\geqslant0$，解得 $m\leqslant1$ 或 $m\geqslant3$，}
\step{又因为命题 $p$ 和命题 $q$ 只有一个为真，所以 $\begin{cases}m<-2\\ 1<m<3\end{cases}$ 或 $\begin{cases}m\geqslant-2\\ m\leqslant1\ \text{或}\ m\geqslant3\end{cases}$，}
\step{解得 $-2\leqslant m\leqslant1$ 或 $m\geqslant3$，所以实数 $m$ 的取值范围为 $[-2,1]\cup[3,+\infty)$.}
\solend
\fi

\end{enumerate}

\end{document}
"
% 紧凑版：xelatex "\def\iscompact{1}% 高一95_金汇中学_10月月考.tex
%   —— 高一上数学“2026-2027 年金汇中学高一上 10 月月考”（讲评版）
% 试卷版：xelatex 高一95_金汇中学_10月月考.tex
% 答案版：xelatex "\def\isanswer{1}% 高一95_金汇中学_10月月考.tex
%   —— 高一上数学“2026-2027 年金汇中学高一上 10 月月考”（讲评版）
% 试卷版：xelatex 高一95_金汇中学_10月月考.tex
% 答案版：xelatex "\def\isanswer{1}\input{高一95_金汇中学_10月月考.tex}"
% 紧凑版：xelatex "\def\iscompact{1}\input{高一95_金汇中学_10月月考.tex}"
%
% 原始材料：微信公众号“魔都数学加油站”文章，连载“27学年高一95”；
% 原文标题“27学年高一95——2026-2027年上海市金汇中学高一上10月月考”；
% 原文链接：https://mp.weixin.qq.com/s/H3yH73FTin3V7rxWQ-9fjw。
% 正文为 8 张 PNG 页（`原图/高一95_金汇中学_10月月考/01..08.png`），均 4762×6735，按页序读取。
%
% 篇幅与题号：原卷自第 3 题起编号（第 1、2 题不在原文内）。填空题 3—12（10 题）、
% 选择题 13—16（4 题）、解答题 17—21（5 题），共 19 题。
%
% 讲评版：第 4、5 题只印【答案】（“$m>0$”“$x=0$ 或 $x=1$”），本稿只给 \ans{}；
% 第 13、14 题无【解析】、答案“D”“B”直接印在题干括号内，本稿试卷版按惯例留空、
% 答案版提 \ans{}；其余 15 题印【解析】，其中第 15、16 题的答案在解析末句“故选 $A$／$C$”。
% 第 19 题解析带原卷的分步分值（“（3 分）”“（6 分）”等），照录。全卷无插图。
%
% 副标题：原卷未印副标题。本稿按内容标作“集合与常用逻辑用语、不等式”。
%
% 原卷笔误/存疑（均照录未改，正文对应处另有 % 注）：
%   ① 第 17 题 (1) 解析首句作“**第二个**不等式可化为 $(2x+5)(x-3)\leqslant0$”，而
%      $-2x^2+x+15\geqslant0$ 才是**第一个**不等式；原卷误标，照录；
%   ② 第 19 题 (2) 题干问“$x\in B$ 是 $x\in A$ 的充分非必要条件”，解析却从“$A\cup B=A$，
%      得 $B\subseteq A$”起论（少了“真子集”一层）；原卷照录，结论 $a\leqslant-1$ 未改。
%
% 与原卷的排版差异：
%   ① 原卷每题下直接印【答案】/【解析】，本稿按全稿惯例将答案与解析置于答案版；
%   ② 原卷未印副标题，本稿补出；填空横线统一排 \blankline[4.4em]；
%   ③ 原卷数集记号排斜体 $R$、$N$、$Z$，本稿按全稿惯例改用 $\R$、$\N$、$\Z$；
%   ④ 第 17—21 题的子问用行内“（1）（2）（3）”段，解析中的子问标号用 \stepm；
%   ⑤ 解答题按惯例加 \ansspace；选择题括号排 \mbox{（\quad）}。
%
% 行数核对：逐题数 \step/\stepm 条数，合计 87 条（Q3：2；Q4：0；Q5：0；Q6：1；Q7：2；
% Q8：4；Q9：3；Q10：3；Q11：7；Q12：3；Q13：0；Q14：0；Q15：2；Q16：14；Q17：7；
% Q18：5；Q19：12；Q20：10；Q21：12）。第 4、5 题只印答案，第 13、14 题无解析，
% 均不计入 \step 账。原卷解析的换行按语意合并，未逐行强拆。
%
\documentclass[a4paper,11pt]{article}
\usepackage{common}

\begin{document}

\examtitlest[7]{2026--2027 年金汇中学高一上 10 月月考}{集合与常用逻辑用语、不等式}

\bigsec{一、填空题}

\begin{enumerate}
\setcounter{enumi}{2}

\item 命题：如果 $ab>c$，那么 $a>\dfrac cb$，是\blankline[4.4em]（选填“真”、“假”）命题.

\ifanswer
\solbegin
\step{取 $a=1,b=-2,c=-3$，则 $ab=-2>-3=c$，但 $a=1<\dfrac cb=\dfrac32$，}
\step{所以原命题是假命题.}
\solend
\fi

\item 若“$x=0$”是“$x<m$”的充分条件，则实数 $m$ 的取值范围是\blankline[4.4em].

\ifanswer
\ans{$m>0$}
\fi

\item “$x\neq0$ 且 $x\neq1$”的否定形式为\blankline[4.4em].

\ifanswer
\ans{$x=0$ 或 $x=1$}
\fi

\item 若集合 $M=\{(x,y)\mid x+y=3\}$，$N=\left\{(x,y)\left|y=\dfrac12x\right.\right\}$，则 $M\cap N=\blankline[4.4em]$.

\ifanswer
\solbegin
\step{$x+\dfrac12x=3\Rightarrow x=2,y=1$，所以 $M\cap N=\{(2,1)\}$.}
\solend
\fi

\item 集合 $\{x\in\N\mid0<\lvert x-1\rvert<3\}$ 的真子集个数为\blankline[4.4em].

\ifanswer
\solbegin
\step{因为 $\{x\in\N\mid0<\lvert x-1\rvert<3\}=\{x\in\N\mid-3<x-1<3,x\neq1\}$}
\step{$=\{x\in\N\mid-2<x<4,x\neq1\}=\{0,2,3\}$，所以真子集的个数为 $2^3-1=7$.}
\solend
\fi

\item 设 $A=\{x\mid x^2-3x-4=0\}$，$B=\{x\mid ax-1=0\}$，若 $A\cap B=B$，则实数 $a$ 的值可以为\blankline[4.4em].

\ifanswer
\solbegin
\step{$A=\{-1,4\}$，因为 $A\cap B=B$，所以 $B\subseteq A$，}
\step{①$a=0$ 时，$B=\varnothing$，满足 $B\subseteq A$；}
\step{②$a\neq0$ 时，$B=\left\{\dfrac1a\right\}$，则 $\dfrac1a=-1$ 或 $\dfrac1a=4$，所以 $a=-1$ 或 $\dfrac14$；}
\step{所以 $a$ 的值可以为 $0$、$-1$、$\dfrac14$.}
\solend
\fi

\item 关于 $x$ 的不等式 $ax^2-(a+1)x+1\leqslant0(a<0)$ 的解集为\blankline[4.4em].

\ifanswer
\solbegin
\step{原不等式即 $(ax-1)(x-1)\leqslant0$，因为 $a<0$，所以 $\left(x-\dfrac1a\right)(x-1)\geqslant0$.}
\step{又 $\dfrac1a<0<1$，所以 $x\leqslant\dfrac1a$ 或 $x\geqslant1$，}
\step{故解集为 $\left(-\infty,\dfrac1a\right]\cup[1,+\infty)$.}
\solend
\fi

\item 若不等式 $(a-1)x^2+(2a-2)x-3<0$ 的解集是 $\R$，则 $a$ 的取值范围是\blankline[4.4em].

\ifanswer
\solbegin
\step{当 $a=1$ 时，不等式为 $-3<0$，对任意实数 $x$ 均成立.}
\step{当 $a\neq1$ 时，$\begin{cases}a-1<0\\ \Delta=(2a-2)^2+12(a-1)=4(a-1)(a+2)<0\end{cases}$，解得 $-2<a<1$.}
\step{合并得 $a\in(-2,1]$.}
\solend
\fi

\item 若关于 $x$ 的不等式 $x^2-(m+2)x+2m<0$ 的解集中恰有 $3$ 个整数，则实数 $m$ 的取值范围为\blankline[4.4em].

\ifanswer
\solbegin
\step{将不等式 $x^2-(m+2)x+2m<0$ 因式分解为 $(x-2)(x-m)<0$，}
\step{当 $m<2$ 时，解集为 $(m,2)$.}
\step{要使该区间恰有 $3$ 个整数，这 $3$ 个整数为 $-1,0,1$，故 $m$ 需满足 $-2\leqslant m<-1$.}
\step{当 $m>2$ 时，解集为 $(2,m)$. 要使该区间恰有 $3$ 个整数，}
\step{这 $3$ 个整数为 $3$，$4$，$5$，故 $m$ 需满足 $5<m\leqslant6$.}
\step{当 $m=2$ 时，不等式无解，不满足条件.}
\step{综上，实数 $m$ 的取值范围是 $[-2,-1)\cup(5,6]$.}
\solend
\fi

\item 已知集合 $A=\{x\mid x^2-3x+2\leqslant0\},B=\{y\mid y=x^2-2x+a,x\in\R\}$，若对任何一个 $x\in A$，都有 $x\in B$，则实数 $a$ 的取值范围是\blankline[4.4em].

\ifanswer
\solbegin
\step{$(x-1)(x-2)\leqslant0$，所以 $A=[1,2]$.}
\step{由 $y=(x-1)^2+a-1,x\in\R$，得 $B=[a-1,+\infty)$.}
\step{因此 $A\subseteq B\Leftrightarrow a-1\leqslant1\Leftrightarrow a\leqslant2$，所求范围为 $(-\infty,2]$.}
\solend
\fi

\end{enumerate}

\bigsec{二、选择题}

\begin{enumerate}
\setcounter{enumi}{12}

\item 已知 $a=\sqrt2$，集合 $A=\{x\mid x\leqslant2\}$，则下列表示正确的是\mbox{（\quad）}% 原卷括号内已印“D”，试卷版按惯例留空

\keepwithprev
A. $\{a\}\in A$\qquad B. $a\notin A$\qquad C. $a\subset A$\qquad D. $a\in A$

\ifanswer
\ans{$D$}
\fi

\item 下列四个命题中，真命题是\mbox{（\quad）}% 原卷括号内已印“B”，试卷版按惯例留空

\keepwithprev
A. 若 $a>b$，则 $ac^2>bc^2$\qquad B. 若 $a>b,c>d$，则 $a-d>b-c$

C. 若 $\lvert a\rvert>b$，则 $a^2>b^2$\qquad D. 若 $a>b$，则 $\dfrac1a<\dfrac1b$

\ifanswer
\ans{$B$}
\fi

\item “$x>1$”是“$\dfrac1x<1$”的\mbox{（\quad）}

\keepwithprev
A. 充分不必要条件\qquad B. 必要不充分条件

C. 充要条件\qquad D. 既不充分也不必要条件

\ifanswer
\solbegin
\step{$x>1\Rightarrow0<\dfrac1x<1$，充分性成立.}
\step{当 $x=-1$ 时，$\dfrac1x=-1<1$，但 $x>1$ 不成立，必要性不成立. 故选 $A$.}
\solend
\fi

\item 对于集合 $M=\{a\mid a=x^2-y^2,x\in\Z,y\in\Z\}$，给出如下三个结论：①如果 $P=\{b\mid b=2n+1,n\in\Z\}$，那么 $P\subseteq M$；②如果 $c=4n+2,n\in\Z$，那么 $c\notin M$；③如果 $a_1\in M,a_2\in M$，那么 $a_1a_2\in M$. 其中正确结论的个数是\mbox{（\quad）}

\keepwithprev
A. $1$\qquad B. $2$\qquad C. $3$\qquad D. $0$

\ifanswer
\solbegin
\step{对于①：因为 $2n+1=(n+1)^2-n^2,n\in\Z$，所以 $b=2n+1\in M$，故 $P\subseteq M$，}
\step{故①正确；}
\step{对于②：因为 $c=4n+2=2(2n+1),n\in\Z$，所以 $c$ 为偶数，且不能被 $4$ 整除，}
\step{若 $c\in M$，则存在 $x\in\Z,y\in\Z$，得 $c=x^2-y^2=(x+y)(x-y),x\in\Z,y\in\Z$，}
\step{因为 $x+y$ 和 $x-y$ 同奇或同偶，}
\step{若 $x+y$ 和 $x-y$ 同奇，则 $c=(x+y)(x-y)$ 为奇数，矛盾，不符合，}
\step{若 $x+y$ 和 $x-y$ 同偶，则 $c=(x+y)(x-y)$ 能被 $4$ 整除，矛盾，不符合，}
\step{所以 $c\notin M$，故②正确；}
\step{对于③：因为 $a_1\in M,a_2\in M$，}
\step{所以存在 $x_1,x_2\in\Z,y_1,y_2\in\Z$ 使得 $a_1=x_1^2-y_1^2,a_2=x_2^2-y_2^2$，}
\step{所以 $a_1a_2=(x_1^2-y_1^2)(x_2^2-y_2^2)=(x_1x_2)^2+(y_1y_2)^2-(x_1y_2)^2-(x_2y_1)^2$}
\step{$=(x_1x_2+y_1y_2)^2-(x_1y_2+x_2y_1)^2$，}
\step{因为 $x_1x_2+y_1y_2,x_1y_2+x_2y_1\in\Z$，所以 $a_1a_2\in M$，故③正确.}
\step{故选 $C$.}
\solend
\fi

\end{enumerate}

\bigsec{三、解答题}

\begin{enumerate}
\setcounter{enumi}{16}

\item （1）求不等式组 $\begin{cases}-2x^2+x+15\geqslant0\\ (x+1)(3x-1)>0\end{cases}$ 的解集；

（2）讨论解关于 $x$ 的不等式 $mx+4<m^2+2x$.

\ansspace*[12]

\ifanswer
\solbegin
\stepm{（1）}{第二个不等式可化为 $(2x+5)(x-3)\leqslant0$，解得 $-\dfrac52\leqslant x\leqslant3$.}% 原卷把“第一个”误作“第二个”，照录
\step{第二个不等式 $(x+1)(3x-1)>0$ 的解集为 $(-\infty,-1)\cup\left(\dfrac13,+\infty\right)$.}
\step{原不等式组的解集为 $\left[-\dfrac52,-1\right)\cup\left(\dfrac13,3\right]$.}
\stepm{（2）}{移项并因式分解，得 $(m-2)x<m^2-4\Leftrightarrow(m-2)x<(m-2)(m+2)$.}
\step{当 $m>2$ 时，$m-2>0$，解集为 $(-\infty,m+2)$.}
\step{当 $m<2$ 时，$m-2<0$，解集为 $(m+2,+\infty)$.}
\step{当 $m=2$ 时，不等式化为 $0<0$，不成立，解集为 $\varnothing$.}
\solend
\fi

\item （1）设 $x$ 是实数，比较 $x^2+9$ 与 $6x$ 的值的大小；

（2）已知 $a>0$ 且 $a\neq1$，比较 $\dfrac{a^2+1}{a^2-1}$ 与 $\dfrac{a+1}{a-1}$ 的大小.

\ansspace*[11]

\ifanswer
\solbegin
\stepm{（1）}{作差得 $(x^2+9)-6x=(x-3)^2\geqslant0$，因此 $x^2+9\geqslant6x$；}
\stepm{（2）}{$\dfrac{a^2+1}{a^2-1}-\dfrac{a+1}{a-1}=\dfrac{-2a}{a^2-1}$，}
\step{当 $a>1$ 时，$-2a<0$，$a^2-1>0$，则 $\dfrac{-2a}{a^2-1}<0$，即 $\dfrac{a^2+1}{a^2-1}<\dfrac{a+1}{a-1}$；}
\step{当 $0<a<1$ 时，$-2a<0$，$a^2-1<0$，则 $\dfrac{-2a}{a^2-1}>0$，即 $\dfrac{a^2+1}{a^2-1}>\dfrac{a+1}{a-1}$；}
\step{综上，$a>1$ 时，$\dfrac{a^2+1}{a^2-1}<\dfrac{a+1}{a-1}$；$0<a<1$ 时，$\dfrac{a^2+1}{a^2-1}>\dfrac{a+1}{a-1}$.}
\solend
\fi

\item 已知 $a\in\R$，集合 $A=\left\{x\left|\dfrac{2-x}{x+1}\leqslant0\right.\right\},B=\{x\mid x^2+(3-a)x-3a<0\}$.

（1）当 $a=\dfrac13$ 时，求 $A$ 和 $B$；

（2）“$x\in B$”是“$x\in A$”的充分非必要条件，求实数 $a$ 的取值范围.

\ansspace*[13]

\ifanswer
\solbegin
\stepm{（1）}{由 $\dfrac{2-x}{x+1}\leqslant0$，得 $\begin{cases}(2-x)(x+1)\leqslant0\\ x\neq-1\end{cases}$，解得 $x<-1$ 或 $x\geqslant2$，}
\step{则 $A=\{x\mid x<-1$ 或 $x\geqslant2\}$；$\cdots\cdots\cdots\cdots\cdots\cdots\cdots\cdots\cdots\cdots\cdots\cdots\cdots\cdots\cdots\cdots$（3 分）}
\step{$B=\{x\mid x^2+(3-a)x-3a<0\}=\{x\mid(x+3)(x-a)<0\}$，}
\step{当 $a=\dfrac13$ 时，$B=\left\{x\left|(x+2)\left(x-\dfrac13\right)<0\right.\right\}$，}
\step{解得 $B=\left\{x\left|-3<x<\dfrac13\right.\right\}\cdots\cdots\cdots\cdots\cdots\cdots$（6 分）}
\stepm{（2）}{由 $A\cup B=A$，得 $B\subseteq A$，}% 原卷从 A∪B=A 起论，与题干“充分非必要”差一层，照录
\step{①当 $B=\varnothing$ 时，得 $a=-3$，符合题意；$\cdots\cdots\cdots\cdots\cdots$（2 分）}
\step{②当 $B\neq\varnothing$ 时，若 $a>-3$，则 $B=\{x\mid-3<x<a\}$，}
\step{由 $B\subseteq A$，得 $a\leqslant-1$，此时 $-3<a\leqslant-1$；$\cdots\cdots\cdots\cdots\cdots$（4 分）}
\step{若 $a<-3$，则 $B=\{x\mid a<x<-3\}$，此时 $B\subseteq A$ 恒成立，}
\step{故 $a<-3$ 符合题意；$\cdots\cdots\cdots$（6 分）}
\step{综上所述，实数 $a$ 的取值范围为 $a\leqslant-1$. $\cdots\cdots\cdots\cdots\cdots$（8 分）}
\solend
\fi

\item 已知集合 $A=\{x\mid x<-2$ 或 $x\geqslant3\}$，$U=R$.

（1）求 $\overline A$；

（2）若 $B=\{x\mid mx-1>0\}$，当 $B\subseteq A$ 时，求实数 $m$ 的取值范围；

（3）若 $B=\{x\mid2m-6\leqslant x\leqslant m+4\}$，当 $B\cap\overline A=\varnothing$ 时，求实数 $m$ 的取值范围.

\ansspace*[14]

\ifanswer
\solbegin
\stepm{（1）}{因为集合 $A=\{x\mid x<-2$ 或 $x\geqslant3\}$，$U=R$，所以 $\overline A=\{x\mid-2\leqslant x<3\}$；}
\stepm{（2）}{当 $m=0$ 时，$B=\varnothing$，符合题意，}
\step{当 $m>0$ 时，$B=\left\{x\left|x>\dfrac1m\right.\right\}$，则 $\dfrac1m\geqslant3$，解得 $0<m\leqslant\dfrac13$，}
\step{当 $m<0$ 时，$B=\left\{x\left|x<\dfrac1m\right.\right\}$，则 $\dfrac1m\leqslant-2$，解得 $-\dfrac12\leqslant m<0$，}
\step{综上所述，实数 $m$ 的取值范围 $\left[-\dfrac12,\dfrac13\right]$；}
\stepm{（3）}{由题意得 $B=\{x\mid2m-6\leqslant x\leqslant m+4\}$，$\overline A=\{x\mid-2\leqslant x<3\}$，且 $B\cap\overline A=\varnothing$，}
\step{当 $B=\varnothing$ 时，$2m-6>m+4$，解得 $m>10$，}
\step{当 $B\neq\varnothing$ 时，则 $\begin{cases}2m-6\leqslant m+4\\ m+4<-2\end{cases}$ 或 $\begin{cases}2m-6\leqslant m+4\\ 2m-6\geqslant3\end{cases}$，}
\step{解得 $m<-6$ 或 $\dfrac92\leqslant m\leqslant10$，}
\step{综上所述，实数 $m$ 的取值范围为 $(-\infty,-6)\cup\left[\dfrac92,+\infty\right)$.}
\solend
\fi

\item 已知命题 $p$：关于 $x$ 的方程 $x^2+mx+1=0$ 有两个不同的正根；命题 $q$：关于 $x$ 的方程 $4x^2+4(m-2)x=-1$ 有实根.

（1）若命题 $q$ 为假命题，求整数 $m$ 的解集；

（2）若命题 $p$ 为真命题，且关于 $x$ 的方程 $x^2+mx+1=0$ 的两个实数根 $x_1$、$x_2$ 满足 $3x_1+2x_2=5$，求实数 $m$ 的取值集合；

（3）若命题 $p$ 和命题 $q$ 有且仅有一个成立，求实数 $m$ 的取值范围.

\ansspace*[16]

\ifanswer
\solbegin
\stepm{（1）}{命题 $q$：关于 $x$ 的方程 $4x^2+4(m-2)x=-1$ 有实根为假命题，}
\step{即 $4x^2+4(m-2)x+1=0$ 无实根，}
\step{则 $\Delta=16(m-2)^2-16<0$，则 $1<m<3$，则整数 $m$ 的解集为 $\{2\}$；}
\stepm{（2）}{若命题 $p$ 为真命题，且关于 $x$ 的方程 $x^2+mx+1=0$ 的两个实数根 $x_1$、$x_2$}
\step{满足 $3x_1+2x_2=5$，则方程 $x^2+mx+1=0$ 有两个不同的正根，}
\step{则 $\begin{cases}\Delta=m^2-4>0\\ x_1+x_2=-m>0\end{cases}$，则 $m<-2$，又 $\begin{cases}x_1+x_2=-m\\ x_1x_2=1\\ 3x_1+2x_2=5\end{cases}$，则 $\begin{cases}x_1=5+2m\\ x_2=-3m-5\end{cases}$，}
\step{则 $(5+2m)(-3m-5)=1$，得 $m=-2$ 或 $m=-\dfrac{13}6$，}
\step{又因为 $m<-2$，所以 $m=-\dfrac{13}6$，所以实数 $m$ 的取值集合为 $\left\{-\dfrac{13}6\right\}$；}
\stepm{（3）}{由（2）得，当 $p$ 为真时，则有 $m<-2$，}
\step{当 $q$ 为真时，$\Delta=16(m-2)^2-16\geqslant0$，解得 $m\leqslant1$ 或 $m\geqslant3$，}
\step{又因为命题 $p$ 和命题 $q$ 只有一个为真，所以 $\begin{cases}m<-2\\ 1<m<3\end{cases}$ 或 $\begin{cases}m\geqslant-2\\ m\leqslant1\ \text{或}\ m\geqslant3\end{cases}$，}
\step{解得 $-2\leqslant m\leqslant1$ 或 $m\geqslant3$，所以实数 $m$ 的取值范围为 $[-2,1]\cup[3,+\infty)$.}
\solend
\fi

\end{enumerate}

\end{document}
"
% 紧凑版：xelatex "\def\iscompact{1}% 高一95_金汇中学_10月月考.tex
%   —— 高一上数学“2026-2027 年金汇中学高一上 10 月月考”（讲评版）
% 试卷版：xelatex 高一95_金汇中学_10月月考.tex
% 答案版：xelatex "\def\isanswer{1}\input{高一95_金汇中学_10月月考.tex}"
% 紧凑版：xelatex "\def\iscompact{1}\input{高一95_金汇中学_10月月考.tex}"
%
% 原始材料：微信公众号“魔都数学加油站”文章，连载“27学年高一95”；
% 原文标题“27学年高一95——2026-2027年上海市金汇中学高一上10月月考”；
% 原文链接：https://mp.weixin.qq.com/s/H3yH73FTin3V7rxWQ-9fjw。
% 正文为 8 张 PNG 页（`原图/高一95_金汇中学_10月月考/01..08.png`），均 4762×6735，按页序读取。
%
% 篇幅与题号：原卷自第 3 题起编号（第 1、2 题不在原文内）。填空题 3—12（10 题）、
% 选择题 13—16（4 题）、解答题 17—21（5 题），共 19 题。
%
% 讲评版：第 4、5 题只印【答案】（“$m>0$”“$x=0$ 或 $x=1$”），本稿只给 \ans{}；
% 第 13、14 题无【解析】、答案“D”“B”直接印在题干括号内，本稿试卷版按惯例留空、
% 答案版提 \ans{}；其余 15 题印【解析】，其中第 15、16 题的答案在解析末句“故选 $A$／$C$”。
% 第 19 题解析带原卷的分步分值（“（3 分）”“（6 分）”等），照录。全卷无插图。
%
% 副标题：原卷未印副标题。本稿按内容标作“集合与常用逻辑用语、不等式”。
%
% 原卷笔误/存疑（均照录未改，正文对应处另有 % 注）：
%   ① 第 17 题 (1) 解析首句作“**第二个**不等式可化为 $(2x+5)(x-3)\leqslant0$”，而
%      $-2x^2+x+15\geqslant0$ 才是**第一个**不等式；原卷误标，照录；
%   ② 第 19 题 (2) 题干问“$x\in B$ 是 $x\in A$ 的充分非必要条件”，解析却从“$A\cup B=A$，
%      得 $B\subseteq A$”起论（少了“真子集”一层）；原卷照录，结论 $a\leqslant-1$ 未改。
%
% 与原卷的排版差异：
%   ① 原卷每题下直接印【答案】/【解析】，本稿按全稿惯例将答案与解析置于答案版；
%   ② 原卷未印副标题，本稿补出；填空横线统一排 \blankline[4.4em]；
%   ③ 原卷数集记号排斜体 $R$、$N$、$Z$，本稿按全稿惯例改用 $\R$、$\N$、$\Z$；
%   ④ 第 17—21 题的子问用行内“（1）（2）（3）”段，解析中的子问标号用 \stepm；
%   ⑤ 解答题按惯例加 \ansspace；选择题括号排 \mbox{（\quad）}。
%
% 行数核对：逐题数 \step/\stepm 条数，合计 87 条（Q3：2；Q4：0；Q5：0；Q6：1；Q7：2；
% Q8：4；Q9：3；Q10：3；Q11：7；Q12：3；Q13：0；Q14：0；Q15：2；Q16：14；Q17：7；
% Q18：5；Q19：12；Q20：10；Q21：12）。第 4、5 题只印答案，第 13、14 题无解析，
% 均不计入 \step 账。原卷解析的换行按语意合并，未逐行强拆。
%
\documentclass[a4paper,11pt]{article}
\usepackage{common}

\begin{document}

\examtitlest[7]{2026--2027 年金汇中学高一上 10 月月考}{集合与常用逻辑用语、不等式}

\bigsec{一、填空题}

\begin{enumerate}
\setcounter{enumi}{2}

\item 命题：如果 $ab>c$，那么 $a>\dfrac cb$，是\blankline[4.4em]（选填“真”、“假”）命题.

\ifanswer
\solbegin
\step{取 $a=1,b=-2,c=-3$，则 $ab=-2>-3=c$，但 $a=1<\dfrac cb=\dfrac32$，}
\step{所以原命题是假命题.}
\solend
\fi

\item 若“$x=0$”是“$x<m$”的充分条件，则实数 $m$ 的取值范围是\blankline[4.4em].

\ifanswer
\ans{$m>0$}
\fi

\item “$x\neq0$ 且 $x\neq1$”的否定形式为\blankline[4.4em].

\ifanswer
\ans{$x=0$ 或 $x=1$}
\fi

\item 若集合 $M=\{(x,y)\mid x+y=3\}$，$N=\left\{(x,y)\left|y=\dfrac12x\right.\right\}$，则 $M\cap N=\blankline[4.4em]$.

\ifanswer
\solbegin
\step{$x+\dfrac12x=3\Rightarrow x=2,y=1$，所以 $M\cap N=\{(2,1)\}$.}
\solend
\fi

\item 集合 $\{x\in\N\mid0<\lvert x-1\rvert<3\}$ 的真子集个数为\blankline[4.4em].

\ifanswer
\solbegin
\step{因为 $\{x\in\N\mid0<\lvert x-1\rvert<3\}=\{x\in\N\mid-3<x-1<3,x\neq1\}$}
\step{$=\{x\in\N\mid-2<x<4,x\neq1\}=\{0,2,3\}$，所以真子集的个数为 $2^3-1=7$.}
\solend
\fi

\item 设 $A=\{x\mid x^2-3x-4=0\}$，$B=\{x\mid ax-1=0\}$，若 $A\cap B=B$，则实数 $a$ 的值可以为\blankline[4.4em].

\ifanswer
\solbegin
\step{$A=\{-1,4\}$，因为 $A\cap B=B$，所以 $B\subseteq A$，}
\step{①$a=0$ 时，$B=\varnothing$，满足 $B\subseteq A$；}
\step{②$a\neq0$ 时，$B=\left\{\dfrac1a\right\}$，则 $\dfrac1a=-1$ 或 $\dfrac1a=4$，所以 $a=-1$ 或 $\dfrac14$；}
\step{所以 $a$ 的值可以为 $0$、$-1$、$\dfrac14$.}
\solend
\fi

\item 关于 $x$ 的不等式 $ax^2-(a+1)x+1\leqslant0(a<0)$ 的解集为\blankline[4.4em].

\ifanswer
\solbegin
\step{原不等式即 $(ax-1)(x-1)\leqslant0$，因为 $a<0$，所以 $\left(x-\dfrac1a\right)(x-1)\geqslant0$.}
\step{又 $\dfrac1a<0<1$，所以 $x\leqslant\dfrac1a$ 或 $x\geqslant1$，}
\step{故解集为 $\left(-\infty,\dfrac1a\right]\cup[1,+\infty)$.}
\solend
\fi

\item 若不等式 $(a-1)x^2+(2a-2)x-3<0$ 的解集是 $\R$，则 $a$ 的取值范围是\blankline[4.4em].

\ifanswer
\solbegin
\step{当 $a=1$ 时，不等式为 $-3<0$，对任意实数 $x$ 均成立.}
\step{当 $a\neq1$ 时，$\begin{cases}a-1<0\\ \Delta=(2a-2)^2+12(a-1)=4(a-1)(a+2)<0\end{cases}$，解得 $-2<a<1$.}
\step{合并得 $a\in(-2,1]$.}
\solend
\fi

\item 若关于 $x$ 的不等式 $x^2-(m+2)x+2m<0$ 的解集中恰有 $3$ 个整数，则实数 $m$ 的取值范围为\blankline[4.4em].

\ifanswer
\solbegin
\step{将不等式 $x^2-(m+2)x+2m<0$ 因式分解为 $(x-2)(x-m)<0$，}
\step{当 $m<2$ 时，解集为 $(m,2)$.}
\step{要使该区间恰有 $3$ 个整数，这 $3$ 个整数为 $-1,0,1$，故 $m$ 需满足 $-2\leqslant m<-1$.}
\step{当 $m>2$ 时，解集为 $(2,m)$. 要使该区间恰有 $3$ 个整数，}
\step{这 $3$ 个整数为 $3$，$4$，$5$，故 $m$ 需满足 $5<m\leqslant6$.}
\step{当 $m=2$ 时，不等式无解，不满足条件.}
\step{综上，实数 $m$ 的取值范围是 $[-2,-1)\cup(5,6]$.}
\solend
\fi

\item 已知集合 $A=\{x\mid x^2-3x+2\leqslant0\},B=\{y\mid y=x^2-2x+a,x\in\R\}$，若对任何一个 $x\in A$，都有 $x\in B$，则实数 $a$ 的取值范围是\blankline[4.4em].

\ifanswer
\solbegin
\step{$(x-1)(x-2)\leqslant0$，所以 $A=[1,2]$.}
\step{由 $y=(x-1)^2+a-1,x\in\R$，得 $B=[a-1,+\infty)$.}
\step{因此 $A\subseteq B\Leftrightarrow a-1\leqslant1\Leftrightarrow a\leqslant2$，所求范围为 $(-\infty,2]$.}
\solend
\fi

\end{enumerate}

\bigsec{二、选择题}

\begin{enumerate}
\setcounter{enumi}{12}

\item 已知 $a=\sqrt2$，集合 $A=\{x\mid x\leqslant2\}$，则下列表示正确的是\mbox{（\quad）}% 原卷括号内已印“D”，试卷版按惯例留空

\keepwithprev
A. $\{a\}\in A$\qquad B. $a\notin A$\qquad C. $a\subset A$\qquad D. $a\in A$

\ifanswer
\ans{$D$}
\fi

\item 下列四个命题中，真命题是\mbox{（\quad）}% 原卷括号内已印“B”，试卷版按惯例留空

\keepwithprev
A. 若 $a>b$，则 $ac^2>bc^2$\qquad B. 若 $a>b,c>d$，则 $a-d>b-c$

C. 若 $\lvert a\rvert>b$，则 $a^2>b^2$\qquad D. 若 $a>b$，则 $\dfrac1a<\dfrac1b$

\ifanswer
\ans{$B$}
\fi

\item “$x>1$”是“$\dfrac1x<1$”的\mbox{（\quad）}

\keepwithprev
A. 充分不必要条件\qquad B. 必要不充分条件

C. 充要条件\qquad D. 既不充分也不必要条件

\ifanswer
\solbegin
\step{$x>1\Rightarrow0<\dfrac1x<1$，充分性成立.}
\step{当 $x=-1$ 时，$\dfrac1x=-1<1$，但 $x>1$ 不成立，必要性不成立. 故选 $A$.}
\solend
\fi

\item 对于集合 $M=\{a\mid a=x^2-y^2,x\in\Z,y\in\Z\}$，给出如下三个结论：①如果 $P=\{b\mid b=2n+1,n\in\Z\}$，那么 $P\subseteq M$；②如果 $c=4n+2,n\in\Z$，那么 $c\notin M$；③如果 $a_1\in M,a_2\in M$，那么 $a_1a_2\in M$. 其中正确结论的个数是\mbox{（\quad）}

\keepwithprev
A. $1$\qquad B. $2$\qquad C. $3$\qquad D. $0$

\ifanswer
\solbegin
\step{对于①：因为 $2n+1=(n+1)^2-n^2,n\in\Z$，所以 $b=2n+1\in M$，故 $P\subseteq M$，}
\step{故①正确；}
\step{对于②：因为 $c=4n+2=2(2n+1),n\in\Z$，所以 $c$ 为偶数，且不能被 $4$ 整除，}
\step{若 $c\in M$，则存在 $x\in\Z,y\in\Z$，得 $c=x^2-y^2=(x+y)(x-y),x\in\Z,y\in\Z$，}
\step{因为 $x+y$ 和 $x-y$ 同奇或同偶，}
\step{若 $x+y$ 和 $x-y$ 同奇，则 $c=(x+y)(x-y)$ 为奇数，矛盾，不符合，}
\step{若 $x+y$ 和 $x-y$ 同偶，则 $c=(x+y)(x-y)$ 能被 $4$ 整除，矛盾，不符合，}
\step{所以 $c\notin M$，故②正确；}
\step{对于③：因为 $a_1\in M,a_2\in M$，}
\step{所以存在 $x_1,x_2\in\Z,y_1,y_2\in\Z$ 使得 $a_1=x_1^2-y_1^2,a_2=x_2^2-y_2^2$，}
\step{所以 $a_1a_2=(x_1^2-y_1^2)(x_2^2-y_2^2)=(x_1x_2)^2+(y_1y_2)^2-(x_1y_2)^2-(x_2y_1)^2$}
\step{$=(x_1x_2+y_1y_2)^2-(x_1y_2+x_2y_1)^2$，}
\step{因为 $x_1x_2+y_1y_2,x_1y_2+x_2y_1\in\Z$，所以 $a_1a_2\in M$，故③正确.}
\step{故选 $C$.}
\solend
\fi

\end{enumerate}

\bigsec{三、解答题}

\begin{enumerate}
\setcounter{enumi}{16}

\item （1）求不等式组 $\begin{cases}-2x^2+x+15\geqslant0\\ (x+1)(3x-1)>0\end{cases}$ 的解集；

（2）讨论解关于 $x$ 的不等式 $mx+4<m^2+2x$.

\ansspace*[12]

\ifanswer
\solbegin
\stepm{（1）}{第二个不等式可化为 $(2x+5)(x-3)\leqslant0$，解得 $-\dfrac52\leqslant x\leqslant3$.}% 原卷把“第一个”误作“第二个”，照录
\step{第二个不等式 $(x+1)(3x-1)>0$ 的解集为 $(-\infty,-1)\cup\left(\dfrac13,+\infty\right)$.}
\step{原不等式组的解集为 $\left[-\dfrac52,-1\right)\cup\left(\dfrac13,3\right]$.}
\stepm{（2）}{移项并因式分解，得 $(m-2)x<m^2-4\Leftrightarrow(m-2)x<(m-2)(m+2)$.}
\step{当 $m>2$ 时，$m-2>0$，解集为 $(-\infty,m+2)$.}
\step{当 $m<2$ 时，$m-2<0$，解集为 $(m+2,+\infty)$.}
\step{当 $m=2$ 时，不等式化为 $0<0$，不成立，解集为 $\varnothing$.}
\solend
\fi

\item （1）设 $x$ 是实数，比较 $x^2+9$ 与 $6x$ 的值的大小；

（2）已知 $a>0$ 且 $a\neq1$，比较 $\dfrac{a^2+1}{a^2-1}$ 与 $\dfrac{a+1}{a-1}$ 的大小.

\ansspace*[11]

\ifanswer
\solbegin
\stepm{（1）}{作差得 $(x^2+9)-6x=(x-3)^2\geqslant0$，因此 $x^2+9\geqslant6x$；}
\stepm{（2）}{$\dfrac{a^2+1}{a^2-1}-\dfrac{a+1}{a-1}=\dfrac{-2a}{a^2-1}$，}
\step{当 $a>1$ 时，$-2a<0$，$a^2-1>0$，则 $\dfrac{-2a}{a^2-1}<0$，即 $\dfrac{a^2+1}{a^2-1}<\dfrac{a+1}{a-1}$；}
\step{当 $0<a<1$ 时，$-2a<0$，$a^2-1<0$，则 $\dfrac{-2a}{a^2-1}>0$，即 $\dfrac{a^2+1}{a^2-1}>\dfrac{a+1}{a-1}$；}
\step{综上，$a>1$ 时，$\dfrac{a^2+1}{a^2-1}<\dfrac{a+1}{a-1}$；$0<a<1$ 时，$\dfrac{a^2+1}{a^2-1}>\dfrac{a+1}{a-1}$.}
\solend
\fi

\item 已知 $a\in\R$，集合 $A=\left\{x\left|\dfrac{2-x}{x+1}\leqslant0\right.\right\},B=\{x\mid x^2+(3-a)x-3a<0\}$.

（1）当 $a=\dfrac13$ 时，求 $A$ 和 $B$；

（2）“$x\in B$”是“$x\in A$”的充分非必要条件，求实数 $a$ 的取值范围.

\ansspace*[13]

\ifanswer
\solbegin
\stepm{（1）}{由 $\dfrac{2-x}{x+1}\leqslant0$，得 $\begin{cases}(2-x)(x+1)\leqslant0\\ x\neq-1\end{cases}$，解得 $x<-1$ 或 $x\geqslant2$，}
\step{则 $A=\{x\mid x<-1$ 或 $x\geqslant2\}$；$\cdots\cdots\cdots\cdots\cdots\cdots\cdots\cdots\cdots\cdots\cdots\cdots\cdots\cdots\cdots\cdots$（3 分）}
\step{$B=\{x\mid x^2+(3-a)x-3a<0\}=\{x\mid(x+3)(x-a)<0\}$，}
\step{当 $a=\dfrac13$ 时，$B=\left\{x\left|(x+2)\left(x-\dfrac13\right)<0\right.\right\}$，}
\step{解得 $B=\left\{x\left|-3<x<\dfrac13\right.\right\}\cdots\cdots\cdots\cdots\cdots\cdots$（6 分）}
\stepm{（2）}{由 $A\cup B=A$，得 $B\subseteq A$，}% 原卷从 A∪B=A 起论，与题干“充分非必要”差一层，照录
\step{①当 $B=\varnothing$ 时，得 $a=-3$，符合题意；$\cdots\cdots\cdots\cdots\cdots$（2 分）}
\step{②当 $B\neq\varnothing$ 时，若 $a>-3$，则 $B=\{x\mid-3<x<a\}$，}
\step{由 $B\subseteq A$，得 $a\leqslant-1$，此时 $-3<a\leqslant-1$；$\cdots\cdots\cdots\cdots\cdots$（4 分）}
\step{若 $a<-3$，则 $B=\{x\mid a<x<-3\}$，此时 $B\subseteq A$ 恒成立，}
\step{故 $a<-3$ 符合题意；$\cdots\cdots\cdots$（6 分）}
\step{综上所述，实数 $a$ 的取值范围为 $a\leqslant-1$. $\cdots\cdots\cdots\cdots\cdots$（8 分）}
\solend
\fi

\item 已知集合 $A=\{x\mid x<-2$ 或 $x\geqslant3\}$，$U=R$.

（1）求 $\overline A$；

（2）若 $B=\{x\mid mx-1>0\}$，当 $B\subseteq A$ 时，求实数 $m$ 的取值范围；

（3）若 $B=\{x\mid2m-6\leqslant x\leqslant m+4\}$，当 $B\cap\overline A=\varnothing$ 时，求实数 $m$ 的取值范围.

\ansspace*[14]

\ifanswer
\solbegin
\stepm{（1）}{因为集合 $A=\{x\mid x<-2$ 或 $x\geqslant3\}$，$U=R$，所以 $\overline A=\{x\mid-2\leqslant x<3\}$；}
\stepm{（2）}{当 $m=0$ 时，$B=\varnothing$，符合题意，}
\step{当 $m>0$ 时，$B=\left\{x\left|x>\dfrac1m\right.\right\}$，则 $\dfrac1m\geqslant3$，解得 $0<m\leqslant\dfrac13$，}
\step{当 $m<0$ 时，$B=\left\{x\left|x<\dfrac1m\right.\right\}$，则 $\dfrac1m\leqslant-2$，解得 $-\dfrac12\leqslant m<0$，}
\step{综上所述，实数 $m$ 的取值范围 $\left[-\dfrac12,\dfrac13\right]$；}
\stepm{（3）}{由题意得 $B=\{x\mid2m-6\leqslant x\leqslant m+4\}$，$\overline A=\{x\mid-2\leqslant x<3\}$，且 $B\cap\overline A=\varnothing$，}
\step{当 $B=\varnothing$ 时，$2m-6>m+4$，解得 $m>10$，}
\step{当 $B\neq\varnothing$ 时，则 $\begin{cases}2m-6\leqslant m+4\\ m+4<-2\end{cases}$ 或 $\begin{cases}2m-6\leqslant m+4\\ 2m-6\geqslant3\end{cases}$，}
\step{解得 $m<-6$ 或 $\dfrac92\leqslant m\leqslant10$，}
\step{综上所述，实数 $m$ 的取值范围为 $(-\infty,-6)\cup\left[\dfrac92,+\infty\right)$.}
\solend
\fi

\item 已知命题 $p$：关于 $x$ 的方程 $x^2+mx+1=0$ 有两个不同的正根；命题 $q$：关于 $x$ 的方程 $4x^2+4(m-2)x=-1$ 有实根.

（1）若命题 $q$ 为假命题，求整数 $m$ 的解集；

（2）若命题 $p$ 为真命题，且关于 $x$ 的方程 $x^2+mx+1=0$ 的两个实数根 $x_1$、$x_2$ 满足 $3x_1+2x_2=5$，求实数 $m$ 的取值集合；

（3）若命题 $p$ 和命题 $q$ 有且仅有一个成立，求实数 $m$ 的取值范围.

\ansspace*[16]

\ifanswer
\solbegin
\stepm{（1）}{命题 $q$：关于 $x$ 的方程 $4x^2+4(m-2)x=-1$ 有实根为假命题，}
\step{即 $4x^2+4(m-2)x+1=0$ 无实根，}
\step{则 $\Delta=16(m-2)^2-16<0$，则 $1<m<3$，则整数 $m$ 的解集为 $\{2\}$；}
\stepm{（2）}{若命题 $p$ 为真命题，且关于 $x$ 的方程 $x^2+mx+1=0$ 的两个实数根 $x_1$、$x_2$}
\step{满足 $3x_1+2x_2=5$，则方程 $x^2+mx+1=0$ 有两个不同的正根，}
\step{则 $\begin{cases}\Delta=m^2-4>0\\ x_1+x_2=-m>0\end{cases}$，则 $m<-2$，又 $\begin{cases}x_1+x_2=-m\\ x_1x_2=1\\ 3x_1+2x_2=5\end{cases}$，则 $\begin{cases}x_1=5+2m\\ x_2=-3m-5\end{cases}$，}
\step{则 $(5+2m)(-3m-5)=1$，得 $m=-2$ 或 $m=-\dfrac{13}6$，}
\step{又因为 $m<-2$，所以 $m=-\dfrac{13}6$，所以实数 $m$ 的取值集合为 $\left\{-\dfrac{13}6\right\}$；}
\stepm{（3）}{由（2）得，当 $p$ 为真时，则有 $m<-2$，}
\step{当 $q$ 为真时，$\Delta=16(m-2)^2-16\geqslant0$，解得 $m\leqslant1$ 或 $m\geqslant3$，}
\step{又因为命题 $p$ 和命题 $q$ 只有一个为真，所以 $\begin{cases}m<-2\\ 1<m<3\end{cases}$ 或 $\begin{cases}m\geqslant-2\\ m\leqslant1\ \text{或}\ m\geqslant3\end{cases}$，}
\step{解得 $-2\leqslant m\leqslant1$ 或 $m\geqslant3$，所以实数 $m$ 的取值范围为 $[-2,1]\cup[3,+\infty)$.}
\solend
\fi

\end{enumerate}

\end{document}
"
%
% 原始材料：微信公众号“魔都数学加油站”文章，连载“27学年高一95”；
% 原文标题“27学年高一95——2026-2027年上海市金汇中学高一上10月月考”；
% 原文链接：https://mp.weixin.qq.com/s/H3yH73FTin3V7rxWQ-9fjw。
% 正文为 8 张 PNG 页（`原图/高一95_金汇中学_10月月考/01..08.png`），均 4762×6735，按页序读取。
%
% 篇幅与题号：原卷自第 3 题起编号（第 1、2 题不在原文内）。填空题 3—12（10 题）、
% 选择题 13—16（4 题）、解答题 17—21（5 题），共 19 题。
%
% 讲评版：第 4、5 题只印【答案】（“$m>0$”“$x=0$ 或 $x=1$”），本稿只给 \ans{}；
% 第 13、14 题无【解析】、答案“D”“B”直接印在题干括号内，本稿试卷版按惯例留空、
% 答案版提 \ans{}；其余 15 题印【解析】，其中第 15、16 题的答案在解析末句“故选 $A$／$C$”。
% 第 19 题解析带原卷的分步分值（“（3 分）”“（6 分）”等），照录。全卷无插图。
%
% 副标题：原卷未印副标题。本稿按内容标作“集合与常用逻辑用语、不等式”。
%
% 原卷笔误/存疑（均照录未改，正文对应处另有 % 注）：
%   ① 第 17 题 (1) 解析首句作“**第二个**不等式可化为 $(2x+5)(x-3)\leqslant0$”，而
%      $-2x^2+x+15\geqslant0$ 才是**第一个**不等式；原卷误标，照录；
%   ② 第 19 题 (2) 题干问“$x\in B$ 是 $x\in A$ 的充分非必要条件”，解析却从“$A\cup B=A$，
%      得 $B\subseteq A$”起论（少了“真子集”一层）；原卷照录，结论 $a\leqslant-1$ 未改。
%
% 与原卷的排版差异：
%   ① 原卷每题下直接印【答案】/【解析】，本稿按全稿惯例将答案与解析置于答案版；
%   ② 原卷未印副标题，本稿补出；填空横线统一排 \blankline[4.4em]；
%   ③ 原卷数集记号排斜体 $R$、$N$、$Z$，本稿按全稿惯例改用 $\R$、$\N$、$\Z$；
%   ④ 第 17—21 题的子问用行内“（1）（2）（3）”段，解析中的子问标号用 \stepm；
%   ⑤ 解答题按惯例加 \ansspace；选择题括号排 \mbox{（\quad）}。
%
% 行数核对：逐题数 \step/\stepm 条数，合计 87 条（Q3：2；Q4：0；Q5：0；Q6：1；Q7：2；
% Q8：4；Q9：3；Q10：3；Q11：7；Q12：3；Q13：0；Q14：0；Q15：2；Q16：14；Q17：7；
% Q18：5；Q19：12；Q20：10；Q21：12）。第 4、5 题只印答案，第 13、14 题无解析，
% 均不计入 \step 账。原卷解析的换行按语意合并，未逐行强拆。
%
\documentclass[a4paper,11pt]{article}
\usepackage{common}

\begin{document}

\examtitlest[7]{2026--2027 年金汇中学高一上 10 月月考}{集合与常用逻辑用语、不等式}

\bigsec{一、填空题}

\begin{enumerate}
\setcounter{enumi}{2}

\item 命题：如果 $ab>c$，那么 $a>\dfrac cb$，是\blankline[4.4em]（选填“真”、“假”）命题.

\ifanswer
\solbegin
\step{取 $a=1,b=-2,c=-3$，则 $ab=-2>-3=c$，但 $a=1<\dfrac cb=\dfrac32$，}
\step{所以原命题是假命题.}
\solend
\fi

\item 若“$x=0$”是“$x<m$”的充分条件，则实数 $m$ 的取值范围是\blankline[4.4em].

\ifanswer
\ans{$m>0$}
\fi

\item “$x\neq0$ 且 $x\neq1$”的否定形式为\blankline[4.4em].

\ifanswer
\ans{$x=0$ 或 $x=1$}
\fi

\item 若集合 $M=\{(x,y)\mid x+y=3\}$，$N=\left\{(x,y)\left|y=\dfrac12x\right.\right\}$，则 $M\cap N=\blankline[4.4em]$.

\ifanswer
\solbegin
\step{$x+\dfrac12x=3\Rightarrow x=2,y=1$，所以 $M\cap N=\{(2,1)\}$.}
\solend
\fi

\item 集合 $\{x\in\N\mid0<\lvert x-1\rvert<3\}$ 的真子集个数为\blankline[4.4em].

\ifanswer
\solbegin
\step{因为 $\{x\in\N\mid0<\lvert x-1\rvert<3\}=\{x\in\N\mid-3<x-1<3,x\neq1\}$}
\step{$=\{x\in\N\mid-2<x<4,x\neq1\}=\{0,2,3\}$，所以真子集的个数为 $2^3-1=7$.}
\solend
\fi

\item 设 $A=\{x\mid x^2-3x-4=0\}$，$B=\{x\mid ax-1=0\}$，若 $A\cap B=B$，则实数 $a$ 的值可以为\blankline[4.4em].

\ifanswer
\solbegin
\step{$A=\{-1,4\}$，因为 $A\cap B=B$，所以 $B\subseteq A$，}
\step{①$a=0$ 时，$B=\varnothing$，满足 $B\subseteq A$；}
\step{②$a\neq0$ 时，$B=\left\{\dfrac1a\right\}$，则 $\dfrac1a=-1$ 或 $\dfrac1a=4$，所以 $a=-1$ 或 $\dfrac14$；}
\step{所以 $a$ 的值可以为 $0$、$-1$、$\dfrac14$.}
\solend
\fi

\item 关于 $x$ 的不等式 $ax^2-(a+1)x+1\leqslant0(a<0)$ 的解集为\blankline[4.4em].

\ifanswer
\solbegin
\step{原不等式即 $(ax-1)(x-1)\leqslant0$，因为 $a<0$，所以 $\left(x-\dfrac1a\right)(x-1)\geqslant0$.}
\step{又 $\dfrac1a<0<1$，所以 $x\leqslant\dfrac1a$ 或 $x\geqslant1$，}
\step{故解集为 $\left(-\infty,\dfrac1a\right]\cup[1,+\infty)$.}
\solend
\fi

\item 若不等式 $(a-1)x^2+(2a-2)x-3<0$ 的解集是 $\R$，则 $a$ 的取值范围是\blankline[4.4em].

\ifanswer
\solbegin
\step{当 $a=1$ 时，不等式为 $-3<0$，对任意实数 $x$ 均成立.}
\step{当 $a\neq1$ 时，$\begin{cases}a-1<0\\ \Delta=(2a-2)^2+12(a-1)=4(a-1)(a+2)<0\end{cases}$，解得 $-2<a<1$.}
\step{合并得 $a\in(-2,1]$.}
\solend
\fi

\item 若关于 $x$ 的不等式 $x^2-(m+2)x+2m<0$ 的解集中恰有 $3$ 个整数，则实数 $m$ 的取值范围为\blankline[4.4em].

\ifanswer
\solbegin
\step{将不等式 $x^2-(m+2)x+2m<0$ 因式分解为 $(x-2)(x-m)<0$，}
\step{当 $m<2$ 时，解集为 $(m,2)$.}
\step{要使该区间恰有 $3$ 个整数，这 $3$ 个整数为 $-1,0,1$，故 $m$ 需满足 $-2\leqslant m<-1$.}
\step{当 $m>2$ 时，解集为 $(2,m)$. 要使该区间恰有 $3$ 个整数，}
\step{这 $3$ 个整数为 $3$，$4$，$5$，故 $m$ 需满足 $5<m\leqslant6$.}
\step{当 $m=2$ 时，不等式无解，不满足条件.}
\step{综上，实数 $m$ 的取值范围是 $[-2,-1)\cup(5,6]$.}
\solend
\fi

\item 已知集合 $A=\{x\mid x^2-3x+2\leqslant0\},B=\{y\mid y=x^2-2x+a,x\in\R\}$，若对任何一个 $x\in A$，都有 $x\in B$，则实数 $a$ 的取值范围是\blankline[4.4em].

\ifanswer
\solbegin
\step{$(x-1)(x-2)\leqslant0$，所以 $A=[1,2]$.}
\step{由 $y=(x-1)^2+a-1,x\in\R$，得 $B=[a-1,+\infty)$.}
\step{因此 $A\subseteq B\Leftrightarrow a-1\leqslant1\Leftrightarrow a\leqslant2$，所求范围为 $(-\infty,2]$.}
\solend
\fi

\end{enumerate}

\bigsec{二、选择题}

\begin{enumerate}
\setcounter{enumi}{12}

\item 已知 $a=\sqrt2$，集合 $A=\{x\mid x\leqslant2\}$，则下列表示正确的是\mbox{（\quad）}% 原卷括号内已印“D”，试卷版按惯例留空

\keepwithprev
A. $\{a\}\in A$\qquad B. $a\notin A$\qquad C. $a\subset A$\qquad D. $a\in A$

\ifanswer
\ans{$D$}
\fi

\item 下列四个命题中，真命题是\mbox{（\quad）}% 原卷括号内已印“B”，试卷版按惯例留空

\keepwithprev
A. 若 $a>b$，则 $ac^2>bc^2$\qquad B. 若 $a>b,c>d$，则 $a-d>b-c$

C. 若 $\lvert a\rvert>b$，则 $a^2>b^2$\qquad D. 若 $a>b$，则 $\dfrac1a<\dfrac1b$

\ifanswer
\ans{$B$}
\fi

\item “$x>1$”是“$\dfrac1x<1$”的\mbox{（\quad）}

\keepwithprev
A. 充分不必要条件\qquad B. 必要不充分条件

C. 充要条件\qquad D. 既不充分也不必要条件

\ifanswer
\solbegin
\step{$x>1\Rightarrow0<\dfrac1x<1$，充分性成立.}
\step{当 $x=-1$ 时，$\dfrac1x=-1<1$，但 $x>1$ 不成立，必要性不成立. 故选 $A$.}
\solend
\fi

\item 对于集合 $M=\{a\mid a=x^2-y^2,x\in\Z,y\in\Z\}$，给出如下三个结论：①如果 $P=\{b\mid b=2n+1,n\in\Z\}$，那么 $P\subseteq M$；②如果 $c=4n+2,n\in\Z$，那么 $c\notin M$；③如果 $a_1\in M,a_2\in M$，那么 $a_1a_2\in M$. 其中正确结论的个数是\mbox{（\quad）}

\keepwithprev
A. $1$\qquad B. $2$\qquad C. $3$\qquad D. $0$

\ifanswer
\solbegin
\step{对于①：因为 $2n+1=(n+1)^2-n^2,n\in\Z$，所以 $b=2n+1\in M$，故 $P\subseteq M$，}
\step{故①正确；}
\step{对于②：因为 $c=4n+2=2(2n+1),n\in\Z$，所以 $c$ 为偶数，且不能被 $4$ 整除，}
\step{若 $c\in M$，则存在 $x\in\Z,y\in\Z$，得 $c=x^2-y^2=(x+y)(x-y),x\in\Z,y\in\Z$，}
\step{因为 $x+y$ 和 $x-y$ 同奇或同偶，}
\step{若 $x+y$ 和 $x-y$ 同奇，则 $c=(x+y)(x-y)$ 为奇数，矛盾，不符合，}
\step{若 $x+y$ 和 $x-y$ 同偶，则 $c=(x+y)(x-y)$ 能被 $4$ 整除，矛盾，不符合，}
\step{所以 $c\notin M$，故②正确；}
\step{对于③：因为 $a_1\in M,a_2\in M$，}
\step{所以存在 $x_1,x_2\in\Z,y_1,y_2\in\Z$ 使得 $a_1=x_1^2-y_1^2,a_2=x_2^2-y_2^2$，}
\step{所以 $a_1a_2=(x_1^2-y_1^2)(x_2^2-y_2^2)=(x_1x_2)^2+(y_1y_2)^2-(x_1y_2)^2-(x_2y_1)^2$}
\step{$=(x_1x_2+y_1y_2)^2-(x_1y_2+x_2y_1)^2$，}
\step{因为 $x_1x_2+y_1y_2,x_1y_2+x_2y_1\in\Z$，所以 $a_1a_2\in M$，故③正确.}
\step{故选 $C$.}
\solend
\fi

\end{enumerate}

\bigsec{三、解答题}

\begin{enumerate}
\setcounter{enumi}{16}

\item （1）求不等式组 $\begin{cases}-2x^2+x+15\geqslant0\\ (x+1)(3x-1)>0\end{cases}$ 的解集；

（2）讨论解关于 $x$ 的不等式 $mx+4<m^2+2x$.

\ansspace*[12]

\ifanswer
\solbegin
\stepm{（1）}{第二个不等式可化为 $(2x+5)(x-3)\leqslant0$，解得 $-\dfrac52\leqslant x\leqslant3$.}% 原卷把“第一个”误作“第二个”，照录
\step{第二个不等式 $(x+1)(3x-1)>0$ 的解集为 $(-\infty,-1)\cup\left(\dfrac13,+\infty\right)$.}
\step{原不等式组的解集为 $\left[-\dfrac52,-1\right)\cup\left(\dfrac13,3\right]$.}
\stepm{（2）}{移项并因式分解，得 $(m-2)x<m^2-4\Leftrightarrow(m-2)x<(m-2)(m+2)$.}
\step{当 $m>2$ 时，$m-2>0$，解集为 $(-\infty,m+2)$.}
\step{当 $m<2$ 时，$m-2<0$，解集为 $(m+2,+\infty)$.}
\step{当 $m=2$ 时，不等式化为 $0<0$，不成立，解集为 $\varnothing$.}
\solend
\fi

\item （1）设 $x$ 是实数，比较 $x^2+9$ 与 $6x$ 的值的大小；

（2）已知 $a>0$ 且 $a\neq1$，比较 $\dfrac{a^2+1}{a^2-1}$ 与 $\dfrac{a+1}{a-1}$ 的大小.

\ansspace*[11]

\ifanswer
\solbegin
\stepm{（1）}{作差得 $(x^2+9)-6x=(x-3)^2\geqslant0$，因此 $x^2+9\geqslant6x$；}
\stepm{（2）}{$\dfrac{a^2+1}{a^2-1}-\dfrac{a+1}{a-1}=\dfrac{-2a}{a^2-1}$，}
\step{当 $a>1$ 时，$-2a<0$，$a^2-1>0$，则 $\dfrac{-2a}{a^2-1}<0$，即 $\dfrac{a^2+1}{a^2-1}<\dfrac{a+1}{a-1}$；}
\step{当 $0<a<1$ 时，$-2a<0$，$a^2-1<0$，则 $\dfrac{-2a}{a^2-1}>0$，即 $\dfrac{a^2+1}{a^2-1}>\dfrac{a+1}{a-1}$；}
\step{综上，$a>1$ 时，$\dfrac{a^2+1}{a^2-1}<\dfrac{a+1}{a-1}$；$0<a<1$ 时，$\dfrac{a^2+1}{a^2-1}>\dfrac{a+1}{a-1}$.}
\solend
\fi

\item 已知 $a\in\R$，集合 $A=\left\{x\left|\dfrac{2-x}{x+1}\leqslant0\right.\right\},B=\{x\mid x^2+(3-a)x-3a<0\}$.

（1）当 $a=\dfrac13$ 时，求 $A$ 和 $B$；

（2）“$x\in B$”是“$x\in A$”的充分非必要条件，求实数 $a$ 的取值范围.

\ansspace*[13]

\ifanswer
\solbegin
\stepm{（1）}{由 $\dfrac{2-x}{x+1}\leqslant0$，得 $\begin{cases}(2-x)(x+1)\leqslant0\\ x\neq-1\end{cases}$，解得 $x<-1$ 或 $x\geqslant2$，}
\step{则 $A=\{x\mid x<-1$ 或 $x\geqslant2\}$；$\cdots\cdots\cdots\cdots\cdots\cdots\cdots\cdots\cdots\cdots\cdots\cdots\cdots\cdots\cdots\cdots$（3 分）}
\step{$B=\{x\mid x^2+(3-a)x-3a<0\}=\{x\mid(x+3)(x-a)<0\}$，}
\step{当 $a=\dfrac13$ 时，$B=\left\{x\left|(x+2)\left(x-\dfrac13\right)<0\right.\right\}$，}
\step{解得 $B=\left\{x\left|-3<x<\dfrac13\right.\right\}\cdots\cdots\cdots\cdots\cdots\cdots$（6 分）}
\stepm{（2）}{由 $A\cup B=A$，得 $B\subseteq A$，}% 原卷从 A∪B=A 起论，与题干“充分非必要”差一层，照录
\step{①当 $B=\varnothing$ 时，得 $a=-3$，符合题意；$\cdots\cdots\cdots\cdots\cdots$（2 分）}
\step{②当 $B\neq\varnothing$ 时，若 $a>-3$，则 $B=\{x\mid-3<x<a\}$，}
\step{由 $B\subseteq A$，得 $a\leqslant-1$，此时 $-3<a\leqslant-1$；$\cdots\cdots\cdots\cdots\cdots$（4 分）}
\step{若 $a<-3$，则 $B=\{x\mid a<x<-3\}$，此时 $B\subseteq A$ 恒成立，}
\step{故 $a<-3$ 符合题意；$\cdots\cdots\cdots$（6 分）}
\step{综上所述，实数 $a$ 的取值范围为 $a\leqslant-1$. $\cdots\cdots\cdots\cdots\cdots$（8 分）}
\solend
\fi

\item 已知集合 $A=\{x\mid x<-2$ 或 $x\geqslant3\}$，$U=R$.

（1）求 $\overline A$；

（2）若 $B=\{x\mid mx-1>0\}$，当 $B\subseteq A$ 时，求实数 $m$ 的取值范围；

（3）若 $B=\{x\mid2m-6\leqslant x\leqslant m+4\}$，当 $B\cap\overline A=\varnothing$ 时，求实数 $m$ 的取值范围.

\ansspace*[14]

\ifanswer
\solbegin
\stepm{（1）}{因为集合 $A=\{x\mid x<-2$ 或 $x\geqslant3\}$，$U=R$，所以 $\overline A=\{x\mid-2\leqslant x<3\}$；}
\stepm{（2）}{当 $m=0$ 时，$B=\varnothing$，符合题意，}
\step{当 $m>0$ 时，$B=\left\{x\left|x>\dfrac1m\right.\right\}$，则 $\dfrac1m\geqslant3$，解得 $0<m\leqslant\dfrac13$，}
\step{当 $m<0$ 时，$B=\left\{x\left|x<\dfrac1m\right.\right\}$，则 $\dfrac1m\leqslant-2$，解得 $-\dfrac12\leqslant m<0$，}
\step{综上所述，实数 $m$ 的取值范围 $\left[-\dfrac12,\dfrac13\right]$；}
\stepm{（3）}{由题意得 $B=\{x\mid2m-6\leqslant x\leqslant m+4\}$，$\overline A=\{x\mid-2\leqslant x<3\}$，且 $B\cap\overline A=\varnothing$，}
\step{当 $B=\varnothing$ 时，$2m-6>m+4$，解得 $m>10$，}
\step{当 $B\neq\varnothing$ 时，则 $\begin{cases}2m-6\leqslant m+4\\ m+4<-2\end{cases}$ 或 $\begin{cases}2m-6\leqslant m+4\\ 2m-6\geqslant3\end{cases}$，}
\step{解得 $m<-6$ 或 $\dfrac92\leqslant m\leqslant10$，}
\step{综上所述，实数 $m$ 的取值范围为 $(-\infty,-6)\cup\left[\dfrac92,+\infty\right)$.}
\solend
\fi

\item 已知命题 $p$：关于 $x$ 的方程 $x^2+mx+1=0$ 有两个不同的正根；命题 $q$：关于 $x$ 的方程 $4x^2+4(m-2)x=-1$ 有实根.

（1）若命题 $q$ 为假命题，求整数 $m$ 的解集；

（2）若命题 $p$ 为真命题，且关于 $x$ 的方程 $x^2+mx+1=0$ 的两个实数根 $x_1$、$x_2$ 满足 $3x_1+2x_2=5$，求实数 $m$ 的取值集合；

（3）若命题 $p$ 和命题 $q$ 有且仅有一个成立，求实数 $m$ 的取值范围.

\ansspace*[16]

\ifanswer
\solbegin
\stepm{（1）}{命题 $q$：关于 $x$ 的方程 $4x^2+4(m-2)x=-1$ 有实根为假命题，}
\step{即 $4x^2+4(m-2)x+1=0$ 无实根，}
\step{则 $\Delta=16(m-2)^2-16<0$，则 $1<m<3$，则整数 $m$ 的解集为 $\{2\}$；}
\stepm{（2）}{若命题 $p$ 为真命题，且关于 $x$ 的方程 $x^2+mx+1=0$ 的两个实数根 $x_1$、$x_2$}
\step{满足 $3x_1+2x_2=5$，则方程 $x^2+mx+1=0$ 有两个不同的正根，}
\step{则 $\begin{cases}\Delta=m^2-4>0\\ x_1+x_2=-m>0\end{cases}$，则 $m<-2$，又 $\begin{cases}x_1+x_2=-m\\ x_1x_2=1\\ 3x_1+2x_2=5\end{cases}$，则 $\begin{cases}x_1=5+2m\\ x_2=-3m-5\end{cases}$，}
\step{则 $(5+2m)(-3m-5)=1$，得 $m=-2$ 或 $m=-\dfrac{13}6$，}
\step{又因为 $m<-2$，所以 $m=-\dfrac{13}6$，所以实数 $m$ 的取值集合为 $\left\{-\dfrac{13}6\right\}$；}
\stepm{（3）}{由（2）得，当 $p$ 为真时，则有 $m<-2$，}
\step{当 $q$ 为真时，$\Delta=16(m-2)^2-16\geqslant0$，解得 $m\leqslant1$ 或 $m\geqslant3$，}
\step{又因为命题 $p$ 和命题 $q$ 只有一个为真，所以 $\begin{cases}m<-2\\ 1<m<3\end{cases}$ 或 $\begin{cases}m\geqslant-2\\ m\leqslant1\ \text{或}\ m\geqslant3\end{cases}$，}
\step{解得 $-2\leqslant m\leqslant1$ 或 $m\geqslant3$，所以实数 $m$ 的取值范围为 $[-2,1]\cup[3,+\infty)$.}
\solend
\fi

\end{enumerate}

\end{document}
"
%
% 原始材料：微信公众号“魔都数学加油站”文章，连载“27学年高一95”；
% 原文标题“27学年高一95——2026-2027年上海市金汇中学高一上10月月考”；
% 原文链接：https://mp.weixin.qq.com/s/H3yH73FTin3V7rxWQ-9fjw。
% 正文为 8 张 PNG 页（`原图/高一95_金汇中学_10月月考/01..08.png`），均 4762×6735，按页序读取。
%
% 篇幅与题号：原卷自第 3 题起编号（第 1、2 题不在原文内）。填空题 3—12（10 题）、
% 选择题 13—16（4 题）、解答题 17—21（5 题），共 19 题。
%
% 讲评版：第 4、5 题只印【答案】（“$m>0$”“$x=0$ 或 $x=1$”），本稿只给 \ans{}；
% 第 13、14 题无【解析】、答案“D”“B”直接印在题干括号内，本稿试卷版按惯例留空、
% 答案版提 \ans{}；其余 15 题印【解析】，其中第 15、16 题的答案在解析末句“故选 $A$／$C$”。
% 第 19 题解析带原卷的分步分值（“（3 分）”“（6 分）”等），照录。全卷无插图。
%
% 副标题：原卷未印副标题。本稿按内容标作“集合与常用逻辑用语、不等式”。
%
% 原卷笔误/存疑（均照录未改，正文对应处另有 % 注）：
%   ① 第 17 题 (1) 解析首句作“**第二个**不等式可化为 $(2x+5)(x-3)\leqslant0$”，而
%      $-2x^2+x+15\geqslant0$ 才是**第一个**不等式；原卷误标，照录；
%   ② 第 19 题 (2) 题干问“$x\in B$ 是 $x\in A$ 的充分非必要条件”，解析却从“$A\cup B=A$，
%      得 $B\subseteq A$”起论（少了“真子集”一层）；原卷照录，结论 $a\leqslant-1$ 未改。
%
% 与原卷的排版差异：
%   ① 原卷每题下直接印【答案】/【解析】，本稿按全稿惯例将答案与解析置于答案版；
%   ② 原卷未印副标题，本稿补出；填空横线统一排 \blankline[4.4em]；
%   ③ 原卷数集记号排斜体 $R$、$N$、$Z$，本稿按全稿惯例改用 $\R$、$\N$、$\Z$；
%   ④ 第 17—21 题的子问用行内“（1）（2）（3）”段，解析中的子问标号用 \stepm；
%   ⑤ 解答题按惯例加 \ansspace；选择题括号排 \mbox{（\quad）}。
%
% 行数核对：逐题数 \step/\stepm 条数，合计 87 条（Q3：2；Q4：0；Q5：0；Q6：1；Q7：2；
% Q8：4；Q9：3；Q10：3；Q11：7；Q12：3；Q13：0；Q14：0；Q15：2；Q16：14；Q17：7；
% Q18：5；Q19：12；Q20：10；Q21：12）。第 4、5 题只印答案，第 13、14 题无解析，
% 均不计入 \step 账。原卷解析的换行按语意合并，未逐行强拆。
%
\documentclass[a4paper,11pt]{article}
\usepackage{common}

\begin{document}

\examtitlest[7]{2026--2027 年金汇中学高一上 10 月月考}{集合与常用逻辑用语、不等式}

\bigsec{一、填空题}

\begin{enumerate}
\setcounter{enumi}{2}

\item 命题：如果 $ab>c$，那么 $a>\dfrac cb$，是\blankline[4.4em]（选填“真”、“假”）命题.

\ifanswer
\solbegin
\step{取 $a=1,b=-2,c=-3$，则 $ab=-2>-3=c$，但 $a=1<\dfrac cb=\dfrac32$，}
\step{所以原命题是假命题.}
\solend
\fi

\item 若“$x=0$”是“$x<m$”的充分条件，则实数 $m$ 的取值范围是\blankline[4.4em].

\ifanswer
\ans{$m>0$}
\fi

\item “$x\neq0$ 且 $x\neq1$”的否定形式为\blankline[4.4em].

\ifanswer
\ans{$x=0$ 或 $x=1$}
\fi

\item 若集合 $M=\{(x,y)\mid x+y=3\}$，$N=\left\{(x,y)\left|y=\dfrac12x\right.\right\}$，则 $M\cap N=\blankline[4.4em]$.

\ifanswer
\solbegin
\step{$x+\dfrac12x=3\Rightarrow x=2,y=1$，所以 $M\cap N=\{(2,1)\}$.}
\solend
\fi

\item 集合 $\{x\in\N\mid0<\lvert x-1\rvert<3\}$ 的真子集个数为\blankline[4.4em].

\ifanswer
\solbegin
\step{因为 $\{x\in\N\mid0<\lvert x-1\rvert<3\}=\{x\in\N\mid-3<x-1<3,x\neq1\}$}
\step{$=\{x\in\N\mid-2<x<4,x\neq1\}=\{0,2,3\}$，所以真子集的个数为 $2^3-1=7$.}
\solend
\fi

\item 设 $A=\{x\mid x^2-3x-4=0\}$，$B=\{x\mid ax-1=0\}$，若 $A\cap B=B$，则实数 $a$ 的值可以为\blankline[4.4em].

\ifanswer
\solbegin
\step{$A=\{-1,4\}$，因为 $A\cap B=B$，所以 $B\subseteq A$，}
\step{①$a=0$ 时，$B=\varnothing$，满足 $B\subseteq A$；}
\step{②$a\neq0$ 时，$B=\left\{\dfrac1a\right\}$，则 $\dfrac1a=-1$ 或 $\dfrac1a=4$，所以 $a=-1$ 或 $\dfrac14$；}
\step{所以 $a$ 的值可以为 $0$、$-1$、$\dfrac14$.}
\solend
\fi

\item 关于 $x$ 的不等式 $ax^2-(a+1)x+1\leqslant0(a<0)$ 的解集为\blankline[4.4em].

\ifanswer
\solbegin
\step{原不等式即 $(ax-1)(x-1)\leqslant0$，因为 $a<0$，所以 $\left(x-\dfrac1a\right)(x-1)\geqslant0$.}
\step{又 $\dfrac1a<0<1$，所以 $x\leqslant\dfrac1a$ 或 $x\geqslant1$，}
\step{故解集为 $\left(-\infty,\dfrac1a\right]\cup[1,+\infty)$.}
\solend
\fi

\item 若不等式 $(a-1)x^2+(2a-2)x-3<0$ 的解集是 $\R$，则 $a$ 的取值范围是\blankline[4.4em].

\ifanswer
\solbegin
\step{当 $a=1$ 时，不等式为 $-3<0$，对任意实数 $x$ 均成立.}
\step{当 $a\neq1$ 时，$\begin{cases}a-1<0\\ \Delta=(2a-2)^2+12(a-1)=4(a-1)(a+2)<0\end{cases}$，解得 $-2<a<1$.}
\step{合并得 $a\in(-2,1]$.}
\solend
\fi

\item 若关于 $x$ 的不等式 $x^2-(m+2)x+2m<0$ 的解集中恰有 $3$ 个整数，则实数 $m$ 的取值范围为\blankline[4.4em].

\ifanswer
\solbegin
\step{将不等式 $x^2-(m+2)x+2m<0$ 因式分解为 $(x-2)(x-m)<0$，}
\step{当 $m<2$ 时，解集为 $(m,2)$.}
\step{要使该区间恰有 $3$ 个整数，这 $3$ 个整数为 $-1,0,1$，故 $m$ 需满足 $-2\leqslant m<-1$.}
\step{当 $m>2$ 时，解集为 $(2,m)$. 要使该区间恰有 $3$ 个整数，}
\step{这 $3$ 个整数为 $3$，$4$，$5$，故 $m$ 需满足 $5<m\leqslant6$.}
\step{当 $m=2$ 时，不等式无解，不满足条件.}
\step{综上，实数 $m$ 的取值范围是 $[-2,-1)\cup(5,6]$.}
\solend
\fi

\item 已知集合 $A=\{x\mid x^2-3x+2\leqslant0\},B=\{y\mid y=x^2-2x+a,x\in\R\}$，若对任何一个 $x\in A$，都有 $x\in B$，则实数 $a$ 的取值范围是\blankline[4.4em].

\ifanswer
\solbegin
\step{$(x-1)(x-2)\leqslant0$，所以 $A=[1,2]$.}
\step{由 $y=(x-1)^2+a-1,x\in\R$，得 $B=[a-1,+\infty)$.}
\step{因此 $A\subseteq B\Leftrightarrow a-1\leqslant1\Leftrightarrow a\leqslant2$，所求范围为 $(-\infty,2]$.}
\solend
\fi

\end{enumerate}

\bigsec{二、选择题}

\begin{enumerate}
\setcounter{enumi}{12}

\item 已知 $a=\sqrt2$，集合 $A=\{x\mid x\leqslant2\}$，则下列表示正确的是\mbox{（\quad）}% 原卷括号内已印“D”，试卷版按惯例留空

\keepwithprev
A. $\{a\}\in A$\qquad B. $a\notin A$\qquad C. $a\subset A$\qquad D. $a\in A$

\ifanswer
\ans{$D$}
\fi

\item 下列四个命题中，真命题是\mbox{（\quad）}% 原卷括号内已印“B”，试卷版按惯例留空

\keepwithprev
A. 若 $a>b$，则 $ac^2>bc^2$\qquad B. 若 $a>b,c>d$，则 $a-d>b-c$

C. 若 $\lvert a\rvert>b$，则 $a^2>b^2$\qquad D. 若 $a>b$，则 $\dfrac1a<\dfrac1b$

\ifanswer
\ans{$B$}
\fi

\item “$x>1$”是“$\dfrac1x<1$”的\mbox{（\quad）}

\keepwithprev
A. 充分不必要条件\qquad B. 必要不充分条件

C. 充要条件\qquad D. 既不充分也不必要条件

\ifanswer
\solbegin
\step{$x>1\Rightarrow0<\dfrac1x<1$，充分性成立.}
\step{当 $x=-1$ 时，$\dfrac1x=-1<1$，但 $x>1$ 不成立，必要性不成立. 故选 $A$.}
\solend
\fi

\item 对于集合 $M=\{a\mid a=x^2-y^2,x\in\Z,y\in\Z\}$，给出如下三个结论：①如果 $P=\{b\mid b=2n+1,n\in\Z\}$，那么 $P\subseteq M$；②如果 $c=4n+2,n\in\Z$，那么 $c\notin M$；③如果 $a_1\in M,a_2\in M$，那么 $a_1a_2\in M$. 其中正确结论的个数是\mbox{（\quad）}

\keepwithprev
A. $1$\qquad B. $2$\qquad C. $3$\qquad D. $0$

\ifanswer
\solbegin
\step{对于①：因为 $2n+1=(n+1)^2-n^2,n\in\Z$，所以 $b=2n+1\in M$，故 $P\subseteq M$，}
\step{故①正确；}
\step{对于②：因为 $c=4n+2=2(2n+1),n\in\Z$，所以 $c$ 为偶数，且不能被 $4$ 整除，}
\step{若 $c\in M$，则存在 $x\in\Z,y\in\Z$，得 $c=x^2-y^2=(x+y)(x-y),x\in\Z,y\in\Z$，}
\step{因为 $x+y$ 和 $x-y$ 同奇或同偶，}
\step{若 $x+y$ 和 $x-y$ 同奇，则 $c=(x+y)(x-y)$ 为奇数，矛盾，不符合，}
\step{若 $x+y$ 和 $x-y$ 同偶，则 $c=(x+y)(x-y)$ 能被 $4$ 整除，矛盾，不符合，}
\step{所以 $c\notin M$，故②正确；}
\step{对于③：因为 $a_1\in M,a_2\in M$，}
\step{所以存在 $x_1,x_2\in\Z,y_1,y_2\in\Z$ 使得 $a_1=x_1^2-y_1^2,a_2=x_2^2-y_2^2$，}
\step{所以 $a_1a_2=(x_1^2-y_1^2)(x_2^2-y_2^2)=(x_1x_2)^2+(y_1y_2)^2-(x_1y_2)^2-(x_2y_1)^2$}
\step{$=(x_1x_2+y_1y_2)^2-(x_1y_2+x_2y_1)^2$，}
\step{因为 $x_1x_2+y_1y_2,x_1y_2+x_2y_1\in\Z$，所以 $a_1a_2\in M$，故③正确.}
\step{故选 $C$.}
\solend
\fi

\end{enumerate}

\bigsec{三、解答题}

\begin{enumerate}
\setcounter{enumi}{16}

\item （1）求不等式组 $\begin{cases}-2x^2+x+15\geqslant0\\ (x+1)(3x-1)>0\end{cases}$ 的解集；

（2）讨论解关于 $x$ 的不等式 $mx+4<m^2+2x$.

\ansspace*[12]

\ifanswer
\solbegin
\stepm{（1）}{第二个不等式可化为 $(2x+5)(x-3)\leqslant0$，解得 $-\dfrac52\leqslant x\leqslant3$.}% 原卷把“第一个”误作“第二个”，照录
\step{第二个不等式 $(x+1)(3x-1)>0$ 的解集为 $(-\infty,-1)\cup\left(\dfrac13,+\infty\right)$.}
\step{原不等式组的解集为 $\left[-\dfrac52,-1\right)\cup\left(\dfrac13,3\right]$.}
\stepm{（2）}{移项并因式分解，得 $(m-2)x<m^2-4\Leftrightarrow(m-2)x<(m-2)(m+2)$.}
\step{当 $m>2$ 时，$m-2>0$，解集为 $(-\infty,m+2)$.}
\step{当 $m<2$ 时，$m-2<0$，解集为 $(m+2,+\infty)$.}
\step{当 $m=2$ 时，不等式化为 $0<0$，不成立，解集为 $\varnothing$.}
\solend
\fi

\item （1）设 $x$ 是实数，比较 $x^2+9$ 与 $6x$ 的值的大小；

（2）已知 $a>0$ 且 $a\neq1$，比较 $\dfrac{a^2+1}{a^2-1}$ 与 $\dfrac{a+1}{a-1}$ 的大小.

\ansspace*[11]

\ifanswer
\solbegin
\stepm{（1）}{作差得 $(x^2+9)-6x=(x-3)^2\geqslant0$，因此 $x^2+9\geqslant6x$；}
\stepm{（2）}{$\dfrac{a^2+1}{a^2-1}-\dfrac{a+1}{a-1}=\dfrac{-2a}{a^2-1}$，}
\step{当 $a>1$ 时，$-2a<0$，$a^2-1>0$，则 $\dfrac{-2a}{a^2-1}<0$，即 $\dfrac{a^2+1}{a^2-1}<\dfrac{a+1}{a-1}$；}
\step{当 $0<a<1$ 时，$-2a<0$，$a^2-1<0$，则 $\dfrac{-2a}{a^2-1}>0$，即 $\dfrac{a^2+1}{a^2-1}>\dfrac{a+1}{a-1}$；}
\step{综上，$a>1$ 时，$\dfrac{a^2+1}{a^2-1}<\dfrac{a+1}{a-1}$；$0<a<1$ 时，$\dfrac{a^2+1}{a^2-1}>\dfrac{a+1}{a-1}$.}
\solend
\fi

\item 已知 $a\in\R$，集合 $A=\left\{x\left|\dfrac{2-x}{x+1}\leqslant0\right.\right\},B=\{x\mid x^2+(3-a)x-3a<0\}$.

（1）当 $a=\dfrac13$ 时，求 $A$ 和 $B$；

（2）“$x\in B$”是“$x\in A$”的充分非必要条件，求实数 $a$ 的取值范围.

\ansspace*[13]

\ifanswer
\solbegin
\stepm{（1）}{由 $\dfrac{2-x}{x+1}\leqslant0$，得 $\begin{cases}(2-x)(x+1)\leqslant0\\ x\neq-1\end{cases}$，解得 $x<-1$ 或 $x\geqslant2$，}
\step{则 $A=\{x\mid x<-1$ 或 $x\geqslant2\}$；$\cdots\cdots\cdots\cdots\cdots\cdots\cdots\cdots\cdots\cdots\cdots\cdots\cdots\cdots\cdots\cdots$（3 分）}
\step{$B=\{x\mid x^2+(3-a)x-3a<0\}=\{x\mid(x+3)(x-a)<0\}$，}
\step{当 $a=\dfrac13$ 时，$B=\left\{x\left|(x+2)\left(x-\dfrac13\right)<0\right.\right\}$，}
\step{解得 $B=\left\{x\left|-3<x<\dfrac13\right.\right\}\cdots\cdots\cdots\cdots\cdots\cdots$（6 分）}
\stepm{（2）}{由 $A\cup B=A$，得 $B\subseteq A$，}% 原卷从 A∪B=A 起论，与题干“充分非必要”差一层，照录
\step{①当 $B=\varnothing$ 时，得 $a=-3$，符合题意；$\cdots\cdots\cdots\cdots\cdots$（2 分）}
\step{②当 $B\neq\varnothing$ 时，若 $a>-3$，则 $B=\{x\mid-3<x<a\}$，}
\step{由 $B\subseteq A$，得 $a\leqslant-1$，此时 $-3<a\leqslant-1$；$\cdots\cdots\cdots\cdots\cdots$（4 分）}
\step{若 $a<-3$，则 $B=\{x\mid a<x<-3\}$，此时 $B\subseteq A$ 恒成立，}
\step{故 $a<-3$ 符合题意；$\cdots\cdots\cdots$（6 分）}
\step{综上所述，实数 $a$ 的取值范围为 $a\leqslant-1$. $\cdots\cdots\cdots\cdots\cdots$（8 分）}
\solend
\fi

\item 已知集合 $A=\{x\mid x<-2$ 或 $x\geqslant3\}$，$U=R$.

（1）求 $\overline A$；

（2）若 $B=\{x\mid mx-1>0\}$，当 $B\subseteq A$ 时，求实数 $m$ 的取值范围；

（3）若 $B=\{x\mid2m-6\leqslant x\leqslant m+4\}$，当 $B\cap\overline A=\varnothing$ 时，求实数 $m$ 的取值范围.

\ansspace*[14]

\ifanswer
\solbegin
\stepm{（1）}{因为集合 $A=\{x\mid x<-2$ 或 $x\geqslant3\}$，$U=R$，所以 $\overline A=\{x\mid-2\leqslant x<3\}$；}
\stepm{（2）}{当 $m=0$ 时，$B=\varnothing$，符合题意，}
\step{当 $m>0$ 时，$B=\left\{x\left|x>\dfrac1m\right.\right\}$，则 $\dfrac1m\geqslant3$，解得 $0<m\leqslant\dfrac13$，}
\step{当 $m<0$ 时，$B=\left\{x\left|x<\dfrac1m\right.\right\}$，则 $\dfrac1m\leqslant-2$，解得 $-\dfrac12\leqslant m<0$，}
\step{综上所述，实数 $m$ 的取值范围 $\left[-\dfrac12,\dfrac13\right]$；}
\stepm{（3）}{由题意得 $B=\{x\mid2m-6\leqslant x\leqslant m+4\}$，$\overline A=\{x\mid-2\leqslant x<3\}$，且 $B\cap\overline A=\varnothing$，}
\step{当 $B=\varnothing$ 时，$2m-6>m+4$，解得 $m>10$，}
\step{当 $B\neq\varnothing$ 时，则 $\begin{cases}2m-6\leqslant m+4\\ m+4<-2\end{cases}$ 或 $\begin{cases}2m-6\leqslant m+4\\ 2m-6\geqslant3\end{cases}$，}
\step{解得 $m<-6$ 或 $\dfrac92\leqslant m\leqslant10$，}
\step{综上所述，实数 $m$ 的取值范围为 $(-\infty,-6)\cup\left[\dfrac92,+\infty\right)$.}
\solend
\fi

\item 已知命题 $p$：关于 $x$ 的方程 $x^2+mx+1=0$ 有两个不同的正根；命题 $q$：关于 $x$ 的方程 $4x^2+4(m-2)x=-1$ 有实根.

（1）若命题 $q$ 为假命题，求整数 $m$ 的解集；

（2）若命题 $p$ 为真命题，且关于 $x$ 的方程 $x^2+mx+1=0$ 的两个实数根 $x_1$、$x_2$ 满足 $3x_1+2x_2=5$，求实数 $m$ 的取值集合；

（3）若命题 $p$ 和命题 $q$ 有且仅有一个成立，求实数 $m$ 的取值范围.

\ansspace*[16]

\ifanswer
\solbegin
\stepm{（1）}{命题 $q$：关于 $x$ 的方程 $4x^2+4(m-2)x=-1$ 有实根为假命题，}
\step{即 $4x^2+4(m-2)x+1=0$ 无实根，}
\step{则 $\Delta=16(m-2)^2-16<0$，则 $1<m<3$，则整数 $m$ 的解集为 $\{2\}$；}
\stepm{（2）}{若命题 $p$ 为真命题，且关于 $x$ 的方程 $x^2+mx+1=0$ 的两个实数根 $x_1$、$x_2$}
\step{满足 $3x_1+2x_2=5$，则方程 $x^2+mx+1=0$ 有两个不同的正根，}
\step{则 $\begin{cases}\Delta=m^2-4>0\\ x_1+x_2=-m>0\end{cases}$，则 $m<-2$，又 $\begin{cases}x_1+x_2=-m\\ x_1x_2=1\\ 3x_1+2x_2=5\end{cases}$，则 $\begin{cases}x_1=5+2m\\ x_2=-3m-5\end{cases}$，}
\step{则 $(5+2m)(-3m-5)=1$，得 $m=-2$ 或 $m=-\dfrac{13}6$，}
\step{又因为 $m<-2$，所以 $m=-\dfrac{13}6$，所以实数 $m$ 的取值集合为 $\left\{-\dfrac{13}6\right\}$；}
\stepm{（3）}{由（2）得，当 $p$ 为真时，则有 $m<-2$，}
\step{当 $q$ 为真时，$\Delta=16(m-2)^2-16\geqslant0$，解得 $m\leqslant1$ 或 $m\geqslant3$，}
\step{又因为命题 $p$ 和命题 $q$ 只有一个为真，所以 $\begin{cases}m<-2\\ 1<m<3\end{cases}$ 或 $\begin{cases}m\geqslant-2\\ m\leqslant1\ \text{或}\ m\geqslant3\end{cases}$，}
\step{解得 $-2\leqslant m\leqslant1$ 或 $m\geqslant3$，所以实数 $m$ 的取值范围为 $[-2,1]\cup[3,+\infty)$.}
\solend
\fi

\end{enumerate}

\end{document}
"
%
% 原始材料：微信公众号“魔都数学加油站”文章，连载“27学年高一95”；
% 原文标题“27学年高一95——2026-2027年上海市金汇中学高一上10月月考”；
% 原文链接：https://mp.weixin.qq.com/s/H3yH73FTin3V7rxWQ-9fjw。
% 正文为 8 张 PNG 页（`原图/高一95_金汇中学_10月月考/01..08.png`），均 4762×6735，按页序读取。
%
% 篇幅与题号：原卷自第 3 题起编号（第 1、2 题不在原文内）。填空题 3—12（10 题）、
% 选择题 13—16（4 题）、解答题 17—21（5 题），共 19 题。
%
% 讲评版：第 4、5 题只印【答案】（“$m>0$”“$x=0$ 或 $x=1$”），本稿只给 \ans{}；
% 第 13、14 题无【解析】、答案“D”“B”直接印在题干括号内，本稿试卷版按惯例留空、
% 答案版提 \ans{}；其余 15 题印【解析】，其中第 15、16 题的答案在解析末句“故选 $A$／$C$”。
% 第 19 题解析带原卷的分步分值（“（3 分）”“（6 分）”等），照录。全卷无插图。
%
% 副标题：原卷未印副标题。本稿按内容标作“集合与常用逻辑用语、不等式”。
%
% 原卷笔误/存疑（均照录未改，正文对应处另有 % 注）：
%   ① 第 17 题 (1) 解析首句作“**第二个**不等式可化为 $(2x+5)(x-3)\leqslant0$”，而
%      $-2x^2+x+15\geqslant0$ 才是**第一个**不等式；原卷误标，照录；
%   ② 第 19 题 (2) 题干问“$x\in B$ 是 $x\in A$ 的充分非必要条件”，解析却从“$A\cup B=A$，
%      得 $B\subseteq A$”起论（少了“真子集”一层）；原卷照录，结论 $a\leqslant-1$ 未改。
%
% 与原卷的排版差异：
%   ① 原卷每题下直接印【答案】/【解析】，本稿按全稿惯例将答案与解析置于答案版；
%   ② 原卷未印副标题，本稿补出；填空横线统一排 \blankline[4.4em]；
%   ③ 原卷数集记号排斜体 $R$、$N$、$Z$，本稿按全稿惯例改用 $\R$、$\N$、$\Z$；
%   ④ 第 17—21 题的子问用行内“（1）（2）（3）”段，解析中的子问标号用 \stepm；
%   ⑤ 解答题按惯例加 \ansspace；选择题括号排 \mbox{（\quad）}。
%
% 行数核对：逐题数 \step/\stepm 条数，合计 87 条（Q3：2；Q4：0；Q5：0；Q6：1；Q7：2；
% Q8：4；Q9：3；Q10：3；Q11：7；Q12：3；Q13：0；Q14：0；Q15：2；Q16：14；Q17：7；
% Q18：5；Q19：12；Q20：10；Q21：12）。第 4、5 题只印答案，第 13、14 题无解析，
% 均不计入 \step 账。原卷解析的换行按语意合并，未逐行强拆。
%
\documentclass[a4paper,11pt]{article}
\usepackage{common}

\begin{document}

\examtitlest[7]{2026--2027 年金汇中学高一上 10 月月考}{集合与常用逻辑用语、不等式}

\bigsec{一、填空题}

\begin{enumerate}
\setcounter{enumi}{2}

\item 命题：如果 $ab>c$，那么 $a>\dfrac cb$，是\blankline[4.4em]（选填“真”、“假”）命题.

\ifanswer
\solbegin
\step{取 $a=1,b=-2,c=-3$，则 $ab=-2>-3=c$，但 $a=1<\dfrac cb=\dfrac32$，}
\step{所以原命题是假命题.}
\solend
\fi

\item 若“$x=0$”是“$x<m$”的充分条件，则实数 $m$ 的取值范围是\blankline[4.4em].

\ifanswer
\ans{$m>0$}
\fi

\item “$x\neq0$ 且 $x\neq1$”的否定形式为\blankline[4.4em].

\ifanswer
\ans{$x=0$ 或 $x=1$}
\fi

\item 若集合 $M=\{(x,y)\mid x+y=3\}$，$N=\left\{(x,y)\left|y=\dfrac12x\right.\right\}$，则 $M\cap N=\blankline[4.4em]$.

\ifanswer
\solbegin
\step{$x+\dfrac12x=3\Rightarrow x=2,y=1$，所以 $M\cap N=\{(2,1)\}$.}
\solend
\fi

\item 集合 $\{x\in\N\mid0<\lvert x-1\rvert<3\}$ 的真子集个数为\blankline[4.4em].

\ifanswer
\solbegin
\step{因为 $\{x\in\N\mid0<\lvert x-1\rvert<3\}=\{x\in\N\mid-3<x-1<3,x\neq1\}$}
\step{$=\{x\in\N\mid-2<x<4,x\neq1\}=\{0,2,3\}$，所以真子集的个数为 $2^3-1=7$.}
\solend
\fi

\item 设 $A=\{x\mid x^2-3x-4=0\}$，$B=\{x\mid ax-1=0\}$，若 $A\cap B=B$，则实数 $a$ 的值可以为\blankline[4.4em].

\ifanswer
\solbegin
\step{$A=\{-1,4\}$，因为 $A\cap B=B$，所以 $B\subseteq A$，}
\step{①$a=0$ 时，$B=\varnothing$，满足 $B\subseteq A$；}
\step{②$a\neq0$ 时，$B=\left\{\dfrac1a\right\}$，则 $\dfrac1a=-1$ 或 $\dfrac1a=4$，所以 $a=-1$ 或 $\dfrac14$；}
\step{所以 $a$ 的值可以为 $0$、$-1$、$\dfrac14$.}
\solend
\fi

\item 关于 $x$ 的不等式 $ax^2-(a+1)x+1\leqslant0(a<0)$ 的解集为\blankline[4.4em].

\ifanswer
\solbegin
\step{原不等式即 $(ax-1)(x-1)\leqslant0$，因为 $a<0$，所以 $\left(x-\dfrac1a\right)(x-1)\geqslant0$.}
\step{又 $\dfrac1a<0<1$，所以 $x\leqslant\dfrac1a$ 或 $x\geqslant1$，}
\step{故解集为 $\left(-\infty,\dfrac1a\right]\cup[1,+\infty)$.}
\solend
\fi

\item 若不等式 $(a-1)x^2+(2a-2)x-3<0$ 的解集是 $\R$，则 $a$ 的取值范围是\blankline[4.4em].

\ifanswer
\solbegin
\step{当 $a=1$ 时，不等式为 $-3<0$，对任意实数 $x$ 均成立.}
\step{当 $a\neq1$ 时，$\begin{cases}a-1<0\\ \Delta=(2a-2)^2+12(a-1)=4(a-1)(a+2)<0\end{cases}$，解得 $-2<a<1$.}
\step{合并得 $a\in(-2,1]$.}
\solend
\fi

\item 若关于 $x$ 的不等式 $x^2-(m+2)x+2m<0$ 的解集中恰有 $3$ 个整数，则实数 $m$ 的取值范围为\blankline[4.4em].

\ifanswer
\solbegin
\step{将不等式 $x^2-(m+2)x+2m<0$ 因式分解为 $(x-2)(x-m)<0$，}
\step{当 $m<2$ 时，解集为 $(m,2)$.}
\step{要使该区间恰有 $3$ 个整数，这 $3$ 个整数为 $-1,0,1$，故 $m$ 需满足 $-2\leqslant m<-1$.}
\step{当 $m>2$ 时，解集为 $(2,m)$. 要使该区间恰有 $3$ 个整数，}
\step{这 $3$ 个整数为 $3$，$4$，$5$，故 $m$ 需满足 $5<m\leqslant6$.}
\step{当 $m=2$ 时，不等式无解，不满足条件.}
\step{综上，实数 $m$ 的取值范围是 $[-2,-1)\cup(5,6]$.}
\solend
\fi

\item 已知集合 $A=\{x\mid x^2-3x+2\leqslant0\},B=\{y\mid y=x^2-2x+a,x\in\R\}$，若对任何一个 $x\in A$，都有 $x\in B$，则实数 $a$ 的取值范围是\blankline[4.4em].

\ifanswer
\solbegin
\step{$(x-1)(x-2)\leqslant0$，所以 $A=[1,2]$.}
\step{由 $y=(x-1)^2+a-1,x\in\R$，得 $B=[a-1,+\infty)$.}
\step{因此 $A\subseteq B\Leftrightarrow a-1\leqslant1\Leftrightarrow a\leqslant2$，所求范围为 $(-\infty,2]$.}
\solend
\fi

\end{enumerate}

\bigsec{二、选择题}

\begin{enumerate}
\setcounter{enumi}{12}

\item 已知 $a=\sqrt2$，集合 $A=\{x\mid x\leqslant2\}$，则下列表示正确的是\mbox{（\quad）}% 原卷括号内已印“D”，试卷版按惯例留空

\keepwithprev
A. $\{a\}\in A$\qquad B. $a\notin A$\qquad C. $a\subset A$\qquad D. $a\in A$

\ifanswer
\ans{$D$}
\fi

\item 下列四个命题中，真命题是\mbox{（\quad）}% 原卷括号内已印“B”，试卷版按惯例留空

\keepwithprev
A. 若 $a>b$，则 $ac^2>bc^2$\qquad B. 若 $a>b,c>d$，则 $a-d>b-c$

C. 若 $\lvert a\rvert>b$，则 $a^2>b^2$\qquad D. 若 $a>b$，则 $\dfrac1a<\dfrac1b$

\ifanswer
\ans{$B$}
\fi

\item “$x>1$”是“$\dfrac1x<1$”的\mbox{（\quad）}

\keepwithprev
A. 充分不必要条件\qquad B. 必要不充分条件

C. 充要条件\qquad D. 既不充分也不必要条件

\ifanswer
\solbegin
\step{$x>1\Rightarrow0<\dfrac1x<1$，充分性成立.}
\step{当 $x=-1$ 时，$\dfrac1x=-1<1$，但 $x>1$ 不成立，必要性不成立. 故选 $A$.}
\solend
\fi

\item 对于集合 $M=\{a\mid a=x^2-y^2,x\in\Z,y\in\Z\}$，给出如下三个结论：①如果 $P=\{b\mid b=2n+1,n\in\Z\}$，那么 $P\subseteq M$；②如果 $c=4n+2,n\in\Z$，那么 $c\notin M$；③如果 $a_1\in M,a_2\in M$，那么 $a_1a_2\in M$. 其中正确结论的个数是\mbox{（\quad）}

\keepwithprev
A. $1$\qquad B. $2$\qquad C. $3$\qquad D. $0$

\ifanswer
\solbegin
\step{对于①：因为 $2n+1=(n+1)^2-n^2,n\in\Z$，所以 $b=2n+1\in M$，故 $P\subseteq M$，}
\step{故①正确；}
\step{对于②：因为 $c=4n+2=2(2n+1),n\in\Z$，所以 $c$ 为偶数，且不能被 $4$ 整除，}
\step{若 $c\in M$，则存在 $x\in\Z,y\in\Z$，得 $c=x^2-y^2=(x+y)(x-y),x\in\Z,y\in\Z$，}
\step{因为 $x+y$ 和 $x-y$ 同奇或同偶，}
\step{若 $x+y$ 和 $x-y$ 同奇，则 $c=(x+y)(x-y)$ 为奇数，矛盾，不符合，}
\step{若 $x+y$ 和 $x-y$ 同偶，则 $c=(x+y)(x-y)$ 能被 $4$ 整除，矛盾，不符合，}
\step{所以 $c\notin M$，故②正确；}
\step{对于③：因为 $a_1\in M,a_2\in M$，}
\step{所以存在 $x_1,x_2\in\Z,y_1,y_2\in\Z$ 使得 $a_1=x_1^2-y_1^2,a_2=x_2^2-y_2^2$，}
\step{所以 $a_1a_2=(x_1^2-y_1^2)(x_2^2-y_2^2)=(x_1x_2)^2+(y_1y_2)^2-(x_1y_2)^2-(x_2y_1)^2$}
\step{$=(x_1x_2+y_1y_2)^2-(x_1y_2+x_2y_1)^2$，}
\step{因为 $x_1x_2+y_1y_2,x_1y_2+x_2y_1\in\Z$，所以 $a_1a_2\in M$，故③正确.}
\step{故选 $C$.}
\solend
\fi

\end{enumerate}

\bigsec{三、解答题}

\begin{enumerate}
\setcounter{enumi}{16}

\item （1）求不等式组 $\begin{cases}-2x^2+x+15\geqslant0\\ (x+1)(3x-1)>0\end{cases}$ 的解集；

（2）讨论解关于 $x$ 的不等式 $mx+4<m^2+2x$.

\ansspace*[12]

\ifanswer
\solbegin
\stepm{（1）}{第二个不等式可化为 $(2x+5)(x-3)\leqslant0$，解得 $-\dfrac52\leqslant x\leqslant3$.}% 原卷把“第一个”误作“第二个”，照录
\step{第二个不等式 $(x+1)(3x-1)>0$ 的解集为 $(-\infty,-1)\cup\left(\dfrac13,+\infty\right)$.}
\step{原不等式组的解集为 $\left[-\dfrac52,-1\right)\cup\left(\dfrac13,3\right]$.}
\stepm{（2）}{移项并因式分解，得 $(m-2)x<m^2-4\Leftrightarrow(m-2)x<(m-2)(m+2)$.}
\step{当 $m>2$ 时，$m-2>0$，解集为 $(-\infty,m+2)$.}
\step{当 $m<2$ 时，$m-2<0$，解集为 $(m+2,+\infty)$.}
\step{当 $m=2$ 时，不等式化为 $0<0$，不成立，解集为 $\varnothing$.}
\solend
\fi

\item （1）设 $x$ 是实数，比较 $x^2+9$ 与 $6x$ 的值的大小；

（2）已知 $a>0$ 且 $a\neq1$，比较 $\dfrac{a^2+1}{a^2-1}$ 与 $\dfrac{a+1}{a-1}$ 的大小.

\ansspace*[11]

\ifanswer
\solbegin
\stepm{（1）}{作差得 $(x^2+9)-6x=(x-3)^2\geqslant0$，因此 $x^2+9\geqslant6x$；}
\stepm{（2）}{$\dfrac{a^2+1}{a^2-1}-\dfrac{a+1}{a-1}=\dfrac{-2a}{a^2-1}$，}
\step{当 $a>1$ 时，$-2a<0$，$a^2-1>0$，则 $\dfrac{-2a}{a^2-1}<0$，即 $\dfrac{a^2+1}{a^2-1}<\dfrac{a+1}{a-1}$；}
\step{当 $0<a<1$ 时，$-2a<0$，$a^2-1<0$，则 $\dfrac{-2a}{a^2-1}>0$，即 $\dfrac{a^2+1}{a^2-1}>\dfrac{a+1}{a-1}$；}
\step{综上，$a>1$ 时，$\dfrac{a^2+1}{a^2-1}<\dfrac{a+1}{a-1}$；$0<a<1$ 时，$\dfrac{a^2+1}{a^2-1}>\dfrac{a+1}{a-1}$.}
\solend
\fi

\item 已知 $a\in\R$，集合 $A=\left\{x\left|\dfrac{2-x}{x+1}\leqslant0\right.\right\},B=\{x\mid x^2+(3-a)x-3a<0\}$.

（1）当 $a=\dfrac13$ 时，求 $A$ 和 $B$；

（2）“$x\in B$”是“$x\in A$”的充分非必要条件，求实数 $a$ 的取值范围.

\ansspace*[13]

\ifanswer
\solbegin
\stepm{（1）}{由 $\dfrac{2-x}{x+1}\leqslant0$，得 $\begin{cases}(2-x)(x+1)\leqslant0\\ x\neq-1\end{cases}$，解得 $x<-1$ 或 $x\geqslant2$，}
\step{则 $A=\{x\mid x<-1$ 或 $x\geqslant2\}$；$\cdots\cdots\cdots\cdots\cdots\cdots\cdots\cdots\cdots\cdots\cdots\cdots\cdots\cdots\cdots\cdots$（3 分）}
\step{$B=\{x\mid x^2+(3-a)x-3a<0\}=\{x\mid(x+3)(x-a)<0\}$，}
\step{当 $a=\dfrac13$ 时，$B=\left\{x\left|(x+2)\left(x-\dfrac13\right)<0\right.\right\}$，}
\step{解得 $B=\left\{x\left|-3<x<\dfrac13\right.\right\}\cdots\cdots\cdots\cdots\cdots\cdots$（6 分）}
\stepm{（2）}{由 $A\cup B=A$，得 $B\subseteq A$，}% 原卷从 A∪B=A 起论，与题干“充分非必要”差一层，照录
\step{①当 $B=\varnothing$ 时，得 $a=-3$，符合题意；$\cdots\cdots\cdots\cdots\cdots$（2 分）}
\step{②当 $B\neq\varnothing$ 时，若 $a>-3$，则 $B=\{x\mid-3<x<a\}$，}
\step{由 $B\subseteq A$，得 $a\leqslant-1$，此时 $-3<a\leqslant-1$；$\cdots\cdots\cdots\cdots\cdots$（4 分）}
\step{若 $a<-3$，则 $B=\{x\mid a<x<-3\}$，此时 $B\subseteq A$ 恒成立，}
\step{故 $a<-3$ 符合题意；$\cdots\cdots\cdots$（6 分）}
\step{综上所述，实数 $a$ 的取值范围为 $a\leqslant-1$. $\cdots\cdots\cdots\cdots\cdots$（8 分）}
\solend
\fi

\item 已知集合 $A=\{x\mid x<-2$ 或 $x\geqslant3\}$，$U=R$.

（1）求 $\overline A$；

（2）若 $B=\{x\mid mx-1>0\}$，当 $B\subseteq A$ 时，求实数 $m$ 的取值范围；

（3）若 $B=\{x\mid2m-6\leqslant x\leqslant m+4\}$，当 $B\cap\overline A=\varnothing$ 时，求实数 $m$ 的取值范围.

\ansspace*[14]

\ifanswer
\solbegin
\stepm{（1）}{因为集合 $A=\{x\mid x<-2$ 或 $x\geqslant3\}$，$U=R$，所以 $\overline A=\{x\mid-2\leqslant x<3\}$；}
\stepm{（2）}{当 $m=0$ 时，$B=\varnothing$，符合题意，}
\step{当 $m>0$ 时，$B=\left\{x\left|x>\dfrac1m\right.\right\}$，则 $\dfrac1m\geqslant3$，解得 $0<m\leqslant\dfrac13$，}
\step{当 $m<0$ 时，$B=\left\{x\left|x<\dfrac1m\right.\right\}$，则 $\dfrac1m\leqslant-2$，解得 $-\dfrac12\leqslant m<0$，}
\step{综上所述，实数 $m$ 的取值范围 $\left[-\dfrac12,\dfrac13\right]$；}
\stepm{（3）}{由题意得 $B=\{x\mid2m-6\leqslant x\leqslant m+4\}$，$\overline A=\{x\mid-2\leqslant x<3\}$，且 $B\cap\overline A=\varnothing$，}
\step{当 $B=\varnothing$ 时，$2m-6>m+4$，解得 $m>10$，}
\step{当 $B\neq\varnothing$ 时，则 $\begin{cases}2m-6\leqslant m+4\\ m+4<-2\end{cases}$ 或 $\begin{cases}2m-6\leqslant m+4\\ 2m-6\geqslant3\end{cases}$，}
\step{解得 $m<-6$ 或 $\dfrac92\leqslant m\leqslant10$，}
\step{综上所述，实数 $m$ 的取值范围为 $(-\infty,-6)\cup\left[\dfrac92,+\infty\right)$.}
\solend
\fi

\item 已知命题 $p$：关于 $x$ 的方程 $x^2+mx+1=0$ 有两个不同的正根；命题 $q$：关于 $x$ 的方程 $4x^2+4(m-2)x=-1$ 有实根.

（1）若命题 $q$ 为假命题，求整数 $m$ 的解集；

（2）若命题 $p$ 为真命题，且关于 $x$ 的方程 $x^2+mx+1=0$ 的两个实数根 $x_1$、$x_2$ 满足 $3x_1+2x_2=5$，求实数 $m$ 的取值集合；

（3）若命题 $p$ 和命题 $q$ 有且仅有一个成立，求实数 $m$ 的取值范围.

\ansspace*[16]

\ifanswer
\solbegin
\stepm{（1）}{命题 $q$：关于 $x$ 的方程 $4x^2+4(m-2)x=-1$ 有实根为假命题，}
\step{即 $4x^2+4(m-2)x+1=0$ 无实根，}
\step{则 $\Delta=16(m-2)^2-16<0$，则 $1<m<3$，则整数 $m$ 的解集为 $\{2\}$；}
\stepm{（2）}{若命题 $p$ 为真命题，且关于 $x$ 的方程 $x^2+mx+1=0$ 的两个实数根 $x_1$、$x_2$}
\step{满足 $3x_1+2x_2=5$，则方程 $x^2+mx+1=0$ 有两个不同的正根，}
\step{则 $\begin{cases}\Delta=m^2-4>0\\ x_1+x_2=-m>0\end{cases}$，则 $m<-2$，又 $\begin{cases}x_1+x_2=-m\\ x_1x_2=1\\ 3x_1+2x_2=5\end{cases}$，则 $\begin{cases}x_1=5+2m\\ x_2=-3m-5\end{cases}$，}
\step{则 $(5+2m)(-3m-5)=1$，得 $m=-2$ 或 $m=-\dfrac{13}6$，}
\step{又因为 $m<-2$，所以 $m=-\dfrac{13}6$，所以实数 $m$ 的取值集合为 $\left\{-\dfrac{13}6\right\}$；}
\stepm{（3）}{由（2）得，当 $p$ 为真时，则有 $m<-2$，}
\step{当 $q$ 为真时，$\Delta=16(m-2)^2-16\geqslant0$，解得 $m\leqslant1$ 或 $m\geqslant3$，}
\step{又因为命题 $p$ 和命题 $q$ 只有一个为真，所以 $\begin{cases}m<-2\\ 1<m<3\end{cases}$ 或 $\begin{cases}m\geqslant-2\\ m\leqslant1\ \text{或}\ m\geqslant3\end{cases}$，}
\step{解得 $-2\leqslant m\leqslant1$ 或 $m\geqslant3$，所以实数 $m$ 的取值范围为 $[-2,1]\cup[3,+\infty)$.}
\solend
\fi

\end{enumerate}

\end{document}
